\appendix

\section{Proof Details}\label{sec:proof-details}

\subsection{Exact class parameters}\label{app:step-001}
We use the following explicit convention. A complete binary instance tree of depth \(D\) over a domain \(\mathcal X\) has an instance \(x_\sigma\in\mathcal X\) at every node \(\sigma\in\{0,1\}^{<D}\). A binary class \(H\subseteq\{0,1\}^{\mathcal X}\) shatters the tree if, for every \(b=(b_1,\ldots,b_D)\in\{0,1\}^D\), there is \(h_b\in H\) such that

\[
h_b(x_{b_{<r}})=b_r\qquad(1\le r\le D),
\]

where \(b_{<r}=(b_1,\ldots,b_{r-1})\). The Littlestone dimension is the largest shattered depth. This is the convention used throughout the proof; no tree/mistake equivalence is assumed without proof.

\begin{lemma}[Distinct tagged cuts and exact cardinality]\label{lem:step-001-cardinality}
Under the primitive tagged-threshold definitions with integers \(k\ge2\) and \(N\ge2\), the functions \(\tau_t\), \(t\in[N+1]\), are pairwise distinct, the parameterization \(\boldsymbol t\mapsto c_{\boldsymbol t}\) is injective, and

\[
\lvert C_{k,N}\rvert=(N+1)^k.
\]
\end{lemma}

\begin{proof}
If \(t<s\) in \([N+1]\), then \(t\le N\), so \(x=t\) is a valid point of \([N]\). At that point,

\[
\tau_t(t)=1,\qquad \tau_s(t)=0.
\]

Thus all \(N+1\) cuts are distinct. This comparison includes the endpoints: \(\tau_1\equiv1\) and \(\tau_{N+1}\equiv0\) on \([N]\), and they disagree at every point.

If \(\boldsymbol t\ne\boldsymbol s\), choose a tag \(i\) for which \(t_i\ne s_i\). The preceding one-tag argument supplies an \(x\in[N]\) with \(\tau_{t_i}(x)\ne\tau_{s_i}(x)\), hence

\[
c_{\boldsymbol t}(i,x)\ne c_{\boldsymbol s}(i,x).
\]

The parameter map is therefore injective. It is surjective onto \(C_{k,N}\) by the class definition, and its domain has size \((N+1)^k\), proving the formula. \end{proof}

\begin{lemma}[VC dimension of the tagged threshold product]\label{lem:step-001-vc}
Under the primitive tagged-threshold definitions with integers \(k\ge2\) and \(N\ge2\),

\[
\operatorname{VC}(C_{k,N})=k.
\]
\end{lemma}

\begin{proof}
For the lower bound, consider

\[
S=\{(i,1):i\in[k]\}.
\]

Given arbitrary desired labels \(b_1,\ldots,b_k\in\{0,1\}\), set

\[
t_i=\begin{cases}
1,&b_i=1,\\
2,&b_i=0.
\end{cases}
\]

Both choices lie in \([N+1]\) because \(N\ge2\), and \(\tau_{t_i}(1)=b_i\). Hence \(S\) is shattered and \(\operatorname{VC}(C_{k,N})\ge k\).

For the upper bound, any set of \(k+1\) distinct points of \([k]\times[N]\) contains two points \((i,x)\) and \((i,y)\) on the same tag. Relabel them so that \(x<y\). For every threshold \(t_i\),

\[
\tau_{t_i}(x)=1\quad\Longrightarrow\quad x\ge t_i
\quad\Longrightarrow\quad y\ge t_i
\quad\Longrightarrow\quad \tau_{t_i}(y)=1.
\]

Thus the labeling \((1,0)\) on these two points cannot be realized by any member of \(C_{k,N}\). No \((k+1)\)-point set is shattered, so \(\operatorname{VC}(C_{k,N})\le k\). \end{proof}

\begin{lemma}[Exact one-tag Littlestone dimension and halving budget]\label{lem:step-001-one-tag-ld}
Under the primitive one-tag definition

\[
T_N=\{\tau_t:t\in[N+1]\}\subseteq\{0,1\}^{[N]},
\]

with \(N\ge2\), and Lemma~\ref{lem:step-001-cardinality}, define

\[
q=N+1,\qquad d_N=\lfloor\log_2 q\rfloor.
\]

Then \(T_N\) shatters a complete binary instance tree of depth \(d_N\), shatters no tree of depth greater than \(d_N\), and admits a deterministic online prediction rule that makes at most \(d_N\) mistakes on every finite labeled sequence realizable by \(T_N\). Consequently,

\[
\operatorname{LD}(T_N)=d_N.
\]
\end{lemma}

\begin{proof}
First, suppose \(T_N\) shatters a tree of depth \(D\). Choose one witnessing concept for each of the \(2^D\) root-to-leaf label sequences. Two different label sequences first diverge at some common node and prescribe opposite labels to the same queried point there. Their witnessing concepts must therefore be different. Lemma~\ref{lem:step-001-cardinality} gives exactly \(q\) distinct one-tag concepts, so

\[
2^D\le q,
\]

and hence \(D\le\lfloor\log_2q\rfloor=d_N\).

For the matching lower bound, put \(m=2^{d_N}\), so \(m\le q=N+1\) and \(m-1\le N\). We construct a depth-\(d_N\) tree using only the thresholds \(t\in[m]\). At a node whose surviving threshold set is a consecutive interval

\[
I=\{a,a+1,\ldots,a+2^s-1\}
\]

of size \(2^s\), where \(s\ge1\), query

\[
x_I=a+2^{s-1}-1.
\]

This is a valid point of \([N]\): throughout the construction \(1\le x_I\le m-1\le N\). On the edge labeled \(1\), retain the lower half \(\{a,\ldots,x_I\}\); on the edge labeled \(0\), retain the upper half \(\{x_I+1,\ldots,a+2^s-1\}\). In the first case every retained threshold obeys \(t\le x_I\) and therefore \(\tau_t(x_I)=1\); in the second case every retained threshold obeys \(t>x_I\) and therefore \(\tau_t(x_I)=0\). Starting from \(I=[m]\), every root-to-leaf path leaves a singleton after exactly \(d_N\) splits, and that singleton threshold realizes every edge label on the path. The constructed tree is therefore shattered.

It remains to record the online interface needed for the product upper bound. Maintain the version space \(V\subseteq T_N\) of concepts consistent with all previously revealed examples, starting from \(V=T_N\). At a query \(x\), predict a label having the larger number of concepts in

\[
V_y(x)=\{\tau\in V:\tau(x)=y\},\qquad y\in\{0,1\},
\]

with an arbitrary fixed tie rule. After the true label \(y\) is revealed, replace \(V\) by \(V_y(x)\). On a realizable sequence, the realizing threshold remains in \(V\), so \(V\) never becomes empty. Whenever the prediction is wrong, the surviving side has cardinality at most half the previous cardinality. If \(r\) mistakes occur, then

\[
1\le \lvert V\rvert\le q\,2^{-r},
\]

so \(2^r\le q\) and therefore \(r\le d_N\). This also covers the endpoint targets \(t=1\) and \(t=N+1\): each belongs to the initial version space and remains present whenever it is the realizing target. \end{proof}

\begin{lemma}[Concatenated tagged trees give the product lower bound]\label{lem:step-001-product-lower}
Under the primitive definition of \(C_{k,N}\), with \(k\ge2,N\ge2\), and the depth-\(d_N\) one-tag tree from Lemma~\ref{lem:step-001-one-tag-ld}, the product class shatters a complete tree of depth \(kd_N\). Hence

\[
\operatorname{LD}(C_{k,N})\ge kd_N.
\]
\end{lemma}

\begin{proof}
Attach a copy of the one-tag tree on tag \(1\) first. Below every leaf attach a copy on tag \(2\), and continue through tag \(k\). Equivalently, split any path label vector \(b\in\{0,1\}^{kd_N}\) into consecutive blocks

\[
b=(b^{(1)},\ldots,b^{(k)}),\qquad b^{(i)}\in\{0,1\}^{d_N}.
\]

During block \(i\), if the corresponding one-tag copy queries \(x\in[N]\), the product tree queries \((i,x)\). By Lemma~\ref{lem:step-001-one-tag-ld}, for each block \(b^{(i)}\) there is a threshold \(t_i\in[N+1]\) whose one-tag labels realize that entire block. Because the coordinates \(t_1,\ldots,t_k\) are independent in the defining parameter set \([N+1]^k\), the single product concept \(c_{(t_1,\ldots,t_k)}\) realizes all \(k\) blocks of the chosen path. Every length-\(kd_N\) path is therefore realized, proving the stated lower bound. \end{proof}

\begin{lemma}[Summed per-tag realizable mistake budget]\label{lem:step-001-product-budget}
Under the primitive definition of \(C_{k,N}\), with \(k\ge2,N\ge2\), and the one-tag online rule of Lemma~\ref{lem:step-001-one-tag-ld}, there is a deterministic online rule that makes at most \(kd_N\) mistakes on every finite sequence realizable by \(C_{k,N}\).
\end{lemma}

\begin{proof}
Maintain \(k\) independent one-tag version spaces \(V_1,\ldots,V_k\), each initialized to \(T_N\). When the current instance is \((i,x)\), use the one-tag halving rule based only on \(V_i\), reveal the resulting prediction, and update only \(V_i\) after observing the label.

If the full sequence is realized by \(c_{\boldsymbol t}\), then the subsequence on tag \(i\) is realized by \(\tau_{t_i}\). Lemma~\ref{lem:step-001-one-tag-ld} therefore bounds the number \(M_i\) of mistakes charged to tag \(i\) by \(d_N\). Each global mistake occurs on exactly one tag and is charged exactly once, so

\[
M_{\mathrm{total}}=\sum_{i=1}^k M_i
\le\sum_{i=1}^k d_N
=kd_N.
\]

There is no cross-tag condition to maintain: observations on tag \(i\) constrain only \(t_i\). The same argument includes \(k=2\) and includes any coordinate whose realizing threshold is either endpoint. \end{proof}

\begin{lemma}[A realizable mistake budget upper-bounds Littlestone dimension]\label{lem:step-001-mistake-to-ld}
Let \(H\subseteq\{0,1\}^{\mathcal X}\) be any binary class and let \(B\in\mathbb Z_{\ge0}\). If a deterministic online prediction rule makes at most \(B\) mistakes on every finite labeled sequence realizable by \(H\), then

\[
\operatorname{LD}(H)\le B.
\]
\end{lemma}

\begin{proof}
Suppose instead that \(H\) shatters a complete tree of depth \(B+1\). Starting at its root, present the instance at the current node to the rule. If the rule predicts \(p\in\{0,1\}\), reveal the label \(1-p\) and descend along that edge. Repeat for all \(B+1\) levels. By construction the rule makes a mistake at every level. By shattering, the resulting root-to-leaf labeled sequence is realized by some single member of \(H\). This is a realizable sequence with \(B+1\) mistakes, contradicting the assumed budget. \end{proof}

\begin{proposition}[Exact Littlestone dimension of the tagged product]\label{prop:step-001-product-ld}
Under the primitive tagged-threshold definitions with integers \(k\ge2,N\ge2\), and Lemmas~\ref{lem:step-001-product-lower}, \ref{lem:step-001-product-budget}, and \ref{lem:step-001-mistake-to-ld},

\[
\operatorname{LD}(C_{k,N})
=k\lfloor\log_2(N+1)\rfloor.
\]
\end{proposition}

\begin{proof}
Lemma~\ref{lem:step-001-product-lower} gives

\[
\operatorname{LD}(C_{k,N})\ge kd_N.
\]

Lemma~\ref{lem:step-001-product-budget} constructs a deterministic rule with realizable mistake budget \(kd_N\). Applying Lemma~\ref{lem:step-001-mistake-to-ld} with \(H=C_{k,N}\) and \(B=kd_N\) gives

\[
\operatorname{LD}(C_{k,N})\le kd_N.
\]

The bounds agree, and \(d_N=\lfloor\log_2(N+1)\rfloor\). \end{proof}

Together, Lemmas~\ref{lem:step-001-cardinality} and
\ref{lem:step-001-vc} and
Proposition~\ref{prop:step-001-product-ld} establish all three exact
structural identities used below, including both endpoint cuts.

\subsection{The unrestricted one-block threshold lower bound}\label{app:step-002}
For a randomized one-block learner
$B:([N]\times\{0,1\})^M\to\{0,1\}^{[N]}$, write \emph{(W-PAC)} for
\[
\Pr_{\substack{S\sim(Q^{\tau_t})^M\\g\sim B(S)}}
\left\{R_Q(g,\tau_t)\le\frac1{16}\right\}
\ge\frac{15}{16}
\quad\text{for every }(t,Q)\in\mathcal I_N,
\]
and write \emph{(W-LB)} for $M\ge b_*\log_2^*N$.  The external input is
the unrestricted improper threshold theorem of
\citet{AlonLivniMalliarisMoran2019}; all convention and parameter transport
is proved here.
\begin{lemma}[Fixed-constant form of the ALMM v3 threshold theorem]\label{lem:step-002-almm-constants}
The threshold theorem of \citet[Theorem~1]{AlonLivniMalliarisMoran2019},
together with its PAC and privacy definitions and the fixed constants in its
proof, implies that there exist universal constants

\[
b_0>0,\qquad d_0>0,\qquad N_0\in\mathbb Z_{\ge2}
\]

such that, for every finite ordered real domain \(X\) with \(|X|=N\ge N_0\), every integer \(M\ge8\), and every possibly randomized, possibly improper threshold learner on \(X\) that is \((1/16,1/16)\)-accurate in the source PAC sense and

\[
\left(0.1,\frac{d_0}{M^2\log M}\right)
\text{-differentially private},
\]

one has

\[
M\ge b_0\log_2^*N.
\]
\end{lemma}

\begin{proof}
The source theorem states

\[
M=\Omega(\log^*N)
\quad\text{under}\quad
\varepsilon=0.1,\qquad
\delta=O\!\left(\frac1{M^2\log M}\right).
\tag{S2.1}\label{eq:step-002-1}
\]

The constants in both asymptotic symbols are absolute: the theorem fixes the accuracy and privacy constants before quantifying over the domain and learner. This is also visible in the active proof. Its private homogeneous-set lemma uses the explicit empirical-size cap

\[
\delta\le \frac1{10^3m^2\log m},
\tag{S2.2}\label{eq:step-002-2}
\]

and the source's constant-factor empirical reduction increases the original
PAC sample size by at most the absolute factor nine while preserving accuracy
and privacy.

For completeness, the factor-nine change does not leave an unverified denominator. The function

\[
x\longmapsto \frac{\log x}{\log(9x)}
\]

is increasing for \(x>1\), because its derivative is

\[
\frac{\log 9}{x(\log(9x))^2}>0.
\]

Hence, for every \(M\ge8\),

\[
\frac{\log M}{\log(9M)}
\ge \frac{\log8}{\log72}.
\tag{S2.3}\label{eq:step-002-3}
\]

Choose \(d_0\) no larger than the fixed privacy constant represented by the
\(O(\cdot)\) in \eqref{eq:step-002-1}, and shrink it further, if needed, so that

\[
0<d_0\le \frac{\log8}{81000\log72}\,.
\tag{S2.4}\label{eq:step-002-4}
\]

Then it obeys

\[
\frac{d_0}{M^2\log M}
\le
\frac1{10^3(9M)^2\log(9M)}.
\tag{S2.5}\label{eq:step-002-5}
\]

If the empirical reduction uses a size \(m\le9M\), then
\(m^2\log m\le(9M)^2\log(9M)\), so \eqref{eq:step-002-5} also implies \eqref{eq:step-002-2} at that actual
size. Thus the source proof exhibits a strictly positive universal privacy
constant after its own constant-factor reduction; \eqref{eq:step-002-4} is only a conservative
witness and is not claimed as a numerical constant printed in the theorem
statement.

Unpacking the \(\Omega\)-conclusion in \eqref{eq:step-002-1} gives an absolute lower-bound constant and a finite domain threshold. The source's unadorned iterated-log convention and the branch's base-two convention differ by at most an absolute additive number of iterations: for some universal \(K_{\rm base}\),

\[
\log^*_{\rm source}N
\ge \log_2^*N-K_{\rm base}.
\tag{S2.6}\label{eq:step-002-6}
\]

This is the standard fixed-base comparison: multiplying a logarithm by the fixed base-change constant can require only a fixed number of additional iterates before entering \((0,1]\). Increase the finite domain threshold so that

\[
\log_2^*N\ge2K_{\rm base}
\]

throughout the retained regime. Then \eqref{eq:step-002-6} is at least \(\frac12\log_2^*N\), and halving the source lower-bound constant gives a universal \(b_0>0\). Absorb every finite exceptional domain and every fixed source convention into one integer \(N_0\ge2\). This proves the displayed fixed-constant form for all \(M\ge8\). No value of \(b_0\) or \(N_0\) beyond existence is asserted. \end{proof}

\begin{lemma}[Exact order, label, PAC, output, and adjacency transport]\label{lem:step-002-transport}
Under the primitive one-block definitions with \(N\ge2\), and consistently with the endpoint/distinctness certificate in Lemma~\ref{lem:step-001-cardinality}, let

\[
B:([N]\times\{0,1\})^M\longrightarrow\{0,1\}^{[N]}
\]

be any randomized map. There is a rowwise bijective transport to a possibly randomized, unrestricted ALMM learner

\[
\widetilde B:(X\times\{-1,+1\})^M
\longrightarrow\{-1,+1\}^{X}
\]

on an \(N\)-point ordered real set \(X\) such that:

\begin{enumerate}
\item All increasing thresholds, including the all-positive and all-negative
endpoint restrictions, correspond to \(\{\tau_t:t\in[N+1]\}\).
\item Every arbitrary output of \(B\) corresponds bijectively to an
arbitrary output of \(\widetilde B\), with no properness or monotonicity.
\item For every \(t,Q\), the labeled sample laws are exact i.i.d. products
of the same fixed size \(M\), and population risks agree pointwise.
\item Two input datasets differ in one labeled row by replacement before
transport if and only if their images do so after transport.
\item \emph{(W-PAC)} for every \(t,Q\) is exactly the source realizable-PAC
premise for \(\widetilde B\).
\end{enumerate}

The same conclusions hold if a presentation of the source reverses the threshold orientation, after reversing the order or complementing all labels.
\end{lemma}

\begin{proof}
Choose any increasing real sequence

\[
x_1<x_2<\cdots<x_N
\]

and put \(X=\{x_1,\ldots,x_N\}\). Let \(\phi:[N]\to X\) be the order bijection \(\phi(j)=x_j\), and let

\[
\lambda(-1)=0,\qquad \lambda(+1)=1.
\]

For the increasing source threshold with cut index \(t\), write

\[
s_t(x_j)=
\begin{cases}
-1,&j<t,\\
+1,&j\ge t.
\end{cases}
\]

Then, for every \(j\in[N]\) and \(t\in[N+1]\),

\[
\lambda(s_t(\phi(j)))=\mathbf1\{j\ge t\}=\tau_t(j).
\tag{S2.7}\label{eq:step-002-7}
\]

The endpoint \(t=1\) is all positive/all one, and \(t=N+1\) is all negative/all zero. They are restrictions of real half-line thresholds with the cut at or below \(x_1\), respectively above \(x_N\). Lemma~\ref{lem:step-001-cardinality} independently certifies that these are legal, distinct branch members. Even under a convention that lists only interior cut representatives, that source family is a subclass of the branch family, so a learner satisfying (W-PAC) still meets every source target requirement.

Transport a source labeled row by

\[
(x_j,y)\longmapsto(j,\lambda(y)).
\tag{S2.8}\label{eq:step-002-8}
\]

For a source sample \(S\), apply \eqref{eq:step-002-8} row by row to obtain \(S^{\flat}\), run \(B(S^{\flat})\), and define

\[
\widetilde B(S)(x_j)
=\lambda^{-1}\!\left(B(S^{\flat})(j)\right).
\tag{S2.9}\label{eq:step-002-9}
\]

Equation \eqref{eq:step-002-9} is defined for every branch output, including constant, oscillating, and nonmonotone functions. It transports the internal randomness unchanged and uses no projection onto the target class.

Let \(Q\) be a probability law on \([N]\), and let \(\widetilde Q=\phi_{\#}Q\) be its pushforward to \(X\). For every deterministic branch output \(g\), with transported source output \(\widetilde g(x_j)=\lambda^{-1}(g(j))\), \eqref{eq:step-002-7} gives

\[
\begin{aligned}
L_{\widetilde Q^{s_t}}(\widetilde g)
&=\Pr_{J\sim Q}
   [\widetilde g(x_J)\ne s_t(x_J)]\\
&=\Pr_{J\sim Q}[g(J)\ne\tau_t(J)]\\
&=R_Q(g,\tau_t).
\end{aligned}
\tag{S2.10}\label{eq:step-002-10}
\]

The row transport is deterministic and coordinatewise, so

\[
S\sim(\widetilde Q^{s_t})^M
\quad\Longleftrightarrow\quad
S^{\flat}\sim(Q^{\tau_t})^M.
\tag{S2.11}\label{eq:step-002-11}
\]

There is no random or expected sample-size substitution. Every distribution on labeled source examples realizable by a threshold equals \(\widetilde Q^{s_t}\) for its feature marginal \(\widetilde Q\) and its realizing cut (up to changes on zero-mass points). Thus \eqref{eq:step-002-10}--\eqref{eq:step-002-11}, including the learner's unchanged randomness, identify (W-PAC) with the source PAC premise for every realizable threshold distribution.

Finally, \eqref{eq:step-002-8} is a bijection on individual labeled rows and is applied without mixing row indices. Consequently two length-\(M\) samples are equal or differ at exactly one row before transport exactly when the corresponding samples are equal or differ at that same row afterward. Output relabeling in \eqref{eq:step-002-9} is bijective postprocessing. Hence the replacement-DP event inequalities are identical under the event bijection. The source's two-sided indistinguishability is equivalent to the branch's one-sided display quantified over all ordered adjacent pairs, because adjacency is symmetric.

If a threshold convention uses positive labels below the cut, either replace \(\lambda\) by the complemented label bijection or replace \(\phi(j)\) by the order-reversing bijection \(x_{N+1-j}\). These are still rowwise bijections, so \eqref{eq:step-002-10}--\eqref{eq:step-002-11} and adjacency preservation remain exact. \end{proof}

\begin{lemma}[Privacy monotonicity at the source cap]\label{lem:step-002-privacy-monotonicity}
Under the fixed-size replacement convention, let \(M\ge8\), \(d_0>0\),
\(0\le\varepsilon\le0.1\), and

\[
0\le\delta\le\frac{d_0}{M^2\log M}.
\]

If a randomized kernel is \((\varepsilon,\delta)\)-differentially private, then it is also

\[
\left(0.1,\frac{d_0}{M^2\log M}\right)
\text{-differentially private}.
\]

This conclusion includes \(M=8\), \(\varepsilon<0.1\), \(\delta\) strictly below the cap, and \(\delta=0\).
\end{lemma}

\begin{proof}
For adjacent datasets \(S,S'\) and every output event \(E\), privacy gives

\[
\Pr[B(S)\in E]
\le e^{\varepsilon}\Pr[B(S')\in E]+\delta.
\]

Since the event probability is nonnegative, \(e^{\varepsilon}\le e^{0.1}\), and
\(\delta\le d_0/(M^2\log M)\),

\[
\Pr[B(S)\in E]
\le e^{0.1}\Pr[B(S')\in E]
   +\frac{d_0}{M^2\log M}.
\tag{S2.12}\label{eq:step-002-12}
\]

Swapping \(S,S'\) gives the reverse source inequality. There is one application of privacy at the same sample size; no group privacy, composition, subsampling, or adjacency conversion occurs. At \(M=8\), \(M^2\log M=64\log8>0\), so the cap in \eqref{eq:step-002-12} is finite and positive. Decreasing either privacy parameter only strengthens \eqref{eq:step-002-12}. \end{proof}

\begin{proposition}[Current-notation unrestricted ALMM threshold wrapper]\label{prop:step-002-wrapper}
Under Lemmas~\ref{lem:step-002-almm-constants},
\ref{lem:step-002-transport}, and
\ref{lem:step-002-privacy-monotonicity}, there are universal
\(b_*,d_*>0\) and \(N_*\ge2\) such that, for every \(N\ge N_*\), integer \(M\ge8\),
\(0\le\varepsilon\le0.1\), and
\(0\le\delta\le d_*/(M^2\log M)\), every randomized unrestricted map

\[
B:([N]\times\{0,1\})^M\to\{0,1\}^{[N]}
\]

that is \((\varepsilon,\delta)\)-DP under one-row replacement and satisfies (W-PAC) for every \(t\in[N+1]\) and every law \(Q\) on \([N]\) must satisfy (W-LB):

\[
M\ge b_*\log_2^*N.
\]
\end{proposition}

\begin{proof}
Take

\[
b_*=b_0,\qquad d_*=d_0,\qquad N_*=N_0
\]

from Lemma~\ref{lem:step-002-almm-constants}. Lemma~\ref{lem:step-002-privacy-monotonicity} promotes the stated privacy guarantee to

\[
\left(0.1,\frac{d_0}{M^2\log M}\right)\text{-DP}.
\]

Apply the exact row/order/label transport of Lemma~\ref{lem:step-002-transport}. It produces a possibly randomized source learner whose output ranges over all \(\{\pm1\}\)-valued functions, whose input consists of exactly \(M\) iid labeled examples, whose loss is the identical population 0-1 risk, and whose adjacency/privacy guarantee is the identical one-entry replacement guarantee. By (W-PAC) and \eqref{eq:step-002-10}--\eqref{eq:step-002-11}, this transported learner is \((1/16,1/16)\)-accurate for every source threshold-realizable distribution. All hypotheses of Lemma~\ref{lem:step-002-almm-constants} are therefore discharged, so

\[
M\ge b_0\log_2^*N=b_*\log_2^*N.
\]

The proof does not select a hard instance and does not assert expected-risk hardness. Those are conclusions of later contrapositive and minimax steps, not outputs of this proposition. \end{proof}
Lemmas~\ref{lem:step-002-almm-constants}--
\ref{lem:step-002-privacy-monotonicity} yield
Proposition~\ref{prop:step-002-wrapper}, the precise one-block implication
used in the remainder of the proof.

\subsection{A strict expected-loss hard instance}\label{app:step-003}
\begin{lemma}[Expectation-to-PAC conversion at the ALMM constants]\label{lem:step-003-markov-bridge}
Under the primitive one-block definitions, fix integers \(N\ge2\) and
\(M\ge1\), and let

\[
B:([N]\times\{0,1\})^M\to\{0,1\}^{[N]}
\]

be any randomized unrestricted map. If, for every \(t\in[N+1]\) and every
probability law \(Q\) on \([N]\),

\[
\mathbb E_{\substack{S\sim(Q^{\tau_t})^M\\g\sim B(S)}}
  R_Q(g,\tau_t)
\le 2^{-8},
\tag{S3.1}\label{eq:step-003-1}
\]

then, for every such \(t,Q\),

\[
\Pr_{\substack{S\sim(Q^{\tau_t})^M\\g\sim B(S)}}
  \left[R_Q(g,\tau_t)>\frac1{16}\right]
\le\frac1{16}.
\tag{S3.2}\label{eq:step-003-2}
\]

Thus \(B\) satisfies (W-PAC).
\end{lemma}

\begin{proof}
Fix an arbitrary \(t\in[N+1]\) and arbitrary law \(Q\) on \([N]\), while
keeping the same fixed sample size \(M\). Under the joint randomness of
\(S\sim(Q^{\tau_t})^M\) and \(g\sim B(S)\), define the proof-local random
variable

\[
Z_{t,Q}:=R_Q(g,\tau_t).
\]

For every realized output \(g\), population 0-1 risk is a probability, so

\[
0\le Z_{t,Q}\le1.
\tag{S3.3}\label{eq:step-003-3}
\]

Apply Markov's inequality at the positive threshold \(1/16\). The strict
failure event in (W-PAC) is contained in the weak-threshold event, hence

\[
\begin{aligned}
\Pr\!\left(Z_{t,Q}>\frac1{16}\right)
&\le \Pr\!\left(Z_{t,Q}\ge\frac1{16}\right)\\
&\le 16\,\mathbb E Z_{t,Q}\\
&\le 16\cdot2^{-8}\\
&=2^{-4}
=\frac1{16}.
\end{aligned}
\tag{S3.4}\label{eq:step-003-4}
\]

The final inequality remains valid when the expectation in \eqref{eq:step-003-1} equals
\(2^{-8}\); no strict inequality has been inserted. Since \((t,Q)\) was
arbitrary, \eqref{eq:step-003-4} holds simultaneously as a universally quantified statement
for every fixed \((t,Q)\), which is exactly (W-PAC). Randomization and
improperness of \(g\) do not affect \eqref{eq:step-003-3}. \end{proof}

\begin{proposition}[Expected-loss hardness below the one-block threshold bound]\label{prop:step-003-expected-hardness}
Let \(b_*,d_*>0\) and \(N_*\ge2\) be the constants supplied by
Proposition~\ref{prop:step-002-wrapper}. Under that proposition and
Lemma~\ref{lem:step-003-markov-bridge}, for every pair of integers
\(N\ge N_*\), \(M\ge8\) satisfying

\[
M<b_*\log_2^*N,
\]

every randomized unrestricted map

\[
B:([N]\times\{0,1\})^M\to\{0,1\}^{[N]}
\]

that is
\(\left(0.1,d_*/(M^2\log M)\right)\)-DP under one-row replacement admits
some \(t\in[N+1]\) and some probability law \(Q\) on \([N]\) for which

\[
\mathbb E_{\substack{S\sim(Q^{\tau_t})^M\\g\sim B(S)}}
  R_Q(g,\tau_t)
>2^{-8}.
\tag{S3.5}\label{eq:step-003-5}
\]
\end{proposition}

\begin{proof}
Fix \(N,M\) and an arbitrary randomized unrestricted \(B\) satisfying all
the displayed hypotheses. Suppose \eqref{eq:step-003-5} were false. The exact logical
negation of the existential strict inequality is

\[
\text{for every }t\in[N+1]\text{ and every law }Q\text{ on }[N],
\qquad
\mathbb E_{\substack{S\sim(Q^{\tau_t})^M\\g\sim B(S)}}
  R_Q(g,\tau_t)
\le2^{-8}.
\tag{S3.6}\label{eq:step-003-6}
\]

Lemma~\ref{lem:step-003-markov-bridge} applies to the same map \(B\), the
same exact sample size \(M\), and every fixed \((t,Q)\). It converts \eqref{eq:step-003-6} into
(W-PAC), including the allowed equality
\(\Pr(R_Q(g,\tau_t)>1/16)=1/16\).

All remaining hypotheses of Proposition~\ref{prop:step-002-wrapper} now hold:
\(N\ge N_*\), integer \(M\ge8\), the map is arbitrary and improper, and its
privacy parameters equal the permitted endpoints

\[
\varepsilon=0.1,
\qquad
\delta=\frac{d_*}{M^2\log M}.
\]

The denominator is positive even at \(M=8\), since \(\log8>0\). The wrapper therefore yields

\[
M\ge b_*\log_2^*N,
\]

contradicting the strict hypothesis \(M<b_*\log_2^*N\). Hence \eqref{eq:step-003-6} is false,
so a fixed pair \((t,Q)\) satisfying \eqref{eq:step-003-5} exists. The pair may depend on the
fixed learner \(B\), because only worst-instance hardness is asserted here; no
learner-independent prior or minimax conclusion is asserted. \end{proof}
The Markov conversion in Lemma~\ref{lem:step-003-markov-bridge} and the
contrapositive of Proposition~\ref{prop:step-002-wrapper} give
Proposition~\ref{prop:step-003-expected-hardness} with strict gap
$\eta=2^{-8}$.

\subsection{A finite public minimax prior}\label{app:step-004}
Let $b_*,d_*>0$ and $N_*\ge2$ be the constants from
Proposition~\ref{prop:step-003-expected-hardness}, and set
$\eta=2^{-8}$.  Throughout this subsection the scalar hard regime is
explicitly
\begin{equation}
N\ge N_*,\qquad M\in\mathbb Z_{\ge8},\qquad
M<b_*\log_2^*N.
\tag{S4.H}\label{eq:step-004-hard-regime}
\end{equation}
For these fixed $N,M$, the interface to be constructed is one finite law
$\mu_{N,M}$, selected before the learner, such that every randomized
unrestricted one-block map
$B:([N]\times\{0,1\})^M\to\{0,1\}^{[N]}$ satisfying the source privacy cap
obeys the implication
\begin{equation}
\left[
B\text{ is }\left(0.1,\frac{d_*}{M^2\log M}\right)\text{-DP}
\right]
\quad\Longrightarrow\quad
\mathbb E_{(t,Q)\sim\mu_{N,M}}
\mathbb E_{\substack{S\sim(Q^{\tau_t})^M\\g\sim B(S)}}
R_Q(g,\tau_t)>\eta.
\tag{S4.HP}\label{eq:step-004-hard-prior-interface}
\end{equation}
The quantifier order in \eqref{eq:step-004-hard-prior-interface} is
$\exists\mu_{N,M}\ \forall B$.  The objects
$\mathcal K_{N,M}$, $\mathcal I_N$, and their payoff are all finite
dimensional, although $\mathcal I_N$ itself contains a continuum of
distributions.
\begin{lemma}[Exact compact polytope of source-private kernels]\label{lem:step-004-kernel-polytope}
Under the primitive finite one-block definitions, let \(d_*>0\), fix
integers \(N\ge2\), \(M\ge8\), and put
\(\bar\delta_M=d_*/(M^2\log M)\). Define

\[
\mathcal S_{N,M}:=([N]\times\{0,1\})^M,
\qquad
\mathcal G_N:=\{0,1\}^{[N]},
\qquad
D:=|\mathcal S_{N,M}|\,|\mathcal G_N|=(2N)^M2^N.
\]

Let \(\mathcal K_{N,M}\subset\mathbb R^D\) consist of the arrays
\(K=(K_{s,g})_{s\in\mathcal S_{N,M},g\in\mathcal G_N}\) satisfying

\[
K_{s,g}\ge0,
\qquad
\sum_{g\in\mathcal G_N}K_{s,g}=1
\quad(s\in\mathcal S_{N,M}),
\tag{S4.6}\label{eq:step-004-6}
\]

and, for every \emph{ordered} adjacent pair \(s\simeq s'\) and every event
\(E\subseteq\mathcal G_N\),

\[
\sum_{g\in E}K_{s,g}
\le
e^{0.1}\sum_{g\in E}K_{s',g}+\bar\delta_M.
\tag{S4.7}\label{eq:step-004-7}
\]

Then \(\mathcal K_{N,M}\) is a nonempty compact convex polytope, and its
points are in exact bijection with the output kernels of randomized
unrestricted \((0.1,\bar\delta_M)\)-DP maps
\(B:\mathcal S_{N,M}\to\mathcal G_N\).
\end{lemma}

\begin{proof}
Both \(\mathcal S_{N,M}\) and \(\mathcal G_N\) are finite. Consequently,
\eqref{eq:step-004-6} is a finite family of linear equalities and inequalities. Although the
all-event privacy family is large, it is also finite: there are finitely many
ordered adjacent pairs and exactly \(2^{|\mathcal G_N|}\) subsets of the
finite output space. Thus \eqref{eq:step-004-7} is the \emph{exact} all-event approximate-DP
definition expressed as finitely many closed affine halfspaces. In
particular, no invalid replacement of all-event privacy by singleton-only
constraints is made.

The product of row simplexes defined by \eqref{eq:step-004-6} is closed and bounded in
\(\mathbb R^D\), hence compact. Intersecting it with the finitely many closed
halfspaces \eqref{eq:step-004-7} leaves a compact set. All constraints are affine, so the
intersection is convex and is a polytope.

For nonemptiness, fix any \(g_0\in\mathcal G_N\) and set

\[
K^0_{s,g}:=\mathbf 1\{g=g_0\}
\quad\text{for every }s,g.
\]

The output distribution is identical for \(s\) and \(s'\). Hence it is
\((0,0)\)-DP and satisfies \eqref{eq:step-004-7}, including the events containing \(g_0\) and
those not containing it. Thus \(K^0\in\mathcal K_{N,M}\).

If \(B\) is any randomized map, define
\(K_{s,g}=\Pr[B(s)=g]\). Row-stochasticity is \eqref{eq:step-004-6}, and its all-event DP
inequality is exactly \eqref{eq:step-004-7}, so a source-private \(B\) gives a point of
\(\mathcal K_{N,M}\). Conversely, given \(K\in\mathcal K_{N,M}\), on input
\(s\) draw \(g\) from the finite row distribution \((K_{s,g})_g\). This
defines a randomized unrestricted map \(B_K\), and \eqref{eq:step-004-7} proves its exact
source-cap privacy for every event. The full set \(\mathcal G_N\) includes
every improper, nonmonotone hypothesis.

Only the output law on each input matters to both privacy and the one-call
expected loss. Therefore arbitrary internal randomization is encoded without
loss. Also, if a public object is hardwired into a learner's code, the
resulting map still has one such finite kernel and belongs to
\(\mathcal K_{N,M}\) whenever it satisfies source privacy. If a map is
\((\varepsilon,\delta)\)-DP with
\(0\le\varepsilon\le0.1\) and \(0\le\delta\le\bar\delta_M\), its event
inequality is at most the right side of \eqref{eq:step-004-7}, so its kernel also belongs to
\(\mathcal K_{N,M}\). Thus the polytope includes the exact cap and every
stronger privacy pair below it. \end{proof}

\begin{lemma}[Continuous payoff geometry and uniform strict game value]\label{lem:step-004-payoff-value}
Under Lemma~\ref{lem:step-004-kernel-polytope} and Proposition~\ref{prop:step-003-expected-hardness}, fix integers
\(N,M\) satisfying \eqref{eq:step-004-hard-regime}. Give \([N+1]\) the discrete
topology and \(\Delta([N])\) its Euclidean simplex topology, and set
\(\mathcal I_N=[N+1]\times\Delta([N])\). Then:

\begin{enumerate}
\item \(\mathcal I_N\) is compact.
\item The expected one-block loss \(\ell_{N,M}(K;t,Q)\) admits the finite
   coefficient representation
   \[
   \ell_{N,M}(K;t,Q)=\langle K,a_{N,M}(t,Q)\rangle.
   \]
   It is affine and
   continuous in \(K\), and jointly continuous in \((K,t,Q)\). For fixed
   \(K,t\), its dependence on \(Q\) is generally polynomial of degree at
   most \(M+1\), so only continuity, not affinity in \((t,Q)\), is claimed.
\item The function
   \[
   f_{N,M}(K):=\max_{(t,Q)\in\mathcal I_N}
      \ell_{N,M}(K;t,Q)
   \]
   is continuous on \(\mathcal K_{N,M}\), and the attained value
   \[
   v_{N,M}:=\min_{K\in\mathcal K_{N,M}}f_{N,M}(K)
   \tag{S4.8}\label{eq:step-004-8}
   \]
   satisfies \(v_{N,M}>\eta\).
\end{enumerate}
\end{lemma}

\begin{proof}
The simplex \(\Delta([N])\) is closed and bounded in \(\mathbb R^N\), and
\([N+1]\) is finite. Their product \(\mathcal I_N\) is compact. This includes
every boundary face and every point mass of the simplex.

Write an arbitrary labeled dataset as

\[
s=((x_1,y_1),\ldots,(x_M,y_M))\in\mathcal S_{N,M}.
\]

For \((t,Q)\in\mathcal I_N\) and \(g\in\mathcal G_N\), define

\[
p_{t,Q}(s)
:=\left(\prod_{j=1}^M Q(x_j)\right)
  \mathbf 1\{y_j=\tau_t(x_j)\text{ for every }j\},
\tag{S4.9}\label{eq:step-004-9}
\]

and

\[
r_{t,Q}(g)
:=R_Q(g,\tau_t)
=\sum_{x=1}^N Q(x)\mathbf 1\{g(x)\ne\tau_t(x)\}.
\tag{S4.10}\label{eq:step-004-10}
\]

Thus \(p_{t,Q}\) is exactly the mass function of
\((Q^{\tau_t})^M\), including zero mass on every label-inconsistent
dataset. Define the coefficient vector \(a_{N,M}(t,Q)\in\mathbb R^D\) by

\[
a_{N,M}(t,Q)_{s,g}:=p_{t,Q}(s)r_{t,Q}(g).
\tag{S4.11}\label{eq:step-004-11}
\]

For the learner represented by \(K\), direct expansion of the sample and
output expectations gives

\[
\begin{aligned}
\ell_{N,M}(K;t,Q)
&:=\mathbb E_{\substack{S\sim(Q^{\tau_t})^M\\g\sim B_K(S)}}
       R_Q(g,\tau_t)\\
&=\sum_{s\in\mathcal S_{N,M}}p_{t,Q}(s)
  \sum_{g\in\mathcal G_N}K_{s,g}r_{t,Q}(g)\\
&=\sum_{s,g}K_{s,g}a_{N,M}(t,Q)_{s,g}
=\langle K,a_{N,M}(t,Q)\rangle.
\end{aligned}
\tag{S4.12}\label{eq:step-004-12}
\]

For fixed \(t\), every coordinate in \eqref{eq:step-004-11} is a polynomial in the coordinates
of \(Q\): \eqref{eq:step-004-9} has degree \(M\) and \eqref{eq:step-004-10} has degree one. Hence
\(a_{N,M}\) is continuous on each simplex copy and therefore on the finite
disjoint union \(\mathcal I_N\). Equation \eqref{eq:step-004-12} is linear, hence affine and
continuous, in \(K\), and is jointly continuous in \((K,t,Q)\). The product
structure in \eqref{eq:step-004-9} is precisely why no affine dependence on raw \(Q\) is
asserted or needed.

The compact coefficient image
\(a_{N,M}(\mathcal I_N)\) is bounded. Therefore, for any
\(K,K'\in\mathcal K_{N,M}\),

\[
\begin{aligned}
|f_{N,M}(K)-f_{N,M}(K')|
&\le
\max_{(t,Q)\in\mathcal I_N}
  |\langle K-K',a_{N,M}(t,Q)\rangle|\\
&\le \|K-K'\|_2
\max_{(t,Q)\in\mathcal I_N}\|a_{N,M}(t,Q)\|_2.
\end{aligned}
\tag{S4.13}\label{eq:step-004-13}
\]

Thus \(f_{N,M}\) is continuous. Its inner maximum is attained because
\(\mathcal I_N\) is compact, and its minimum in \eqref{eq:step-004-8} is attained because
\(\mathcal K_{N,M}\) is compact.

Now fix any \(K\in\mathcal K_{N,M}\). By
Lemma~\ref{lem:step-004-kernel-polytope}, \(K\) defines an arbitrary
randomized unrestricted source-private map \(B_K\).  Proposition
\ref{prop:step-003-expected-hardness} therefore supplies some
\((t,Q)\in\mathcal I_N\) such that
\[
\ell_{N,M}(K;t,Q)>\eta.
\]
Hence

\[
f_{N,M}(K)>\eta
\quad\text{for every }K\in\mathcal K_{N,M}.
\tag{S4.14}\label{eq:step-004-14}
\]

Let \(K_*\) attain the minimum in \eqref{eq:step-004-8}. Applying \eqref{eq:step-004-14} to this actual
minimizer, rather than taking an unattained infimum, yields

\[
v_{N,M}=f_{N,M}(K_*)>\eta.
\tag{S4.15}\label{eq:step-004-15}
\]

This is the required compactness upgrade from learner-by-learner strict
hardness to one uniform strict game value. \end{proof}

\begin{lemma}[Finite-dimensional strong separation]\label{lem:step-004-strong-separation}
Let $A\subset\mathbb R^m$ be nonempty, compact, and convex, and let
$D\subset\mathbb R^m$ be nonempty, closed, and convex.  If
$A\cap D=\varnothing$, then there are a nonzero
$\lambda\in\mathbb R^m$ and $c\in\mathbb R$ such that, after orienting
$\lambda$ toward $A$,
\[
\sup_{z\in D}\langle\lambda,z\rangle
<c<
\inf_{u\in A}\langle\lambda,u\rangle.
\]
\end{lemma}

\begin{proof}
The Euclidean distance between $A$ and $D$ is positive.  Indeed, otherwise
there would be $a_j\in A$ and $z_j\in D$ with
$\lVert a_j-z_j\rVert_2\to0$; compactness gives a subsequence
$a_j\to a_0\in A$, hence $z_j\to a_0$, and closedness would put
$a_0\in D$.  The distance is attained at some
$(a_0,z_0)\in A\times D$: a minimizing sequence in $A$ has a convergent
subsequence, and the corresponding points of $D$ are bounded and therefore
have a convergent subsequence whose limit lies in $D$.

Put $\lambda=a_0-z_0\ne0$.  For every $a\in A$, convexity permits the
segment $a_0+t(a-a_0)$, and minimality at $t=0$ gives
\[
\langle\lambda,a-a_0\rangle\ge0.
\]
Likewise, for every $z\in D$, the segment $z_0+t(z-z_0)$ and minimality
give
\[
\langle\lambda,z-z_0\rangle\le0.
\]
Consequently
\[
\sup_{z\in D}\langle\lambda,z\rangle
\le\langle\lambda,z_0\rangle
<\langle\lambda,a_0\rangle
\le\inf_{a\in A}\langle\lambda,a\rangle,
\]
where the strict gap is
$\langle\lambda,a_0-z_0\rangle=\lVert\lambda\rVert_2^2>0$.
Any scalar strictly between the middle two quantities is a valid $c$.
\end{proof}

\begin{lemma}[Compact bilinear minimax via finite-dimensional separation]\label{lem:step-004-compact-minimax}
Let \(X\) and \(Y\) be nonempty compact convex subsets of finite-dimensional
real vector spaces, and let \(p:X\times Y\to\mathbb R\) be continuous and
affine in each argument (in particular, a restricted bilinear pairing).
Then both outer extrema below are attained and

\[
\min_{x\in X}\max_{y\in Y}p(x,y)
=
\max_{y\in Y}\min_{x\in X}p(x,y).
\tag{S4.16}\label{eq:step-004-16}
\]
\end{lemma}

\begin{proof}
Continuity on the compact product implies that
\(x\mapsto\max_{y\in Y}p(x,y)\) is continuous; for example, uniform
continuity of \(p\) bounds the change of the maximum by the same uniform
modulus. Hence

\[
\alpha:=\min_{x\in X}\max_{y\in Y}p(x,y)
\]

is attained. The elementary weak minimax inequality gives

\[
\max_{y\in Y}\min_{x\in X}p(x,y)\le\alpha.
\tag{S4.17}\label{eq:step-004-17}
\]

We prove the reverse inequality by separation. Fix any real \(r<\alpha\).
For each \(x\in X\), compactness of \(Y\) supplies some \(y_x\in Y\) with
\(p(x,y_x)>r\). Continuity in \(x\) gives a neighborhood on which this same
strict inequality holds. Compactness of \(X\) yields a finite subcover, so
there are \(y_1,\ldots,y_m\in Y\) such that

\[
\max_{1\le j\le m}p(x,y_j)>r
\quad\text{for every }x\in X.
\tag{S4.18}\label{eq:step-004-18}
\]

Define the affine continuous map

\[
F:X\to\mathbb R^m,
\qquad
F(x)=(p(x,y_1),\ldots,p(x,y_m)),
\]

and the closed convex lower orthant

\[
\mathcal D_r:=(-\infty,r]^m.
\]

The image \(F(X)\) is compact and convex.  Equation
\eqref{eq:step-004-18} says that
\[
F(X)\cap\mathcal D_r=\varnothing.
\]
Applying
Lemma~\ref{lem:step-004-strong-separation} provides a nonzero
\(\lambda\in\mathbb R^m\) oriented so that

\[
\sup_{z\in\mathcal D_r}\langle\lambda,z\rangle
<
\inf_{x\in X}\langle\lambda,F(x)\rangle.
\tag{S4.19}\label{eq:step-004-19}
\]

Every coordinate of \(\lambda\) must be nonnegative: if
\(\lambda_j<0\), sending the \(j\)-th coordinate of \(z\in\mathcal D_r\)
to \(-\infty\) makes the left side of \eqref{eq:step-004-19} infinite. Normalize the nonzero
nonnegative vector so that \(\sum_j\lambda_j=1\). Then the left side of
\eqref{eq:step-004-19} is exactly \(r\), attained at \((r,\ldots,r)\), and hence

\[
\inf_{x\in X}\sum_{j=1}^m\lambda_jp(x,y_j)>r.
\tag{S4.20}\label{eq:step-004-20}
\]

Convexity of \(Y\) puts \(y_r:=\sum_j\lambda_jy_j\) in \(Y\), while
affinity in the second argument turns \eqref{eq:step-004-20} into

\[
\min_{x\in X}p(x,y_r)>r.
\tag{S4.21}\label{eq:step-004-21}
\]

Choose \(r_q\uparrow\alpha\). Compactness of \(Y\) gives a convergent
subsequence \(y_{r_q}\to y_*\in Y\). For every fixed \(x\in X\), \eqref{eq:step-004-21} and
continuity give

\[
p(x,y_*)=\lim_q p(x,y_{r_q})\ge\lim_q r_q=\alpha.
\]

Thus \(\min_xp(x,y_*)\ge\alpha\). Together with \eqref{eq:step-004-17}, this proves \eqref{eq:step-004-16} and
shows that the right-hand maximum is attained at \(y_*\). \end{proof}

\begin{lemma}[Compact coefficient hull and exact finite barycenters]\label{lem:step-004-finite-barycenter}
Let \(I\) be a nonempty compact space and let \(a:I\to\mathbb R^D\) be
continuous. Then

\[
\mathcal C:=\operatorname{conv}(a(I))
\]

is compact and convex. Moreover, every \(c\in\mathcal C\) has an exact
representation

\[
c=\sum_{j=1}^r\lambda_j a(i_j),
\qquad
1\le r\le D+1,
\quad
\lambda_j\ge0,
\quad
\sum_{j=1}^r\lambda_j=1,
\tag{S4.22}\label{eq:step-004-22}
\]

with \(i_j\in I\). Consequently, for every \(K\in\mathbb R^D\),

\[
\langle K,c\rangle
=\sum_{j=1}^r\lambda_j\langle K,a(i_j)\rangle
\tag{S4.23}\label{eq:step-004-23}
\]

exactly.
\end{lemma}

\begin{proof}
Convexity is immediate from the definition. For the support bound, start
from any finite convex representation
\(c=\sum_{j=1}^q\lambda_ja(i_j)\), which exists by the definition of convex
hull, and first remove every zero-weight term. Thus all remaining
\(\lambda_j>0\). If \(q>D+1\), the points
\(a(i_1),\ldots,a(i_q)\) are affinely dependent. Hence there are scalars
\(\theta_j\), not all zero, such that

\[
\sum_{j=1}^q\theta_j=0,
\qquad
\sum_{j=1}^q\theta_ja(i_j)=0.
\tag{S4.24}\label{eq:step-004-24}
\]

After changing all signs if needed, at least one \(\theta_j>0\). Put

\[
\rho:=\min_{j:\theta_j>0}\frac{\lambda_j}{\theta_j}.
\]

Then \(\lambda'_j=\lambda_j-\rho\theta_j\) are nonnegative, sum to one,
represent the same \(c\) by \eqref{eq:step-004-24}, and at least one formerly positive
coefficient becomes zero. Removing zero coefficients and iterating gives
\eqref{eq:step-004-22}. This is the affine-dependence proof of Caratheodory's theorem.

The representation bound also proves compactness. Pad every representation
to exactly \(D+1\) terms using zero weights. Then \(\mathcal C\) is the
image of the compact set

\[
\Delta([D+1])\times I^{D+1}
\]

under the continuous map
\((\lambda,i_1,\ldots,i_{D+1})\mapsto
\sum_j\lambda_ja(i_j)\). Thus \(\mathcal C\) is compact. Equation \eqref{eq:step-004-23}
follows from bilinearity of the Euclidean pairing and introduces no
approximation. \end{proof}

\begin{proposition}[Finite public hard prior for source-private one-block kernels]\label{prop:step-004-finite-hard-prior}
Let \(b_*,d_*>0\) and \(N_*\ge2\) be the constants supplied by Proposition~\ref{prop:step-003-expected-hardness}. Under that proposition and
Lemmas~\ref{lem:step-004-kernel-polytope}--\ref{lem:step-004-finite-barycenter},
for every pair of integers \(N,M\) satisfying
\eqref{eq:step-004-hard-regime}, there
is a probability law \(\mu_{N,M}\) supported on at most

\[
D+1=(2N)^M2^N+1
\tag{S4.25}\label{eq:step-004-25}
\]

pairs in \([N+1]\times\Delta([N])\) such that every randomized unrestricted
\((0.1,d_*/(M^2\log M))\)-DP one-block learner
$B:([N]\times\{0,1\})^M\to\{0,1\}^{[N]}$ satisfies
\[
\mathbb E_{(t,Q)\sim\mu_{N,M}}
\mathbb E_{\substack{S\sim(Q^{\tau_t})^M\\g\sim B(S)}}
R_Q(g,\tau_t)>\eta=2^{-8}.
\]
The prior is selected independently of the learner and may be revealed to it.
\end{proposition}

\begin{proof}
Use the coefficient map in \eqref{eq:step-004-11} and define

\[
\mathcal A_{N,M}:=a_{N,M}(\mathcal I_N),
\qquad
\mathcal C_{N,M}:=\operatorname{conv}(\mathcal A_{N,M}).
\]

Lemma~\ref{lem:step-004-finite-barycenter} makes
\(\mathcal C_{N,M}\) a nonempty compact convex subset of \(\mathbb R^D\).
For every fixed \(K\), linearity of the pairing gives

\[
\max_{c\in\mathcal C_{N,M}}\langle K,c\rangle
=
\max_{a\in\mathcal A_{N,M}}\langle K,a\rangle
=f_{N,M}(K).
\tag{S4.26}\label{eq:step-004-26}
\]

Indeed, a convex combination cannot exceed the largest paired coefficient,
and \(\mathcal A_{N,M}\subseteq\mathcal C_{N,M}\) gives the reverse
inequality.

Apply Lemma~\ref{lem:step-004-compact-minimax} with

\[
X=\mathcal K_{N,M},
\qquad
Y=\mathcal C_{N,M},
\qquad
p(K,c)=\langle K,c\rangle.
\]

The first set is nonempty compact convex by
Lemma~\ref{lem:step-004-kernel-polytope}; the second is nonempty compact
convex by Lemma~\ref{lem:step-004-finite-barycenter}; and the payoff is
continuous bilinear. Thus all hypotheses of the separation-based minimax
theorem are discharged. Equations \eqref{eq:step-004-8}, \eqref{eq:step-004-15}, and \eqref{eq:step-004-26} yield the exact
quantifier swap

\[
\begin{aligned}
v_{N,M}
&=\min_{K\in\mathcal K_{N,M}}
  \max_{c\in\mathcal C_{N,M}}\langle K,c\rangle\\
&=\max_{c\in\mathcal C_{N,M}}
  \min_{K\in\mathcal K_{N,M}}\langle K,c\rangle
>\eta.
\end{aligned}
\tag{S4.27}\label{eq:step-004-27}
\]

Let \(c_*\in\mathcal C_{N,M}\) attain the maximum in \eqref{eq:step-004-27}. By
Lemma~\ref{lem:step-004-finite-barycenter}, there are
\(r\le D+1\), pairs \((t_j,Q_j)\in\mathcal I_N\), and nonnegative weights
\(\lambda_j\) summing to one such that

\[
c_*=\sum_{j=1}^r\lambda_j a_{N,M}(t_j,Q_j).
\tag{S4.28}\label{eq:step-004-28}
\]

Remove zero weights and merge repeated pairs, and define the finite law

\[
\mu_{N,M}:=\sum_{j=1}^r\lambda_j\,\delta_{(t_j,Q_j)}.
\tag{S4.29}\label{eq:step-004-29}
\]

For every \(K\in\mathcal K_{N,M}\), equations \eqref{eq:step-004-12}, \eqref{eq:step-004-27}, and \eqref{eq:step-004-28} give

\[
\begin{aligned}
\mathbb E_{(t,Q)\sim\mu_{N,M}}\ell_{N,M}(K;t,Q)
&=\sum_{j=1}^r\lambda_j
  \langle K,a_{N,M}(t_j,Q_j)\rangle\\
&=\langle K,c_*\rangle\\
&\ge\min_{K'\in\mathcal K_{N,M}}\langle K',c_*\rangle\\
&=v_{N,M}>\eta.
\end{aligned}
\tag{S4.30}\label{eq:step-004-30}
\]

The first equality is the exact finite-prior expectation, not an
approximation to a continuum. Lemma~\ref{lem:step-004-kernel-polytope}
identifies every source-private randomized unrestricted \(B\) with one such
\(K\), so \eqref{eq:step-004-30} is precisely the hard-prior inequality in
the proposition and in \eqref{eq:step-004-hard-prior-interface}.

The construction of \(\mu_{N,M}\) uses the whole fixed strategy set
\(\mathcal K_{N,M}\), not any selected learner. Once \eqref{eq:step-004-29} is fixed and made
public, a learner whose program refers to \eqref{eq:step-004-29} still induces some kernel
\(K\in\mathcal K_{N,M}\) if it is source-private; \eqref{eq:step-004-30} was proved for every
such kernel after the prior was fixed. Hence prior awareness does not change
the quantifier order.

Finally, no interiority was used. The compact instance space contains
\(t=1,N+1\), every boundary distribution, and every point-mass \(Q\).
Equations \eqref{eq:step-004-9}--\eqref{eq:step-004-12} remain polynomial and exact when some \(Q(x)=0\).
Different thresholds that agree on the support of a point mass may give
duplicate coefficient vectors, but duplicates neither invalidate the convex
hull nor change \eqref{eq:step-004-30}. \end{proof}
The compact-kernel and separation arguments culminate in
Proposition~\ref{prop:step-004-finite-hard-prior}: one finite law is chosen
before the universal private-kernel quantifier, and it gives a strict
prior-averaged loss gap for every such kernel.

\subsection{Candidate and simulator calibration}\label{app:step-005}
Fix the constants from Proposition~\ref{prop:step-002-wrapper} and set
\begin{equation}
a=\frac{b_*}{16},
\qquad c_\delta=d_*,
\qquad \varepsilon_0=0.1,
\qquad \alpha_0=\beta_0=2^{-13}.
\tag{S5.C1}\label{eq:calibration}
\end{equation}
For local formulas write $L_N=\log_2^*N$ and $M=m_{n,k}$.
\begin{lemma}[Universal calibration and fixed log-star threshold]\label{lem:step-005-calibration}
Under the universal constants \(b_*,d_*>0\) and
\(N_*\in\mathbb Z_{\ge2}\) supplied by Proposition~\ref{prop:step-002-wrapper}, define

\[
q_*:=\left\lceil\frac{16}{b_*}\right\rceil,
\qquad
N_0:=\max\left\{N_*,\operatorname{Tow}_2(q_*)\right\},
\tag{S5.1}\label{eq:step-005-1}
\]

and choose the constants in \eqref{eq:calibration}. Then

\[
a,c_\delta,\varepsilon_0>0,\qquad
\alpha_0,\beta_0\in(0,1/2),\qquad
N_0\in\mathbb Z_{\ge2},
\]

all choices are universal and independent of \(k,N,n,\varepsilon,\delta\),
and every \(N\ge N_0\) satisfies

\[
N\ge N_*,
\qquad
b_*\log_2^*N\ge16.
\tag{S5.2}\label{eq:step-005-2}
\]

In particular, \eqref{eq:step-005-2} holds at the boundary \(N=N_0\).
\end{lemma}

\begin{proof}
Because \(b_*>0\), \(q_*\) is an integer at least one and

\[
q_*\ge\frac{16}{b_*}.
\tag{S5.3}\label{eq:step-005-3}
\]

The tower convention gives

\[
\log_2^*\operatorname{Tow}_2(q)=q
\qquad(q\in\mathbb Z_{\ge1}).
\tag{S5.4}\label{eq:step-005-4}
\]

Indeed, the claim is immediate for \(q=1\), since
\(\operatorname{Tow}_2(1)=2\). If it holds at \(q\), then the first
base-two logarithm of
\(\operatorname{Tow}_2(q+1)=2^{\operatorname{Tow}_2(q)}\) is
\(\operatorname{Tow}_2(q)\), so exactly one further iteration is needed.

More directly for the comparison needed here, if
\(N\ge\operatorname{Tow}_2(q_*)\), then for every
\(0\le j<q_*\), monotonicity of the iterated base-two logarithm on the
positive range traversed here gives

\[
\log_2^{(j)}N
\ge
\log_2^{(j)}\operatorname{Tow}_2(q_*)
>1.
\]

Hence \(L_N=\log_2^*N\ge q_*\). Combining this with \eqref{eq:step-005-3} yields

\[
b_*L_N\ge b_*q_*\ge16.
\tag{S5.5}\label{eq:step-005-5}
\]

Definition \eqref{eq:step-005-1} also gives \(N\ge N_*\), including when \(N=N_0\).
The remaining assertions follow from the positivity of
\(b_*,d_*\), the exact choices
\(a=b_*/16\), \(c_\delta=d_*\), \(\varepsilon_0=0.1\), and
\(2^{-13}\in(0,1/2)\). Since \(b_*,d_*,N_*\) are universal by
Proposition~\ref{prop:step-002-wrapper},
so are \(q_*\), \(N_0\), and every constant in
\eqref{eq:calibration}. \end{proof}

\begin{lemma}[Strict domination of the exact simulated budget]\label{lem:step-005-budget}
Under Assumption~\ref{assump:candidate-regime} and
Lemma~\ref{lem:step-005-calibration}, if

\[
n<akL_N
\qquad\text{with}\qquad a=\frac{b_*}{16},
\]

then the exact setting-defined budget

\[
M=\max\left\{8,\left\lceil\frac{4n}{k}\right\rceil\right\}
\]

is an integer satisfying

\[
8\le M<b_*L_N.
\tag{S5.6}\label{eq:step-005-6}
\]

This remains strict when the maximum is attained by the floor-eight
branch, when \(4n/k\) is an integer, and at \(N=N_0\).
\end{lemma}

\begin{proof}
The contradiction hypothesis and \(k>0\) give the strict inequality

\[
\frac{4n}{k}
<4aL_N
=\frac14 b_*L_N.
\tag{S5.7}\label{eq:step-005-7}
\]

For every real \(x\), including integral \(x\),
\(\lceil x\rceil<x+1\). By \eqref{eq:step-005-5}, \(1\le b_*L_N/16\). Therefore

\[
\begin{aligned}
\left\lceil\frac{4n}{k}\right\rceil
&<\frac{4n}{k}+1\\
&<\frac14b_*L_N+1\\
&\le\frac{5}{16}b_*L_N\\
&<b_*L_N.
\end{aligned}
\tag{S5.8}\label{eq:step-005-8}
\]

The fixed lower branch is also strictly below the source threshold:

\[
8<16\le b_*L_N.
\tag{S5.9}\label{eq:step-005-9}
\]

Both entries of the maximum are strictly below \(b_*L_N\), so their
maximum is strictly below it. Since both entries are integers and the first
is 8, \(M\in\mathbb Z_{\ge8}\), proving \eqref{eq:step-005-6}.

At the requested smallest tuple \((n,k)=(1,2)\),

\[
\left\lceil\frac{4n}{k}\right\rceil=\lceil2\rceil=2,
\qquad M=8.
\]

Whenever the local contradiction hypothesis holds, \eqref{eq:step-005-9} proves the strict
hard-regime inequality for this tuple as well. At \(N=N_0\), \eqref{eq:step-005-5} still
holds, so no asymptotic or strictly-larger-domain exception is being used.
\end{proof}

\begin{lemma}[Candidate-conjunction preservation and strict source privacy cap]\label{lem:step-005-privacy-cap}
Under Assumption~\ref{assump:candidate-regime}, with
\(\varepsilon_0=0.1\), \(0<c_\delta\le d_*\), and the exact
\(M=m_{n,k}\ge8\), one has simultaneously

\[
0<\varepsilon\le0.1,
\tag{S5.10}\label{eq:step-005-10}
\]

\[
0<\delta\le\frac{1}{n\log(n+1)},
\tag{S5.11}\label{eq:step-005-11}
\]

\[
0<\delta\le\frac{c_\delta}{M^2\log(M+1)}
<\frac{d_*}{M^2\log M}.
\tag{S5.12}\label{eq:step-005-12}
\]

Thus both original \(\delta\)-conjuncts are retained, while the second
strictly implies the source cap.
\end{lemma}

\begin{proof}
Equation \eqref{eq:step-005-10} is the candidate \(\varepsilon\)-condition with the fixed
choice \(\varepsilon_0=0.1\). Since an upper bound by a minimum is an upper
bound by each entry, the candidate \(\delta\)-condition gives \eqref{eq:step-005-11} and the
first inequality in \eqref{eq:step-005-12} separately. Their denominators are well-defined
and positive: \(n\ge1\) gives \(\log(n+1)\ge\log2>0\), while \(M\ge8\)
gives

\[
\log(M+1)>\log M\ge\log8>0.
\tag{S5.13}\label{eq:step-005-13}
\]

The first inequality in \eqref{eq:step-005-13} is strict because the natural logarithm is
strictly increasing and \(M+1>M\). Taking positive reciprocals reverses
that strict comparison. Using \(0<c_\delta\le d_*\),

\[
\frac{c_\delta}{M^2\log(M+1)}
<\frac{c_\delta}{M^2\log M}
\le\frac{d_*}{M^2\log M},
\tag{S5.14}\label{eq:step-005-14}
\]

which proves the strict part of \eqref{eq:step-005-12}. Notice that \eqref{eq:step-005-11} was not used to
derive \eqref{eq:step-005-14} and has not been replaced by it.

At the exact floor boundary \(M=8\), \eqref{eq:step-005-12} reads

\[
0<\delta\le\frac{c_\delta}{64\log9}
<\frac{d_*}{64\log8}.
\tag{S5.15}\label{eq:step-005-15}
\]

At \((n,k)=(1,2)\), where Lemma~\ref{lem:step-005-budget} computes
\(M=8\), the full primitive conjunction is explicitly

\[
0<\delta\le\frac1{\log2}
\quad\text{and}\quad
0<\delta\le\frac{c_\delta}{64\log9},
\]

and \eqref{eq:step-005-15} supplies the source cap. \end{proof}

\begin{proposition}[Hard-regime and candidate-parameter certificate]\label{prop:step-005-certificate}
Under Assumption~\ref{assump:candidate-regime}, Proposition~\ref{prop:step-002-wrapper}, and the local contradiction
hypothesis \(n<ak\log_2^*N\), the universal choices in
\eqref{eq:calibration} and \eqref{eq:step-005-1} imply
the following complete certificate.  The exact \(M=m_{n,k}\) satisfies

\[
N\ge N_*,
\qquad 8\le M<b_*\log_2^*N,
\qquad 0<\varepsilon\le0.1,
\]

\[
0<\delta\le\frac1{n\log(n+1)},
\qquad
0<\delta\le\frac{c_\delta}{M^2\log(M+1)}
<\frac{d_*}{M^2\log M}.
\]

The conclusion includes \(N=N_0\), \(M=8\), and, when the contradiction
hypothesis is active, \((n,k)=(1,2)\).
\end{proposition}

\begin{proof}
Lemma~\ref{lem:step-005-calibration} fixes the universal theorem constants
once and for all and gives \(N\ge N_*\) and \(b_*L_N\ge16\), including at
\(N=N_0\). Lemma~\ref{lem:step-005-budget} applies the local contradiction
hypothesis to the exact ceiling-and-maximum definition of \(M\), proving
\(M\in\mathbb Z_{\ge8}\) and the strict inequality \(M<b_*L_N\).
Lemma~\ref{lem:step-005-privacy-cap} gives \(0<\varepsilon\le0.1\), keeps
both candidate-$\delta$ inequalities, and proves the stronger strict
source-cap membership \(0<\delta<d_*/(M^2\log M)\).
The explicit boundary computations \eqref{eq:step-005-9} and \eqref{eq:step-005-15}, together with the
\((n,k)=(1,2)\) calculation after \eqref{eq:step-005-9}, prove the stated boundary clauses.
\end{proof}
Proposition~\ref{prop:step-005-certificate} therefore retains both
candidate-$\delta$ conditions while placing the same exact integer $M$ in
the strict one-block hard regime under the negated target bound.

\subsection{Pointwise PAC expectation control}\label{app:step-006}
\begin{lemma}[Every fixed product vector is a realizable instance]\label{lem:step-006-fixed-product-realizable}
Under the basic setting definitions, fix integers \(k\ge2\), \(N\ge2\)
and a deterministic vector

\[
\bigl((t_i,Q_i)\bigr)_{i=1}^k
\in\bigl([N+1]\times\Delta([N])\bigr)^k.
\]

Then

\[
P_{\boldsymbol Q}(i,x):=\frac1kQ_i(x)
\qquad((i,x)\in X_{k,N})
\]

is a probability distribution on \(X_{k,N}\),
\(\boldsymbol t=(t_1,\ldots,t_k)\in[N+1]^k\) defines
\(c_{\boldsymbol t}\in C_{k,N}\), and
\(P_{\boldsymbol Q}^{c_{\boldsymbol t}}\) is the fixed realizable
labeled law generated by \(P_{\boldsymbol Q}\) and that concept. These
conclusions also hold when any \(t_i\in\{1,N+1\}\), when any \(Q_i\) is
a point mass, or, more generally, when any \(Q_i\) lies on a boundary face
of \(\Delta([N])\).
\end{lemma}

\begin{proof}
For every \((i,x)\), nonnegativity of \(Q_i(x)\) gives
\(P_{\boldsymbol Q}(i,x)\ge0\). Since each \(Q_i\) is a probability
distribution on \([N]\),

\[
\sum_{(i,x)\in X_{k,N}}P_{\boldsymbol Q}(i,x)
=\frac1k\sum_{i=1}^k\sum_{x=1}^NQ_i(x)
=\frac1k\sum_{i=1}^k1
=1.
\tag{S6.1}\label{eq:step-006-1}
\]

Thus \(P_{\boldsymbol Q}\) is a fixed probability distribution on the
exact domain \(X_{k,N}\). Each fixed \(t_i\) belongs to \([N+1]\), so
the displayed class definition gives

\[
c_{\boldsymbol t}(i,x)=\tau_{t_i}(x)
=\mathbf 1\{x\ge t_i\}.
\tag{S6.2}\label{eq:step-006-2}
\]

Thus the function \(c_{\boldsymbol t}\), rather than its scalar value in
\eqref{eq:step-006-2}, belongs to \(C_{k,N}\).

Drawing \(Z\sim P_{\boldsymbol Q}\) and assigning the deterministic label
\(c_{\boldsymbol t}(Z)\) is, by definition, the labeled law
\(P_{\boldsymbol Q}^{c_{\boldsymbol t}}\). Hence it is realizable by a
member of \(C_{k,N}\).

No interior-support property was used in \eqref{eq:step-006-1} or \eqref{eq:step-006-2}. If \(Q_i\) is a point
mass or has zero coordinates, its coordinates remain nonnegative and sum to
one. If \(t_i=1\), then \(\tau_{t_i}\equiv1\) on \([N]\); if
\(t_i=N+1\), then \(\tau_{t_i}\equiv0\). Both endpoint concepts belong
to the displayed class. Nonuniqueness of two thresholds on the support of a
degenerate \(Q_i\) is irrelevant: the fixed parameter
\(\boldsymbol t\) still specifies one legal concept and therefore one
fixed realizable labeled law. \end{proof}

\begin{proposition}[Pointwise PAC-to-expectation conversion]\label{prop:step-006-pointwise-pac-expectation}
Under Assumption~\ref{assump:distribution-free-realizable-pac}, Lemma~\ref{lem:step-005-calibration}, and
Lemma~\ref{lem:step-006-fixed-product-realizable}, let
\(((t_i,Q_i))_{i=1}^k\in\mathcal I_N^k\) be any fixed deterministic
vector. For the exact fixed instance
\((P_{\boldsymbol Q},c_{\boldsymbol t})\), exact iid sample size \(n\),
and arbitrary randomized output \(A(S)\in\mathcal H_{k,N}\),

\[
\begin{aligned}
&\mathbb E_{\substack{S\sim
  (P_{\boldsymbol Q}^{c_{\boldsymbol t}})^n\\ A}}
  R_{P_{\boldsymbol Q}}\bigl(A(S),c_{\boldsymbol t}\bigr)\\
&\qquad\le
  \alpha_0(1-\beta_0)+\beta_0
  \le\alpha_0+\beta_0
  =2^{-12}.
\end{aligned}
\tag{S6.3}\label{eq:step-006-3}
\]

The quantifier over the fixed vector precedes both the sample/learner
expectation in \eqref{eq:step-006-3} and any later analysis-side prior average.
\end{proposition}

\begin{proof}
Fix the vector. By
Lemma~\ref{lem:step-006-fixed-product-realizable},
\(P_{\boldsymbol Q}\) is a fixed distribution and
\(c_{\boldsymbol t}\in C_{k,N}\) realizes its labels. Therefore the
universal distribution-free PAC premise applies directly to this same
fixed pair. Define the proof-local random variable

\[
\mathcal R
:=R_{P_{\boldsymbol Q}}\bigl(A(S),c_{\boldsymbol t}\bigr),
\]

where the only randomness is
\(S\sim(P_{\boldsymbol Q}^{c_{\boldsymbol t}})^n\) and the internal
randomness of \(A\). Population 0-1 risk is a probability of disagreement,
so for every possible sample and every possible output hypothesis,

\[
0\le\mathcal R\le1.
\tag{S6.4}\label{eq:step-006-4}
\]

This bound uses no properness or monotonicity: it holds for every
\(A(S)\in\{0,1\}^{X_{k,N}}\).

Let

\[
G:=\{\mathcal R\le\alpha_0\},
\qquad p:=\Pr(G^{\mathsf c}).
\]

Assumption~\ref{assump:distribution-free-realizable-pac} gives
\(p\le\beta_0\), and \(1-\alpha_0>0\) because
\(\alpha_0\in(0,1/2)\). Splitting the expectation over the success and
failure events and using \eqref{eq:step-006-4},

\[
\begin{aligned}
\mathbb E\mathcal R
&=\mathbb E[\mathcal R\mathbf 1_G]
  +\mathbb E[\mathcal R\mathbf 1_{G^{\mathsf c}}]\\
&\le \alpha_0\Pr(G)+\Pr(G^{\mathsf c})\\
&=\alpha_0(1-p)+p\\
&=\alpha_0+(1-\alpha_0)p\\
&\le\alpha_0+(1-\alpha_0)\beta_0\\
&=\alpha_0(1-\beta_0)+\beta_0\\
&\le\alpha_0+\beta_0.
\end{aligned}
\tag{S6.5}\label{eq:step-006-5}
\]

Lemma~\ref{lem:step-005-calibration} fixes
\(\alpha_0=\beta_0=2^{-13}\), so

\[
\alpha_0+\beta_0=2^{-13}+2^{-13}=2^{-12}.
\tag{S6.6}\label{eq:step-006-6}
\]

Equations \eqref{eq:step-006-5}--\eqref{eq:step-006-6} prove \eqref{eq:step-006-3}. Because the vector was fixed before invoking
the PAC premise, the proof remains valid for every deterministic vector,
including all endpoint-threshold and degenerate-distribution cases covered
by Lemma~\ref{lem:step-006-fixed-product-realizable}. \end{proof}
Lemma~\ref{lem:step-006-fixed-product-realizable} applies at every fixed
product vector, including endpoints and degenerate block laws.  The
success/failure split in
Proposition~\ref{prop:step-006-pointwise-pac-expectation} then gives the
pointwise expectation ceiling $2^{-12}$ before any prior averaging.

\begin{proof}[Assembly of the pointwise and finite-prior consequences]
For every deterministic vector
\[
((t_i,Q_i))_{i=1}^k\in\mathcal I_N^k,
\]
first fix that vector and hence
fix $P_{\boldsymbol Q}$ and $c_{\boldsymbol t}$.  Lemma~\ref{lem:step-006-fixed-product-realizable}
discharges realizability, and
Proposition~\ref{prop:step-006-pointwise-pac-expectation} then gives
\begin{equation}
\mathbb E_{\substack{S\sim
(P_{\boldsymbol Q}^{c_{\boldsymbol t}})^n\\A}}
R_{P_{\boldsymbol Q}}(A(S),c_{\boldsymbol t})\le2^{-12}.
\tag{S6.7}\label{eq:step-006-7}
\end{equation}
The universal quantifier over the deterministic vector precedes the sample
and learner expectation in \eqref{eq:step-006-7}.

Only after this pointwise statement is established do we use the finite law
from Proposition~\ref{prop:step-004-finite-hard-prior}.  Finite averaging
of \eqref{eq:step-006-7} under
$\mu_{N,M}^{\otimes k}$ gives
\begin{equation}
\mathbb E_{((t_i,Q_i))_{i=1}^k\sim\mu_{N,M}^{\otimes k}}
\left[
\mathbb E_{\substack{S\sim
(P_{\boldsymbol Q}^{c_{\boldsymbol t}})^n\\A}}
R_{P_{\boldsymbol Q}}(A(S),c_{\boldsymbol t})
\right]
\le2^{-12}.
\tag{S6.8}\label{eq:step-006-8}
\end{equation}
Thus \eqref{eq:step-006-8} is a consequence of the fixed-vector PAC
argument, not an application of the PAC premise to a prior mixture.
\end{proof}

\subsection{Block restriction and exact risk}\label{app:step-007}
\begin{lemma}[Legality of arbitrary block restriction]\label{lem:step-007-restriction-legality}
Under the tagged-domain and output-space definitions in
Sections~\ref{sec:setup}--\ref{sec:preliminaries}, if
\(h\in\mathcal H_{k,N}\) and \(j\in[k]\), then the map
\(D_jh:[N]\to\{0,1\}\) defined by \((D_jh)(x)=h(j,x)\) belongs to the full
one-block improper output space \(\{0,1\}^{[N]}\). No monotonicity or
threshold-representation condition on \(h\) or \(D_jh\) is required.
\end{lemma}

\begin{proof}
Fix arbitrary \(h\in\mathcal H_{k,N}\), \(j\in[k]\), and \(x\in[N]\).
Then \((j,x)\in X_{k,N}\). Because
\(h\in\{0,1\}^{X_{k,N}}\), its value \(h(j,x)\) lies in \(\{0,1\}\).
Consequently \(x\mapsto h(j,x)\) is a binary-valued function on all of
\([N]\), hence an element of \(\{0,1\}^{[N]}\).

An improper learner of the threshold targets \(\{\tau_t:t\in[N+1]\}\)
may output any element of this full binary-function space. Thus legality
does not require \(D_jh\) to equal any threshold \(\tau_s\), and the argument
applies unchanged to constant, oscillating, or otherwise nonmonotone \(h\).
\end{proof}

\begin{proposition}[Exact restriction-risk identity]\label{prop:step-007-risk-identity}
Under the primitive threshold and population 0-1 risk definitions and
Lemma~\ref{lem:step-007-restriction-legality}, if
\(h\in\mathcal H_{k,N}\), \(j\in[k]\), and
\((\boldsymbol t,\boldsymbol Q)=((t_i,Q_i))_{i=1}^k\) has
\(t_i\in[N+1]\) and each \(Q_i\) a probability distribution on \([N]\),
then

\[
  \Pr_{Y\sim Q_j}\!\left[(D_jh)(Y)\ne\tau_{t_j}(Y)\right]
  =e_j(h;\boldsymbol t,\boldsymbol Q).
\]

This equality remains valid for \(t_j=1\), for \(t_j=N+1\), and for every
degenerate or non-full-support \(Q_j\).
\end{proposition}

\begin{proof}
For every \(x\in[N]\), the restriction definition and the target definition
give the pointwise identity

\[
  \mathbf1\!\left\{(D_jh)(x)\ne\tau_{t_j}(x)\right\}
  =\mathbf1\!\left\{h(j,x)\ne\tau_{t_j}(x)\right\}
  =\mathbf1\!\left\{h(j,x)\ne c_{\boldsymbol t}(j,x)\right\}.
\]

Taking the finite \(Q_j\)-weighted sum of this pointwise equality yields

\[
\begin{aligned}
  \Pr_{Y\sim Q_j}\!\left[(D_jh)(Y)\ne\tau_{t_j}(Y)\right]
  &=\sum_{x\in[N]}Q_j(x)
       \mathbf1\!\left\{(D_jh)(x)\ne\tau_{t_j}(x)\right\}\\
  &=\sum_{x\in[N]}Q_j(x)
       \mathbf1\!\left\{h(j,x)\ne c_{\boldsymbol t}(j,x)\right\}\\
  &=e_j(h;\boldsymbol t,\boldsymbol Q).
\end{aligned}
\]

There is no inequality or approximation in this calculation. When
\(t_j=1\), the target is the all-one function; when \(t_j=N+1\), it is the
all-zero function. The same pointwise equality still holds. If \(Q_j\) is a
point mass, or has zero mass on part of \([N]\), the finite weighted sum
simply discards the zero-mass coordinates and remains exact.

Finally, let \(J\) be any \([k]\)-valued random variable and let \(H\) be any
random element of \(\mathcal H_{k,N}\), with arbitrary dependence between
them. Lemma~\ref{lem:step-007-restriction-legality} and the displayed risk
identity apply separately to every realization \((J,H)=(j,h)\). Therefore
\(D_JH\) is almost surely a legal improper one-block output and

\[
  \sum_{x\in[N]}Q_J(x)
    \mathbf1\!\left\{(D_JH)(x)\ne\tau_{t_J}(x)\right\}
  =e_J(H;\boldsymbol t,\boldsymbol Q)
\]

holds pathwise. No independence, learner symmetry, or monotonicity premise
is used. \end{proof}
The restriction is legal for every improper output, and
Proposition~\ref{prop:step-007-risk-identity} identifies its one-block risk
exactly, with no projection or residual.

\subsection{The ideal hidden-coordinate experiment}\label{app:step-008}
Fix a finite public law $\mu_{N,M}$ and an arbitrary randomized learner
$A$.  Draw $J\sim\operatorname{Unif}[k]$.  Conditional on every $J=j$,
draw $\Xi_i=(T_i,Q_i)$ independently from $\mu_{N,M}$ for all $i\in[k]$;
then draw independent uniform tags $I_1,\ldots,I_n$, features
$X_\ell\sim Q_{I_\ell}$ conditionally independently, labels
$Y_\ell=\tau_{T_{I_\ell}}(X_\ell)$, and finally
$H=A(S^{\mathrm{id}})$ for
$S^{\mathrm{id}}=(((I_\ell,X_\ell),Y_\ell))_{\ell=1}^n$.
The conclusions to be proved are
\begin{align}
\boldsymbol\Xi&\sim\mu_{N,M}^{\otimes k},
&J&\perp\!\!\!\perp\boldsymbol\Xi,
\tag{S8.2}\label{eq:ideal-instance-law}\\
\mathcal L(S^{\mathrm{id}}\mid\boldsymbol\Xi)
&=(P_{\boldsymbol Q}^{c_{\boldsymbol T}})^n,
\tag{S8.3}\label{eq:ideal-sample-law}\\
\mathbb E e_J(H;\boldsymbol T,\boldsymbol Q)
&=\mathbb E R_{P_{\boldsymbol Q}}(H,c_{\boldsymbol T}).
\tag{S8.4}\label{eq:ideal-risk-identity}
\end{align}
\begin{lemma}[Hidden-coordinate instance factorization]\label{lem:step-008-instance-factorization}
Under Proposition~\ref{prop:step-004-finite-hard-prior} and the
ideal experiment above, if \(J\sim\operatorname{Unif}[k]\) and, conditional
on each \(J=j\), the coordinate instances
\(\Xi_1,\ldots,\Xi_k\) are mutually independent with their common law
\(\mu_{N,M}\), then

\[
\mathcal L(\boldsymbol\Xi\mid J=j)=\mu_{N,M}^{\otimes k}
\quad\text{for every }j\in[k],
\tag{S8.5}\label{eq:step-008-5}
\]

and consequently \eqref{eq:ideal-instance-law} holds. This remains true when \(\mu_{N,M}\) has
singleton support or repeated coordinate draws occur.
\end{lemma}

\begin{proof}
Let \(\mathcal Z=\operatorname{supp}(\mu_{N,M})\), which is a nonempty
finite set because Proposition~\ref{prop:step-004-finite-hard-prior}
supplies a finitely supported
probability law. For every \(j\in[k]\) and
\(\boldsymbol\xi=(\xi_1,\ldots,\xi_k)\in\mathcal Z^k\), the specified
conditional independence gives

\[
\begin{aligned}
\Pr\{J=j,\boldsymbol\Xi=\boldsymbol\xi\}
&=\Pr\{J=j\}
  \Pr\{\boldsymbol\Xi=\boldsymbol\xi\mid J=j\}\\
&=\frac1k\prod_{i=1}^k\mu_{N,M}(\{\xi_i\}).
\end{aligned}
\tag{S8.6}\label{eq:step-008-6}
\]

The product on the right does not depend on \(j\). Summing \eqref{eq:step-008-6} over
\(j\) gives

\[
\Pr\{\boldsymbol\Xi=\boldsymbol\xi\}
=\prod_{i=1}^k\mu_{N,M}(\{\xi_i\}),
\]

which is the product law in \eqref{eq:ideal-instance-law}. Equation \eqref{eq:step-008-6} then factors as

\[
\Pr\{J=j,\boldsymbol\Xi=\boldsymbol\xi\}
=\Pr\{J=j\}\Pr\{\boldsymbol\Xi=\boldsymbol\xi\},
\]

proving independence. If the prior has singleton support, all coordinates
are the same deterministic pair almost surely, but the same factorization
still proves both the product law and independence from \(J\). Calling one
coordinate hidden therefore creates no selection bias. \end{proof}

\begin{lemma}[Conditional i.i.d. tagged-product sample law]\label{lem:step-008-ideal-sample-law}
Under the primitive tagged-product definitions and
Lemma~\ref{lem:step-008-instance-factorization}, suppose that after
\((J,\boldsymbol\Xi)\) is generated, the tags
\(I_1,\ldots,I_n\) are drawn independently from
\(\operatorname{Unif}[k]\), independently of the preceding variables;
conditional on \((\boldsymbol\Xi,I_1,\ldots,I_n)\), the features are
independent with \(X_\ell\sim Q_{I_\ell}\); and
\(Y_\ell=\tau_{T_{I_\ell}}(X_\ell)\). Then for every realized
\(\boldsymbol\xi=((t_i,Q_i))_{i=1}^k\in\mathcal Z^k\) and every
\(j\in[k]\),

\[
\mathcal L\!\left(S^{\mathrm{id}}\mid
  \boldsymbol\Xi=\boldsymbol\xi,J=j\right)
=\left(P_{\boldsymbol Q}^{c_{\boldsymbol t}}\right)^n.
\tag{S8.7}\label{eq:step-008-7}
\]

In particular, the conditional row kernel is independent of \(j\), and \eqref{eq:ideal-sample-law}
holds.
\end{lemma}

\begin{proof}
Fix \(\boldsymbol\xi=((t_i,Q_i))_{i=1}^k\) and \(j\). For any row index
\(\ell\), any \(i\in[k]\), \(x\in[N]\), and \(y\in\{0,1\}\), the tag,
feature, and deterministic-label rules give

\[
\begin{aligned}
&\Pr\left\{((I_\ell,X_\ell),Y_\ell)=((i,x),y)
  \mid\boldsymbol\Xi=\boldsymbol\xi,J=j\right\}\\
&\qquad=\frac1k Q_i(x)\,
  \mathbf1\{y=\tau_{t_i}(x)\}\\
&\qquad=P_{\boldsymbol Q}(i,x)\,
  \mathbf1\{y=c_{\boldsymbol t}(i,x)\}\\
&\qquad=P_{\boldsymbol Q}^{c_{\boldsymbol t}}
  \bigl(\{((i,x),y)\}\bigr).
\end{aligned}
\tag{S8.8}\label{eq:step-008-8}
\]

The first line on the right uses neither \(j\) nor any property of the
learner. The pairs \((I_\ell,X_\ell)\) are mutually independent conditional
on \(\boldsymbol\Xi\): tags are drawn independently, and the conditional
feature draws use independent randomness across rows. Labels are
deterministic rowwise functions. Therefore the \(n\) labeled rows are
conditionally independent and all have the common law in \eqref{eq:step-008-8}, which proves
\eqref{eq:step-008-7}.

By Lemma~\ref{lem:step-008-instance-factorization}, the conditional law of
\(\boldsymbol\Xi\) itself is the same for every \(j\). Since \eqref{eq:step-008-7} is also the
same for every \(j\), summing over \(J\) yields \eqref{eq:ideal-sample-law}. Importantly, \eqref{eq:ideal-sample-law} is a
conditional i.i.d. statement for each fixed product instance. Unconditionally,
the rows share the random vector \(\boldsymbol\Xi\), so no stronger
unconditional i.i.d. assertion is made or needed.

Equation \eqref{eq:step-008-8} remains valid for \(t_i=1\) or \(t_i=N+1\), when the labels in
that block are constant, and when \(Q_i\) is a point mass or has zeros. No
division by \(Q_i(x)\) occurs. It also defines an exact size-\(n\) dataset
when no sampled tag equals the later-designated hidden coordinate \(J\).
\end{proof}

\begin{lemma}[Persistence of selector independence through an asymmetric
learner]\label{lem:step-008-selector-independence}
Under Lemmas~\ref{lem:step-008-instance-factorization} and
\ref{lem:step-008-ideal-sample-law}, if \(A\) is any randomized map with the
domain and codomain in Section~\ref{sec:setup} and is run on
\(S^{\mathrm{id}}\) with fresh internal randomness, then

\[
J\ \perp\!\!\!\perp\
(\boldsymbol\Xi,S^{\mathrm{id}},H).
\tag{S8.9}\label{eq:step-008-9}
\]

Equivalently, for every \(i\in[k]\),

\[
\Pr\{J=i\mid\boldsymbol\Xi,S^{\mathrm{id}},H\}=\frac1k
\quad\text{almost surely}.
\tag{S8.10}\label{eq:step-008-10}
\]

No permutation invariance, tag symmetry, properness, or determinism of
\(A\) is required.
\end{lemma}

\begin{proof}
Because all involved spaces are finite, represent the randomized learner by
its stochastic kernel

\[
K_A(h\mid s):=\Pr\{A(s)=h\},
\qquad h\in\mathcal H_{k,N}.
\]

For a fixed instance vector \(\boldsymbol\xi\), let
\(p_{\boldsymbol\xi}(s)\) denote its product-sample mass.  This is the
mass furnished by Lemma~\ref{lem:step-008-ideal-sample-law}. Equations
\eqref{eq:step-008-6}--\eqref{eq:step-008-8}, followed by the independent
learner call, give for every atom with
\(\boldsymbol\xi\in\mathcal Z^k\),

\[
\begin{aligned}
&\Pr\{J=j,\boldsymbol\Xi=\boldsymbol\xi,
       S^{\mathrm{id}}=s,H=h\}\\
&\qquad=\frac1k
  \left(\prod_{i=1}^k\mu_{N,M}(\{\xi_i\})\right)
  p_{\boldsymbol\xi}(s)K_A(h\mid s).
\end{aligned}
\tag{S8.11}\label{eq:step-008-11}
\]

Every factor after \(1/k\) is independent of \(j\). Summing \eqref{eq:step-008-11} over
\(j\) shows that it equals

\[
\Pr\{J=j\}\Pr\{\boldsymbol\Xi=\boldsymbol\xi,
S^{\mathrm{id}}=s,H=h\},
\]

which proves \eqref{eq:step-008-9}, including zero-mass atoms by omission. The finite
conditional-probability formula gives \eqref{eq:step-008-10} on every positive-mass atom.

The kernel \(K_A(h\mid s)\) may depend arbitrarily on every tag and every
row of \(s\). Thus \eqref{eq:step-008-11} does not assert that the law of \(H\) is invariant
under permutations of tags. It uses only that \(J\) is an analysis-side
variable not supplied as an argument to \(A\), and that the sample kernel
in \eqref{eq:step-008-7} does not depend on \(J\). The learner may also have the public prior
hardwired into its code; this changes \(K_A\) but introduces no \(j\)-factor.
\end{proof}

\begin{proposition}[Deterministic tagged-product risk decomposition]\label{prop:step-008-risk-decomposition}
Under the primitive product-distribution definition and Proposition~\ref{prop:step-007-risk-identity}, for every fixed
\(\boldsymbol t\in[N+1]^k\), every vector of probability distributions
\(\boldsymbol Q=(Q_1,\ldots,Q_k)\) on \([N]\), and every arbitrary
\(h\in\mathcal H_{k,N}\),

\[
\begin{aligned}
R_{P_{\boldsymbol Q}}(h,c_{\boldsymbol t})
&=\frac1k\sum_{i=1}^k e_i(h;\boldsymbol t,\boldsymbol Q)\\
&=\frac1k\sum_{i=1}^k
  R_{Q_i}(D_i h,\tau_{t_i}).
\end{aligned}
\tag{S8.12}\label{eq:step-008-12}
\]

The equality is deterministic and does not require any distributional
symmetry of \(\boldsymbol Q\) or any structural property of \(h\).
\end{proposition}

\begin{proof}
Expanding the population risk on the finite tagged domain and using
\[
P_{\boldsymbol Q}(i,x)=k^{-1}Q_i(x)
\]
gives

\[
\begin{aligned}
R_{P_{\boldsymbol Q}}(h,c_{\boldsymbol t})
&=\sum_{i=1}^k\sum_{x\in[N]}
  \frac1k Q_i(x)
  \mathbf1\{h(i,x)\ne c_{\boldsymbol t}(i,x)\}\\
&=\frac1k\sum_{i=1}^k\sum_{x\in[N]}Q_i(x)
  \mathbf1\{h(i,x)\ne\tau_{t_i}(x)\}\\
&=\frac1k\sum_{i=1}^k e_i(h;\boldsymbol t,\boldsymbol Q).
\end{aligned}
\tag{S8.13}\label{eq:step-008-13}
\]

Proposition~\ref{prop:step-007-risk-identity} identifies the
\(i\)-th inner sum with
\(R_{Q_i}(D_i h,\tau_{t_i})\), proving the second equality in \eqref{eq:step-008-12}. The
calculation remains exact for unequal block distributions, endpoint
thresholds, point masses, and arbitrary nonmonotone or tag-asymmetric
hypotheses. \end{proof}

\begin{lemma}[Finite conditional averaging]\label{lem:step-008-finite-conditional-average}
Let $\mathcal G$ be a sigma-field, let $J$ take values in $[k]$, and suppose
$\Pr(J=i\mid\mathcal G)=1/k$ almost surely for every $i\in[k]$.  If
$a_1,\ldots,a_k$ are integrable $\mathcal G$-measurable random variables,
then
\[
\mathbb E[a_J\mid\mathcal G]=\frac1k\sum_{i=1}^k a_i
\quad\text{almost surely},
\]
and consequently
\[
\mathbb E a_J=\mathbb E\left[\frac1k\sum_{i=1}^k a_i\right].
\]
\end{lemma}

\begin{proof}
The pointwise finite identity
$a_J=\sum_{i=1}^k a_i\mathbf1\{J=i\}$ and
$\mathcal G$-measurability of every $a_i$ give
\[
\mathbb E[a_J\mid\mathcal G]
=\sum_{i=1}^k a_i\mathbb E[\mathbf1\{J=i\}\mid\mathcal G]
=\sum_{i=1}^k a_i\Pr(J=i\mid\mathcal G)
=\frac1k\sum_{i=1}^k a_i.
\]
Taking expectations proves the second identity.  In the ideal experiment
all variables to which this lemma is applied lie in $[0,1]$, so
integrability is automatic.
\end{proof}

\begin{proposition}[Selected-coordinate averaging for arbitrary randomized
learners]\label{prop:step-008-selected-risk-identity}
Under Propositions~\ref{prop:step-007-risk-identity} and
\ref{prop:step-008-risk-decomposition}, and
Lemmas~\ref{lem:step-008-selector-independence} and
\ref{lem:step-008-finite-conditional-average}, for the ideal experiment
and every randomized map \(A\) specified above,

\[
\mathbb E\!\left[
e_J(H;\boldsymbol T,\boldsymbol Q)
\mid\boldsymbol\Xi,S^{\mathrm{id}},H
\right]
=R_{P_{\boldsymbol Q}}(H,c_{\boldsymbol T})
\quad\text{almost surely}.
\tag{S8.14}\label{eq:step-008-14}
\]

Consequently,

\[
\begin{aligned}
\mathbb E\,
R_{Q_J}(D_JH,\tau_{T_J})
&=\mathbb E\,e_J(H;\boldsymbol T,\boldsymbol Q)\\
&=\mathbb E\,R_{P_{\boldsymbol Q}}(H,c_{\boldsymbol T}),
\end{aligned}
\tag{S8.15}\label{eq:step-008-15}
\]

where the expectations retain all prior, sample, selector, and learner
randomness.
\end{proposition}

\begin{proof}
Conditional on \((\boldsymbol\Xi,S^{\mathrm{id}},H)\), every quantity
\(e_i(H;\boldsymbol T,\boldsymbol Q)\) is fixed. By \eqref{eq:step-008-10}, the remaining
coordinate selector is uniform, so
Lemma~\ref{lem:step-008-finite-conditional-average} gives

\[
\begin{aligned}
\mathbb E\!\left[
e_J(H;\boldsymbol T,\boldsymbol Q)
\mid\boldsymbol\Xi,S^{\mathrm{id}},H
\right]
&=\sum_{i=1}^k\frac1k
  e_i(H;\boldsymbol T,\boldsymbol Q)\\
&=R_{P_{\boldsymbol Q}}(H,c_{\boldsymbol T}),
\end{aligned}
\]

where the last equality is the pathwise application of
Proposition~\ref{prop:step-008-risk-decomposition}. Taking expectations and
using the second conclusion of
Lemma~\ref{lem:step-008-finite-conditional-average} proves
the second equality in \eqref{eq:step-008-15}, which is \eqref{eq:ideal-risk-identity}. The first equality in \eqref{eq:step-008-15} is the
pathwise Proposition~\ref{prop:step-007-risk-identity} at the realized
\((J,H,T_J,Q_J)\).

The proof averages over an independent uniform \(J\) after conditioning on
the complete realized vector, sample, and output. It never replaces
\(e_J\) by \(e_1\), never permutes \(A\), and never assumes the errors
\(e_1,\ldots,e_k\) have the same distribution. Thus deliberate tag
asymmetry of \(A\) is fully covered. For \(k=2\) and \(k=3\), the
conditional sums are exactly \((e_1+e_2)/2\) and
\((e_1+e_2+e_3)/3\), respectively. Samples in which the selected tag does
not occur require no special conditioning and remain part of the same
unconditioned identity. \end{proof}
Lemmas~\ref{lem:step-008-instance-factorization}--
\ref{lem:step-008-selector-independence} establish the exact stochastic
factorization, and
Proposition~\ref{prop:step-008-selected-risk-identity} proves
\eqref{eq:ideal-risk-identity} for arbitrary tag-asymmetric learners.

\subsection{Usage and overflow}\label{app:step-009}
Use
\begin{equation}
M=\max\left\{8,\left\lceil\frac{4n}{k}\right\rceil\right\},
\tag{S9.1}\label{eq:overflow-budget}
\end{equation}
and, for the selector and tags above, set
\begin{equation}
U=\sum_{\ell=1}^n\mathbf 1\{I_\ell=J\}.
\tag{S9.2}\label{eq:usage-count}
\end{equation}
The desired certificate is
\begin{equation}
U\sim\operatorname{Bin}(n,1/k),
\qquad \Pr\{U>M\}<2^{-9}.
\tag{S9.3}\label{eq:overflow-certificate}
\end{equation}
\begin{lemma}[Exact hidden-tag usage law]\label{lem:step-009-usage-law}
Under the ideal hidden-coordinate tag interface, for every pair of
integers \(k\ge2,n\ge1\), the count \(U\) in \eqref{eq:usage-count} satisfies, for each
\(j\in[k]\) and each \(u\in\{0,\ldots,n\}\),

\[
\Pr\{U=u\mid J=j\}
=\binom nu\left(\frac1k\right)^u
 \left(1-\frac1k\right)^{n-u}.
\tag{S9.4}\label{eq:step-009-4}
\]

Consequently \(U\sim\operatorname{Bin}(n,1/k)\) unconditionally. The law
includes \(U=0\), and its trial count is exactly \(n\).
\end{lemma}

\begin{proof}
Fix \(j\in[k]\). By the tag interface, conditional on \(J=j\),
the variables

\[
B_\ell:=\mathbf1\{I_\ell=j\},\qquad \ell\in[n],
\]

are mutually independent, and

\[
\Pr\{B_\ell=1\mid J=j\}=\Pr\{I_\ell=j\mid J=j\}=\frac1k.
\]

Thus \(U=\sum_{\ell=1}^nB_\ell\). For any fixed set of exactly \(u\)
indices, the probability that precisely those indicators equal one is
\(k^{-u}(1-1/k)^{n-u}\). The \(\binom nu\) choices are disjoint, proving
\eqref{eq:step-009-4}. The right-hand side does not depend on \(j\). Since
\(\Pr(J=j)=1/k\),

\[
\begin{aligned}
\Pr\{U=u\}
&=\sum_{j=1}^k\Pr\{J=j\}\Pr\{U=u\mid J=j\}\\
&=\binom nu\left(\frac1k\right)^u
 \left(1-\frac1k\right)^{n-u}.
\end{aligned}
\]

This is the asserted binomial mass function. In particular, no random block
instance, feature, label, learner output, or stopping rule changes the law.
\end{proof}

\begin{lemma}[Falling-factorial binomial tail bound]\label{lem:step-009-factorial-tail}
Let \(V\sim\operatorname{Bin}(n,p)\), where \(n\ge1\) and
\(p\in[0,1]\). For every integer \(v\ge0\), set \((v)_0:=1\); for every
integer \(r\ge1\), write

\[
(v)_r:=v(v-1)\cdots(v-r+1),
\]

with \((v)_r=0\) when \(r>v\). Then, for every integer \(r\ge1\),

\[
\mathbb E(V)_r=(n)_r p^r
\tag{S9.5}\label{eq:step-009-5}
\]

and

\[
\Pr\{V\ge r\}
\le \frac{(n)_r p^r}{r!}
\le \frac{(np)^r}{r!}.
\tag{S9.6}\label{eq:step-009-6}
\]

Equations \eqref{eq:step-009-5}--\eqref{eq:step-009-6} remain valid when \(r>n\), in which case both the event
and \((n)_r\) vanish.
\end{lemma}

\begin{proof}
Represent \(V=\sum_{\ell=1}^n Z_\ell\) for independent
Bernoulli\((p)\) variables. For every outcome with \(V=v\), the number of
ordered \(r\)-tuples of pairwise distinct successful trial indices is
exactly \((v)_r\). Hence the pointwise identity

\[
(V)_r
=\sum_{\substack{(\ell_1,\ldots,\ell_r)\in[n]^r\\
                    \ell_1,\ldots,\ell_r\ \mathrm{pairwise\ distinct}}}
  \prod_{s=1}^r Z_{\ell_s}
\tag{S9.7}\label{eq:step-009-7}
\]

holds, with an empty sum if \(r>n\). There are \((n)_r\) ordered tuples,
and independence gives expectation \(p^r\) for every product in \eqref{eq:step-009-7},
proving \eqref{eq:step-009-5}.

If \(V\ge r\), then

\[
(V)_r=r!\binom Vr\ge r!,
\]

whereas if \(V<r\), both \(\mathbf1\{V\ge r\}\) and \((V)_r\) are zero.
Therefore

\[
\mathbf1\{V\ge r\}\le\frac{(V)_r}{r!}
\]

pointwise. Taking expectations and using \eqref{eq:step-009-5} gives the first inequality in
\eqref{eq:step-009-6}. Finally, each of the \(r\) nonnegative factors of \((n)_r\) is at
most \(n\) when \(r\le n\), so \((n)_r\le n^r\). When \(r>n\), the
left side is zero. This proves the second inequality in every case.
\end{proof}

\begin{lemma}[Floor-eight small-mean bound]\label{lem:step-009-small-mean}
Under Lemmas~\ref{lem:step-009-usage-law} and
\ref{lem:step-009-factorial-tail}, define

\[
\lambda:=\frac nk.
\tag{S9.8}\label{eq:step-009-8}
\]

If \(\lambda\le2\), then the exact budget \eqref{eq:overflow-budget} satisfies \(M=8\) and

\[
\Pr\{U>M\}
\le\frac{\lambda^9}{9!}
\le\frac{2^9}{9!}
<2^{-9}.
\tag{S9.9}\label{eq:step-009-9}
\]

If, more specifically, \(n<k\), then \(M=8\) and

\[
\Pr\{U>M\}<\frac1{9!}<2^{-9}.
\tag{S9.10}\label{eq:step-009-10}
\]

Thus the statement includes the exact boundary \(\lambda=2\), every
occurrence of \(M=8\), and the case \(n<k\).
\end{lemma}

\begin{proof}
If \(\lambda\le2\), then \(4\lambda\le8\), so

\[
\lceil4n/k\rceil=\lceil4\lambda\rceil\le8.
\]

Equation \eqref{eq:overflow-budget} therefore gives \(M=8\) exactly. Since \(U\) is integer,

\[
\{U>M\}=\{U>8\}=\{U\ge9\}.
\tag{S9.11}\label{eq:step-009-11}
\]

Apply Lemma~\ref{lem:step-009-factorial-tail} to the exact law in
Lemma~\ref{lem:step-009-usage-law} with \(p=1/k\) and \(r=9\). Because
\(np=n/k=\lambda\),

\[
\Pr\{U>M\}\le\frac{\lambda^9}{9!}
\le\frac{2^9}{9!}.
\tag{S9.12}\label{eq:step-009-12}
\]

The final comparison in \eqref{eq:step-009-9} is strict and completely numerical:

\[
9!=362880>262144=2^{18}.
\tag{S9.13}\label{eq:step-009-13}
\]

Dividing \eqref{eq:step-009-13} by the positive number \(2^9 9!\) gives
\(2^9/9!<2^{-9}\), completing \eqref{eq:step-009-9}. If \(n<9\), the event in \eqref{eq:step-009-11} is
actually empty because \(U\le n\), and the convention in
Lemma~\ref{lem:step-009-factorial-tail} records this as \((n)_9=0\).

If \(n<k\), then \(0<\lambda<1\) and
\(0<4\lambda<4\), whence \(\lceil4\lambda\rceil\le4\) and \(M=8\).
The first inequality of \eqref{eq:step-009-9} is then strict:

\[
\Pr\{U>M\}\le\frac{\lambda^9}{9!}<\frac1{9!}.
\]

Since \(9!=362880>512=2^9\), one has \(1/9!<2^{-9}\), proving \eqref{eq:step-009-10}.
Conversely, if \(M=8\), then \eqref{eq:overflow-budget} implies
\(\lceil4\lambda\rceil\le8\), hence \(4\lambda\le8\) and
\(\lambda\le2\); thus every possible \(M=8\) case is indeed contained
in this lemma. \end{proof}

\begin{lemma}[Monotone factorial envelope]\label{lem:step-009-envelope}
For every integer \(j\ge9\), define

\[
a_j:=\frac{(j/4)^{j+1}}{(j+1)!}.
\tag{S9.14}\label{eq:step-009-14}
\]

Then \((a_j)_{j\ge9}\) is strictly decreasing and

\[
a_j\le a_9
=\frac{(9/4)^{10}}{10!}
<2^{-10}.
\tag{S9.15}\label{eq:step-009-15}
\]
\end{lemma}

\begin{proof}
Direct cancellation gives, for every integer \(j\ge9\),

\[
\frac{a_{j+1}}{a_j}
=\frac{j+1}{4(j+2)}
 \left(1+\frac1j\right)^{j+1}.
\tag{S9.16}\label{eq:step-009-16}
\]

We bound the power in \eqref{eq:step-009-16} without importing an asymptotic exponential
estimate. The binomial theorem gives

\[
\left(1+\frac1j\right)^j
=\sum_{r=0}^j \binom jr j^{-r}
=\sum_{r=0}^j \frac{(j)_r}{r!\,j^r}
\le\sum_{r=0}^j\frac1{r!}
<3.
\tag{S9.17}\label{eq:step-009-17}
\]

For the strict last inequality, \(r!\ge2^{r-1}\) for every \(r\ge2\),
with strict inequality for every \(r\ge3\), so

\[
\sum_{r=2}^{\infty}\frac1{r!}
<\sum_{r=2}^{\infty}\frac1{2^{r-1}}=1.
\]

Adding the \(r=0,1\) terms proves \eqref{eq:step-009-17}. Since \(j\ge9\),

\[
\left(1+\frac1j\right)^{j+1}
<3\left(1+\frac19\right)
=\frac{10}{3}<4.
\tag{S9.18}\label{eq:step-009-18}
\]

Substituting \eqref{eq:step-009-18} into \eqref{eq:step-009-16} yields the strict ratio bound

\[
\frac{a_{j+1}}{a_j}
<\frac{j+1}{j+2}<1.
\tag{S9.19}\label{eq:step-009-19}
\]

Thus the sequence is strictly decreasing and \(a_j\le a_9\). It remains
to verify the endpoint constant. Since \(4^{10}=2^{20}\), the inequality
\(a_9<2^{-10}\) is equivalent to

\[
9^{10}<2^{10}\,10!.
\]

Both integers are explicit:

\[
9^{10}=3486784401
<3715891200=2^{10}\,10!.
\tag{S9.20}\label{eq:step-009-20}
\]

Equations \eqref{eq:step-009-19}--\eqref{eq:step-009-20} prove \eqref{eq:step-009-15}. \end{proof}

\begin{lemma}[Ceiling-controlled large-mean bound]\label{lem:step-009-large-mean}
Under Lemmas~\ref{lem:step-009-usage-law},
\ref{lem:step-009-factorial-tail}, and
\ref{lem:step-009-envelope}, if \(\lambda=n/k>2\) and

\[
j:=\lceil4\lambda\rceil,
\tag{S9.21}\label{eq:step-009-21}
\]

then \(j\) is an integer satisfying \(j\ge9\), the exact budget is
\(M=j\), and

\[
\Pr\{U>M\}
\le\frac{\lambda^{j+1}}{(j+1)!}
\le\frac{(j/4)^{j+1}}{(j+1)!}
\le\frac{(9/4)^{10}}{10!}
<2^{-10}<2^{-9}.
\tag{S9.22}\label{eq:step-009-22}
\]

The conclusion remains valid when \(j+1>n\), when the overflow event is
empty.
\end{lemma}

\begin{proof}
The strict inequality \(\lambda>2\) gives \(4\lambda>8\). Therefore its
integer ceiling obeys

\[
j=\lceil4\lambda\rceil\ge9,
\tag{S9.23}\label{eq:step-009-23}
\]

so \eqref{eq:overflow-budget} gives \(M=\max\{8,j\}=j\). Because \(U\) is integer,

\[
\{U>M\}=\{U\ge j+1\}.
\tag{S9.24}\label{eq:step-009-24}
\]

Apply Lemma~\ref{lem:step-009-factorial-tail} with
\(p=1/k\) and the integer order \(r=j+1\). Together with
Lemma~\ref{lem:step-009-usage-law}, this gives

\[
\Pr\{U>M\}
\le\frac{(n/k)^{j+1}}{(j+1)!}
=\frac{\lambda^{j+1}}{(j+1)!}.
\tag{S9.25}\label{eq:step-009-25}
\]

If \(j+1>n\), the sharper first expression in \eqref{eq:step-009-6} has
\((n)_{j+1}=0\), so \eqref{eq:step-009-24} is empty and \eqref{eq:step-009-25} still holds. No unstated
assumption \(j+1\le n\) is needed.

The ceiling definition \eqref{eq:step-009-21} also gives \(j\ge4\lambda\), hence
\(\lambda\le j/4\). Raising nonnegative quantities to the positive integer
power \(j+1\) and dividing by \((j+1)!>0\) proves the second inequality in
\eqref{eq:step-009-22}. Lemma~\ref{lem:step-009-envelope} supplies the next two inequalities,
and \(2^{-10}<2^{-9}\) follows from \(2^{-10}=\tfrac12 2^{-9}\).
This proves the entire strict chain \eqref{eq:step-009-22}. \end{proof}

\begin{proposition}[Exact finite-budget overflow certificate]\label{prop:step-009-overflow}
Under the basic-setting ranges \(k\ge2,n\ge1\), the exact budget \eqref{eq:overflow-budget}, and
the ideal hidden-coordinate tag interface, the usage count \eqref{eq:usage-count}
satisfies \eqref{eq:overflow-certificate}. More explicitly:

\begin{enumerate}
\item If \(n/k\le2\), then \(M=8\) and \eqref{eq:step-009-9} gives the strict tail bound.
\item If \(n/k>2\), then \(M=\lceil4n/k\rceil\ge9\) and \eqref{eq:step-009-22} gives the
stronger tail bound \(\Pr(U>M)<2^{-10}\).
\item If \(k=2\) or \(k=3\), then \(M\ge n\) and \(\Pr(U>M)=0\).
\item If \(n<k\), then \(M=8\) and the strict chain \eqref{eq:step-009-10} holds.
\item If \(M=8\), necessarily \(n/k\le2\), so the first branch applies.
\end{enumerate}

No contradiction-regime, hard-prior, privacy, PAC, or learner condition is
needed for any clause.
\end{proposition}

\begin{proof}
Lemma~\ref{lem:step-009-usage-law} proves the first assertion in \eqref{eq:overflow-certificate}.
Exactly one of \(n/k\le2\) and \(n/k>2\) holds. In the first case,
Lemma~\ref{lem:step-009-small-mean} proves
\(\Pr(U>M)<2^{-9}\); in the second,
Lemma~\ref{lem:step-009-large-mean} proves the stronger
\(\Pr(U>M)<2^{-10}<2^{-9}\). This proves the second assertion in \eqref{eq:overflow-certificate}
for all \(k,n\).

For the requested small-tag boundaries, if \(k\in\{2,3\}\), then
\(4/k\ge1\), so

\[
\left\lceil\frac{4n}{k}\right\rceil\ge n
\quad\text{and hence}\quad M\ge n.
\tag{S9.26}\label{eq:step-009-26}
\]

The count \eqref{eq:usage-count} always obeys \(0\le U\le n\). Thus \(U>M\) is impossible
and its probability is zero. This direct argument covers every \(n\),
including values for which \(M=8\) and values for which the ceiling branch
exceeds eight.

The \(n<k\) and \(M=8\) clauses, including all ceiling inequalities and
their strict constants, were proved in
Lemma~\ref{lem:step-009-small-mean}. Finally, \(U=0\) never belongs to the
overflow event because \(M\ge8\); it remains an ordinary atom of the exact
binomial law rather than a conditioned-away case. \end{proof}
The exhaustive small- and large-mean estimates in
Proposition~\ref{prop:step-009-overflow} prove
\eqref{eq:overflow-certificate}, including zero overflow for $k=2,3$.

\subsection{The executable one-use simulator}\label{app:step-010}
Fix a candidate and, locally for contradiction, suppose
\begin{equation}
n<ak\log_2^*N.
\tag{S10.2}\label{eq:local-negation}
\end{equation}
With $M=m_{n,k}$, Proposition~\ref{prop:step-004-finite-hard-prior}
provides the finite public law
\begin{equation}
\mu_{N,M}\quad\text{on}\quad\mathcal I_N.
\tag{S10.3}\label{eq:public-prior}
\end{equation}
For an arbitrary $A:\mathsf S_n\to\mathcal H_{k,N}$, the simulator to be
constructed is
\begin{equation}
B_{\mu_{N,M},A}:\mathsf Z_M\longrightarrow\mathsf G_N.
\tag{S10.5}\label{eq:simulator-type}
\end{equation}
\begin{lemma}[Input-independent public preprocessing]\label{lem:step-010-public-preprocessing}
Under Assumption~\ref{assump:candidate-regime}, the contradiction hypothesis~\eqref{eq:local-negation},
Propositions~\ref{prop:step-004-finite-hard-prior},
\ref{prop:step-005-certificate}, and
\ref{prop:step-009-overflow}, the following ordered procedure defines a
preprocessing seed \(\Omega_{\mathrm{pre}}\) before any input coordinate is
evaluated:

\begin{enumerate}
\item Draw \(J\sim\operatorname{Unif}[k]\).
\item Draw \(I_1,\ldots,I_n\) independently and uniformly from \([k]\),
independently of \(J\).
\item Conditional on \(J\), independently draw
   \[
     \Theta_i=(T_i,Q_i)\sim\mu_{N,M},
     \qquad i\in[k]\setminus\{J\};
   \]
\item Conditional on \(J,I_{1:n},(\Theta_i)_{i\ne J}\), independently draw
   \[
     X_r^\circ\sim Q_{I_r}
     \quad\text{for every }r\in[n]\text{ with }I_r\ne J,
   \]
   and set
   \[
     Y_r^\circ:=\tau_{T_{I_r}}(X_r^\circ).
   \]
\end{enumerate}

The seed consists of precisely these realized objects. Its distribution
depends only on the public tuple
\((k,N,n,M,\mu_{N,M})\), is independent of every
\(z\in\mathsf Z_M\), and contains neither a hidden instance
\((t,Q)\), a block-instance draw \(\Theta_J\), nor a synthetic feature or
label at a position whose tag equals \(J\). Conditional on the realized
selector, tags, and nonhidden instances, the nonhidden labeled
rows

\[
  \left((I_r,X_r^\circ),Y_r^\circ\right),
  \qquad r:I_r\ne J,
\tag{S10.6}\label{eq:step-010-6}
\]

are mutually independent, and each row with tag \(i\ne J\) has labeled law
\(Q_i^{\tau_{T_i}}\) after its tag is suppressed.
\end{lemma}

\begin{proof}
Proposition~\ref{prop:step-005-certificate} discharges the exact
\(N,M\) conditions for Proposition~\ref{prop:step-004-finite-hard-prior}, so the finite public law
\(\mu_{N,M}\) exists before the present simulator is defined. The first two
draws use the selector/tag product kernel certified by Proposition~\ref{prop:step-009-overflow}. Once \(J=j\) is fixed, the finite
product law

\[
  \bigotimes_{i\in[k]\setminus\{j\}}\mu_{N,M}
\]

defines all nonhidden instances. Each realized \(Q_i\) is a probability
distribution on the finite set \([N]\), so the further finite conditional
product

\[
  \bigotimes_{\substack{r\in[n]\\ I_r\ne J}}Q_{I_r}
\tag{S10.7}\label{eq:step-010-7}
\]

defines the nonhidden feature draws. For any fixed values
\((x_r)_{r:I_r\ne J}\), \eqref{eq:step-010-7} gives the exact conditional mass

\[
  \prod_{\substack{r\in[n]\\ I_r\ne J}}Q_{I_r}(x_r),
\tag{S10.8}\label{eq:step-010-8}
\]

which proves conditional mutual independence and the stated marginal
kernels. Applying the deterministic threshold label
\(\tau_{T_{I_r}}\) gives exactly the labeled law in the lemma.

All four randomization stages are completed without referring to an input
argument. In particular, the sampling code knows the public distribution
\(\mu_{N,M}\) and its own realized nonhidden pairs
\((T_i,Q_i)\), but it neither draws nor receives \((T_J,Q_J)\).
A later hidden pair \((t,Q)\) is therefore not an input to this procedure.
Endpoint draws \(T_i=1\) and \(T_i=N+1\) are executable: their deterministic
labels are respectively all one and all zero. Point-mass \(Q_i\)'s are also
valid factors in \eqref{eq:step-010-7}.

The internal randomness used by \(A\) is intentionally not a component of
\(\Omega_{\mathrm{pre}}\). It remains the fresh internal randomization of
the one call to the kernel \(A\), if that call is reached. Thus every
random choice belonging to preprocessing is fixed before input access,
without converting randomized \(A\) into a deterministic map.
\end{proof}

\begin{proposition}[Exact hidden-blind size-\(n\) row construction]\label{prop:step-010-row-construction}
Under Lemma~\ref{lem:step-010-public-preprocessing}, fix any realization
\(\omega\) of \(\Omega_{\mathrm{pre}}\) and any arbitrary labeled input

\[
  z=(z_1,\ldots,z_M)\in\mathsf Z_M,
  \qquad z_\ell=(x_\ell,y_\ell).
\tag{S10.9}\label{eq:step-010-9}
\]

Define, using only the selector and tags in \(\omega\),

\[
  L_r:=\sum_{s=1}^r\mathbf1\{I_s=J\},
  \qquad
  U:=L_n.
\tag{S10.10}\label{eq:step-010-10}
\]

The value of \(U\) is computed before any \(z_\ell\) is evaluated. If
\(U>M\), the row-construction procedure stops immediately and evaluates no
input coordinate. If \(U\le M\), it defines an ordered size-\(n\) labeled
product dataset

\[
  \mathcal P_\omega(z)
  =(\widetilde z_1,\ldots,\widetilde z_n)\in\mathsf S_n
\tag{S10.11}\label{eq:step-010-11}
\]

by the exact rule

\[
  \widetilde z_r=
  \begin{cases}
    ((J,x_{L_r}),y_{L_r}), & I_r=J,\\[2mm]
    ((I_r,X_r^\circ),Y_r^\circ), & I_r\ne J.
  \end{cases}
\tag{S10.12}\label{eq:step-010-12}
\]

Thus, on nonoverflow, the \(\ell\)-th input row is used exactly at the
\(\ell\)-th hidden-tag occurrence for every \(\ell\in[U]\), every
nonhidden row is the pre-drawn row \eqref{eq:step-010-6}, and the rule requires no knowledge
of a hidden \(t\) or \(Q\). It is valid for every binary label pattern in
\eqref{eq:step-010-9}, including patterns inconsistent with every threshold.
\end{proposition}

\begin{proof}
The seed is drawn before the input is accessed by
Lemma~\ref{lem:step-010-public-preprocessing}. Both \(L_r\) and \(U\) are
functions only of the already drawn \(J,I_1,\ldots,I_n\). Consequently the
comparison \(U>M\) can be resolved before evaluating any coordinate of
\eqref{eq:step-010-9}. On overflow, the procedure takes no later row-construction action, so
it never attempts to access a nonexistent \((M+1)\)-st record.

Suppose \(U\le M\). For each position \(r\) with \(I_r=J\), \eqref{eq:step-010-10} gives

\[
  1\le L_r\le U\le M.
\tag{S10.13}\label{eq:step-010-13}
\]

Hence \(z_{L_r}\) is an available input record. Since
\(x_{L_r}\in[N]\), \(J\in[k]\), and \(y_{L_r}\in\{0,1\}\), the first line
of \eqref{eq:step-010-12} lies in \(X_{k,N}\times\{0,1\}\), regardless of whether
\(y_{L_r}=\tau_t(x_{L_r})\) for any \(t\).

For each position with \(I_r\ne J\), the seed already contains
\(X_r^\circ\in[N]\) and
\(Y_r^\circ=\tau_{T_{I_r}}(X_r^\circ)\in\{0,1\}\), so the second line of
\eqref{eq:step-010-12} has the same required type. There is one and only one applicable line
of \eqref{eq:step-010-12} for every \(r\in[n]\). Thus \eqref{eq:step-010-11} has exactly \(n\) ordered labeled
rows and belongs to the full domain \(\mathsf S_n\) of \(A\).

If \(r\) is the \(\ell\)-th occurrence of tag \(J\), then by definition
\(L_r=\ell\), so the first line of \eqref{eq:step-010-12} uses \(z_\ell\), including its raw
label. This is why no hidden threshold \(t\) is needed. The feature
distribution \(Q\) of a later hidden input law is likewise unnecessary:
the procedure takes hidden features directly from the input records.
Nonhidden features and labels use only the simulator's publicly drawn
instances, whose exact conditional kernels were proved in
Lemma~\ref{lem:step-010-public-preprocessing}.
\end{proof}

\begin{lemma}[Ordered one-use input incidence]\label{lem:step-010-one-use}
Under Proposition~\ref{prop:step-010-row-construction}, for every fixed
preprocessing seed \(\omega\) and every arbitrary input
\(z\in\mathsf Z_M\), the following statements hold pathwise.

\begin{itemize}
\item If \(U>M\), no coordinate of \(z\) is evaluated.
\item If \(U\le M\), let
  \[
    r_\ell:=\min\left\{r\in[n]:
      \sum_{s=1}^r\mathbf1\{I_s=J\}=\ell\right\},
    \qquad \ell=1,\ldots,U.
  \tag{S10.14}\label{eq:step-010-14}
  \]
  Then
  \[
    1\le r_1<\cdots<r_U\le n,
    \qquad
    \widetilde z_{r_\ell}=((J,x_\ell),y_\ell).
  \tag{S10.15}\label{eq:step-010-15}
  \]
  No row of \(\mathcal P_\omega(z)\) other than row \(r_\ell\) depends on
  input coordinate \(z_\ell\); coordinates \(z_{U+1},\ldots,z_M\) are not
  evaluated at all.
\end{itemize}

In particular, each input record has at most one image row. If \(U=0\), the
input-access set is empty and all \(n\) rows are nonhidden pre-drawn rows.
If \(U=M\), every input record has exactly one image row. The statement is
about row incidence only and makes no differential-privacy claim.
\end{lemma}

\begin{proof}
On overflow, the first bullet is part of the control flow proved in
Proposition~\ref{prop:step-010-row-construction}. Now fix \(U\le M\). The
running count \(L_r\) in \eqref{eq:step-010-10} starts at zero, is nondecreasing, and
increases by exactly one precisely at the positions with \(I_r=J\).
Because its terminal value is \(U\), it hits each integer
\(\ell=1,\ldots,U\) at a unique first position \(r_\ell\), and those
positions are strictly increasing. At \(r_\ell\), one has
\(I_{r_\ell}=J\) and \(L_{r_\ell}=\ell\), so the first line of \eqref{eq:step-010-12} gives
the identity in \eqref{eq:step-010-15}.

At a hidden position \(r\), the only input coordinate named by \eqref{eq:step-010-12} is
\(z_{L_r}\). The map from hidden positions to positive running counts is
the bijection \(r_\ell\mapsto\ell\). At a nonhidden position, the second
line of \eqref{eq:step-010-12} is already fixed by \(\omega\) and names no input coordinate.
Therefore \(z_\ell\), for \(\ell\le U\), occurs in exactly row \(r_\ell\),
and every \(\ell>U\) occurs in no row. This proves the one-use incidence
claim for arbitrary features and labels, without any realizability premise.
\end{proof}

\begin{proposition}[Total executable improper one-block simulator]\label{prop:step-010-simulator}
Under Assumption~\ref{assump:candidate-regime}, the contradiction hypothesis~\eqref{eq:local-negation},
Propositions~\ref{prop:step-004-finite-hard-prior},
\ref{prop:step-005-certificate}, and
\ref{prop:step-009-overflow}, Lemma~\ref{lem:step-007-restriction-legality}, and
Lemma~\ref{lem:step-010-public-preprocessing},
Proposition~\ref{prop:step-010-row-construction}, and
Lemma~\ref{lem:step-010-one-use}, the following algorithm defines a
total randomized kernel
\[
  B_{\mu_{N,M},A}:\mathsf Z_M\longrightarrow\mathsf G_N
\]
for every arbitrary randomized map
\(A:\mathsf S_n\to\mathcal H_{k,N}\).

On input \(z\in\mathsf Z_M\):

\begin{enumerate}
\item Sample the complete seed \(\Omega_{\mathrm{pre}}\) in the order stated
by Lemma~\ref{lem:step-010-public-preprocessing}.
\item Compute \(U\) from \(J,I_1,\ldots,I_n\), without evaluating \(z\).
\item If \(U>M\), return the fixed hypothesis
   \[
     g_0(x):=0\quad(x\in[N])
   \tag{S10.16}\label{eq:step-010-16}
   \]
   immediately, without evaluating any coordinate of \(z\) and without
   invoking \(A\);
\item If \(U\le M\), construct the exact
   \(\mathcal P_{\Omega_{\mathrm{pre}}}(z)\in\mathsf S_n\) from \eqref{eq:step-010-12}, draw
   \[
     H\sim A(\mathcal P_{\Omega_{\mathrm{pre}}}(z))
   \]
   using \(A\)'s own fresh internal randomness, and return \(D_JH\).
\end{enumerate}

This simulator has exact input size \(M\), makes at most one call to \(A\)
and only with exact input size \(n\), uses each input row at most once, and
is defined on every arbitrary size-\(M\) labeled input. It does not know a
hidden \(t,Q\), does not require \(A\) to be proper or tag-symmetric, and
does not infer privacy or an ideal-law coupling.
\end{proposition}

\begin{proof}
The seed law is a valid input-independent probability kernel by
Lemma~\ref{lem:step-010-public-preprocessing}. On the overflow branch,
\(g_0\) is a fixed member of \(\mathsf G_N\); the control flow reaches \eqref{eq:step-010-16}
before any input dereference or call to \(A\). Thus this branch is total and
input-independent even when the supplied input has arbitrary corrupt
labels.

On the nonoverflow branch,
Proposition~\ref{prop:step-010-row-construction} proves that the argument of
\(A\) is a valid member of its complete domain \(\mathsf S_n\) and has
exactly \(n\) rows. Therefore the randomized map \(A\) returns some
\(H\in\mathcal H_{k,N}\). Since \(J\in[k]\), Lemma~\ref{lem:step-007-restriction-legality} gives
\(D_JH\in\mathsf G_N\), regardless of whether \(H\) is a threshold,
monotone, proper, symmetric, deterministic, or tag-asymmetric.

Equivalently, if \(E\subseteq\mathsf G_N\) and
\(A(s,h)\) denotes the output probability assigned by \(A\) to
\(h\in\mathcal H_{k,N}\) on \(s\in\mathsf S_n\), then the simulator kernel
is the well-defined finite expectation

\[
B_{\mu_{N,M},A}(z,E)
=
\mathbb E_{\Omega_{\mathrm{pre}}}
\left[
\begin{cases}
  \mathbf1\{g_0\in E\}, & U>M,\\[1mm]
  \displaystyle
  \sum_{h\in\mathcal H_{k,N}}
    A(\mathcal P_{\Omega_{\mathrm{pre}}}(z),h)
    \mathbf1\{D_Jh\in E\}, & U\le M.
\end{cases}
\right].
\tag{S10.17}\label{eq:step-010-17}
\]

Every term is nonnegative, the expression equals one when
\(E=\mathsf G_N\), and it is additive over disjoint events. Because all
input and output spaces are finite, \eqref{eq:step-010-17} is a total randomized kernel on
every \(z\in\mathsf Z_M\). Formula \eqref{eq:step-010-17} is a mathematical description of
the procedural branches; on \(U>M\) it contains no evaluation of
\(\mathcal P_\omega(z)\) and no \(A\)-term.

Lemma~\ref{lem:step-010-one-use} proves the exact access pattern. In
particular:

\begin{itemize}
\item At \(U=0\), no input row is read, all \(n\) rows come from the realized
  nonhidden instances, and \(A\) is called once on that size-\(n\) dataset;
\item At \(1\le U\le M\), precisely the first \(U\) input rows are read once
  each, in their original order, and all trailing rows are ignored;
\item At \(U=M\), all \(M\) rows are used once if that event is possible.
\item At \(U>M\), no row is read and \(A\) is not called.
\end{itemize}

The same clauses apply to endpoint labels, point-mass nonhidden
distributions, and arbitrary nonrealizable input labels. For \(k=2,3\),
Proposition~\ref{prop:step-009-overflow} says the overflow branch
has probability zero, but the branch remains defined. No step of this proof
uses Assumption~\ref{assump:central-dp}, compares two input datasets,
identifies an ideal sample law, or evaluates population risk. Those are
separate downstream obligations.
\end{proof}
The ordered construction is total on arbitrary labels.  In particular,
Lemma~\ref{lem:step-010-one-use} proves the pathwise incidence property,
and Proposition~\ref{prop:step-010-simulator} packages the construction as
the kernel in \eqref{eq:simulator-type}.

\subsection{Actual-to-ideal coupling}\label{app:step-011}
Fix $(t,Q)\in\mathcal I_N$ and let
\begin{equation}
Z^{\mathrm{in}}\sim(Q^{\tau_t})^M.
\tag{S11.3}\label{eq:coupling-input-law}
\end{equation}
The fixed-hidden and prior-averaged targets are, respectively,
\begin{align}
\mathbb E_{\mathrm{id},(t,Q)}R_Q(G^{\mathrm{id}},\tau_t)
&\ge
\mathbb E R_Q(B_{\mu_{N,M},A}(Z^{\mathrm{in}}),\tau_t)
-\Pr\{U>M\},
\tag{S11.4a}\label{eq:fixed-hidden-transfer}\\
\mathbb E_{\mathrm{ideal}}R_{Q_J}(D_JH^{\mathrm{id}},\tau_{T_J})
&\ge
\mathbb E_{\substack{(T,Q)\sim\mu_{N,M}\\
Z^{\mathrm{in}}\sim(Q^{\tau_T})^M}}
R_Q(B_{\mu_{N,M},A}(Z^{\mathrm{in}}),\tau_T)
-\Pr\{U>M\}.
\tag{S11.4b}\label{eq:prior-transfer}
\end{align}
\begin{lemma}[Shared realization of a finite randomized learner]\label{lem:step-011-shared-randomization}
For every randomized map
\(A:(X_{k,N}\times\{0,1\})^n\to\mathcal H_{k,N}\), there is a
deterministic function

\[
  \Phi_A:(X_{k,N}\times\{0,1\})^n\times(0,1)
  \longrightarrow\mathcal H_{k,N}
\]

such that, for \(V\sim\operatorname{Unif}(0,1)\),
\(\Phi_A(s,V)\) has exactly the law \(A(s)\) for every
\(s\in(X_{k,N}\times\{0,1\})^n\). For the same \(V\), equal input datasets give equal
outputs pathwise. No properness, determinism, or tag symmetry of \(A\) is
required.
\end{lemma}

\begin{proof}
Enumerate the finite set
\(\mathcal H_{k,N}=\{h_1,\ldots,h_q\}\). For each dataset \(s\), put

\[
  c_m(s):=\sum_{a=1}^m\Pr\{A(s)=h_a\},\qquad c_0(s):=0,
\]

and let \(\Phi_A(s,v)=h_m\) at the first index satisfying
\(c_m(s)\ge v\). Since \(c_q(s)=1\), this is total, and

\[
  \Pr\{\Phi_A(s,V)=h_m\}=c_m(s)-c_{m-1}(s)
  =\Pr\{A(s)=h_m\}.
\]

This proves the exact learner marginal. Equal datasets and the same \(V\)
give the same two arguments to the same deterministic function, proving
the pathwise synchronization statement. \end{proof}

\begin{lemma}[Fixed-hidden common-randomness coupling]\label{lem:step-011-coupling-kernel}
Under Propositions~\ref{prop:step-009-overflow} and
\ref{prop:step-010-simulator} and
Lemma~\ref{lem:step-011-shared-randomization}, fix any
\((t,Q)\in\mathcal I_N\)
and any arbitrary randomized \(A\). There exists a joint probability
kernel carrying

\begin{enumerate}
\item An infinite sequence
   \[
     W_\ell=(X_\ell^\star,Y_\ell^\star),\qquad \ell\ge1,
   \tag{S11.7}\label{eq:step-011-7}
   \]
   of mutually independent draws from \(Q^{\tau_t}\), so
   \(X_\ell^\star\sim Q\) and
   \(Y_\ell^\star=\tau_t(X_\ell^\star)\);
\item Independently of \eqref{eq:step-011-7}, the complete input-independent preprocessing seed
   of Lemma~\ref{lem:step-010-public-preprocessing}; and
\item Independently of both, one \(V\sim\operatorname{Unif}(0,1)\).
\end{enumerate}

such that the following statements hold.

\begin{itemize}
\item The actual input is
  \[
    Z^{\mathrm{in}}_\ell=W_\ell,\qquad \ell=1,\ldots,M,
  \tag{S11.8}\label{eq:step-011-8}
  \]
  and hence has exactly the law \eqref{eq:coupling-input-law}, independently of simulator
  preprocessing.
\item Put \(\Theta_J=(T_J,Q_J):=(t,Q)\), retain the seed's
  \(\Theta_i=(T_i,Q_i)\) for \(i\ne J\), write
  \(\boldsymbol\Xi=(\Theta_1,\ldots,\Theta_k)\), and define the always-present
  ideal ordered dataset
  \[
    S^{\mathrm{id}}=(\widehat z_1,\ldots,\widehat z_n)
  \tag{S11.9}\label{eq:step-011-9}
  \]
  by
  \[
    \widehat z_r=
    \begin{cases}
      ((J,X_{L_r}^\star),Y_{L_r}^\star),&I_r=J,\\[1mm]
      ((I_r,X_r^\circ),Y_r^\circ),&I_r\ne J,
    \end{cases}
    \qquad r\in[n].
  \tag{S11.10}\label{eq:step-011-10}
  \]
  The infinite sequence makes \eqref{eq:step-011-10} well defined even when \(U>M\).
\item There is a deterministic realization \(\Phi_A\) of the learner kernel
  for which
  \[
    H^{\mathrm{id}}:=\Phi_A(S^{\mathrm{id}},V),
    \qquad
    G^{\mathrm{id}}:=D_JH^{\mathrm{id}}
  \tag{S11.11}\label{eq:step-011-11}
  \]
  has the exact ideal call-and-restrict marginal.
\item Define the actual output by the simulator control flow:
  \[
    G^{\mathrm{act}}:=
    \begin{cases}
      g_0,&U>M,\\
      D_J\Phi_A(\mathcal P_{\Omega_{\mathrm{pre}}}
                    (Z^{\mathrm{in}}),V),&U\le M.
    \end{cases}
  \tag{S11.12}\label{eq:step-011-12}
  \]
  Its marginal is exactly
  \(B_{\mu_{N,M},A}(Z^{\mathrm{in}})\) with input law \eqref{eq:coupling-input-law}.
  In the first line of \eqref{eq:step-011-12}, \(V\), \eqref{eq:step-011-7}, and every other coupling-only
  variable are ignored: operationally the actual simulator still reads no
  input and makes no call to \(A\).
\end{itemize}

Thus \eqref{eq:step-011-7}--\eqref{eq:step-011-12} give a coupling of the exact actual and ideal experiments,
with every expectation and source of randomness specified.
\end{lemma}

\begin{proof}
The finite law \(Q^{\tau_t}\) defines the product sequence \eqref{eq:step-011-7}. Only its
first \(M\) coordinates are exposed as the actual input, so \eqref{eq:step-011-8} is exactly
an iid size-\(M\) input rather than a stopped or conditioned sample. By
construction this sequence is independent of the preprocessing seed, as
required by Lemma~\ref{lem:step-010-public-preprocessing}.

The running indices satisfy \(1\le L_r\le U\le n\) at every hidden
position. Hence the infinite sequence gives every hidden row in \eqref{eq:step-011-10},
including on overflow. Nonhidden rows in \eqref{eq:step-011-10} are the identical realized
rows already present in the preprocessing seed. Therefore \eqref{eq:step-011-9} is
always an ordered member of the full size-\(n\) learner domain.

Lemma~\ref{lem:step-011-shared-randomization} makes \eqref{eq:step-011-11} an exact call to
the randomized kernel \(A\), followed by
the same restriction used by the simulator. In \eqref{eq:step-011-12}, the
nonoverflow branch applies the same realization to the exact row
map, while the overflow branch is the exact fixed output from Proposition~\ref{prop:step-010-simulator}. Sampling an unused coupling
variable does not create an operational learner call. Equations \eqref{eq:step-011-8} and
\eqref{eq:step-011-12}, together with the input-independent seed law and that lemma, prove the
actual marginal claim. Equations \eqref{eq:step-011-9}--\eqref{eq:step-011-11} define the ideal marginal.
\end{proof}

\begin{lemma}[Identification with the ideal experiment]\label{lem:step-011-ideal-marginal}
Under Lemmas~\ref{lem:step-008-instance-factorization},
\ref{lem:step-008-ideal-sample-law}, and
\ref{lem:step-008-selector-independence}, and under
Lemma~\ref{lem:step-011-coupling-kernel}, fix \(j\in[k]\) and a complete
vector

\[
  \boldsymbol\theta=((t_1,Q_1),\ldots,(t_k,Q_k))
\tag{S11.14}\label{eq:step-011-14}
\]

with \((t_j,Q_j)=(t,Q)\). Conditional on \(J=j\) and on the nonhidden
seed instances taking the values in \eqref{eq:step-011-14}, but without conditioning on
\(U\), \(U\le M\), or any tag-count event, the ideal dataset \eqref{eq:step-011-9} satisfies

\[
  S^{\mathrm{id}}\sim
  \left(P_{\boldsymbol Q}^{c_{\boldsymbol t}}\right)^n.
\tag{S11.15}\label{eq:step-011-15}
\]

Equivalently, the \(n\) coupled ideal global labeled rows are mutually independent,
and each has the fixed realizable law that first selects a tag uniformly
and then draws its feature from that tag's fixed \(Q_i\). If the nonhidden
instances are not fixed, their marginal is only the corresponding mixture
of the product laws in \eqref{eq:step-011-15}; this lemma does not call that mixture iid from
one deterministic product instance. Moreover, if the hidden pair
\((t,Q)\) is first drawn from \(\mu_{N,M}\), the complete joint marginal
\((J,\boldsymbol\Xi,S^{\mathrm{id}},H^{\mathrm{id}})\) is exactly the
joint law characterized by Lemmas~\ref{lem:step-008-instance-factorization}--
\ref{lem:step-008-selector-independence}, and

\[
  \mathbb E_{(t,Q)\sim\mu_{N,M}}
    \mathbb E_{\mathrm{id},(t,Q)}R_Q(G^{\mathrm{id}},\tau_t)
  =\mathbb E_{\mathrm{ideal}}
    R_{Q_J}(D_JH^{\mathrm{id}},\tau_{T_J}).
\tag{S11.15a}\label{eq:step-011-15a}
\]
\end{lemma}

\begin{proof}
Fix ordered tags \(i_1,\ldots,i_n\in[k]\) and features
\(x_1,\ldots,x_n\in[N]\). The tag kernel gives

\[
  \Pr\{I_1=i_1,\ldots,I_n=i_n\mid J=j\}=k^{-n}.
\tag{S11.16}\label{eq:step-011-16}
\]

Conditional on this tag vector and \eqref{eq:step-011-14}, enumerate the positions with
tag \(j\) in increasing order. Formula \eqref{eq:step-011-10} assigns them the distinct
sequence variables \(X_1^\star,\ldots,X_u^\star\), where \(u\) is the
number of such positions. Their joint mass is the product of the
corresponding \(Q_j\) factors. At every position with tag \(i_r\ne j\),
Lemma~\ref{lem:step-010-public-preprocessing} supplies an
independent feature with law \(Q_{i_r}\). Those nonhidden features are
conditionally mutually independent and are independent of the separately
drawn hidden sequence \eqref{eq:step-011-7}. Therefore

\[
  \Pr\{\text{the ideal feature at row }r\text{ is }x_r
        \text{ for all }r
       \mid J=j,\boldsymbol\theta,I_{1:n}=i_{1:n}\}
  =\prod_{r=1}^n Q_{i_r}(x_r).
\tag{S11.17}\label{eq:step-011-17}
\]

All labels in \eqref{eq:step-011-10} are deterministic and equal
\(\tau_{t_{i_r}}(x_r)\). Multiplying \eqref{eq:step-011-16} and \eqref{eq:step-011-17} gives

\[
  \prod_{r=1}^n\left[
    \frac1k Q_{i_r}(x_r)
    \mathbf1\{y_r=\tau_{t_{i_r}}(x_r)\}
  \right]
\tag{S11.18}\label{eq:step-011-18}
\]

for every ordered labeled row vector. Expression \eqref{eq:step-011-18} is exactly the
product mass function in \eqref{eq:step-011-15}. This calculation sums over all tag vectors,
including those with \(U>M\); no overflow conditioning appears.

Now draw the hidden pair \((t,Q)\sim\mu_{N,M}\) independently of the
selector before applying the fixed-hidden kernel. Conditional on \(J=j\),
this pair occupies coordinate \(j\), while
Lemma~\ref{lem:step-010-public-preprocessing} draws every coordinate
\(i\ne j\) independently from the
same \(\mu_{N,M}\). These are precisely the hypotheses of Lemma~\ref{lem:step-008-instance-factorization}; hence the full vector has
law \(\mu_{N,M}^{\otimes k}\) and is independent of \(J\). The row kernel
verified in \eqref{eq:step-011-16}--\eqref{eq:step-011-18} is precisely the kernel in Lemma~\ref{lem:step-008-ideal-sample-law}. Finally,
Lemma~\ref{lem:step-011-shared-randomization} gives the exact same learner
kernel \(A\). Thus the complete coupled ideal marginal equals the ideal experiment, and its selected restriction is \(D_JH^{\mathrm{id}}\),
which proves \eqref{eq:step-011-15a}. Lemma~\ref{lem:step-008-selector-independence} and
Proposition~\ref{prop:step-008-selected-risk-identity} therefore remain
valid for this marginal; their exchangeability and product-risk
conclusions are consumed as facts and are not rederived here.

Conditioning \eqref{eq:step-011-15} further on \(U\le M\) would restrict the tag vectors in
\eqref{eq:step-011-16} and generally destroy the iid tag law. The proof never makes that
conditional-iid assertion: later equality on \(U\le M\) is pathwise under
the coupling, while \eqref{eq:step-011-15} is the unconditioned ideal sample law for the fixed
block vector. \end{proof}

\begin{proposition}[Exact row and restricted-output identity off overflow]\label{prop:step-011-no-overflow-identity}
Under Lemma~\ref{lem:step-011-coupling-kernel}, Proposition~\ref{prop:step-010-row-construction}, and Lemma~\ref{lem:step-010-one-use}, on the event

\[
  E:=\{U\le M\},
\tag{S11.19}\label{eq:step-011-19}
\]

the actual ordered dataset and ideal ordered dataset satisfy, row by row,

\[
  \mathcal P_{\Omega_{\mathrm{pre}}}(Z^{\mathrm{in}})_r
  =\widehat z_r,
  \qquad r=1,\ldots,n.
\tag{S11.20}\label{eq:step-011-20}
\]

Consequently the shared randomized learner outputs and restrictions in
\eqref{eq:step-011-11}--\eqref{eq:step-011-12} agree pathwise on \(E\), and

\[
  R_Q(G^{\mathrm{id}},\tau_t)
  =R_Q(G^{\mathrm{act}},\tau_t)
  \qquad\text{on }E.
\tag{S11.21}\label{eq:step-011-21}
\]

This includes \(U=0\) and \(U=M\). No dataset, output, or risk equality is
claimed on \(E^c=\{U>M\}\).
\end{proposition}

\begin{proof}
Fix an outcome in \(E\) and a row \(r\in[n]\). If \(I_r=J\), then

\[
  1\le L_r\le U\le M.
\tag{S11.22}\label{eq:step-011-22}
\]

The exact occurrence map in
Proposition~\ref{prop:step-010-row-construction} uses actual input row
\(L_r\).
By \eqref{eq:step-011-8}, this row is

\[
  Z^{\mathrm{in}}_{L_r}
  =(X_{L_r}^\star,Y_{L_r}^\star).
\]

The hidden-row line of the row rule in
Proposition~\ref{prop:step-010-row-construction} is therefore

\[
  ((J,X_{L_r}^\star),Y_{L_r}^\star),
\]

which is exactly the first line of the ideal rule \eqref{eq:step-011-10}. If \(I_r\ne J\),
the nonhidden-row line of
Proposition~\ref{prop:step-010-row-construction} and the nonhidden-row line
of \eqref{eq:step-011-10} are literally the same realized
nonhidden row \(((I_r,X_r^\circ),Y_r^\circ)\). This exhaustive row split
proves \eqref{eq:step-011-20}, with ordering preserved.

At \(U=0\), there is no hidden row and the second-line identity applies at
all \(n\) positions; both experiments call \(A\) on the identical all-
nonhidden dataset. At \(U=M\), \eqref{eq:step-011-22} still holds, and the exact occurrence
map may use all and only the first \(M\) sequence records. Thus neither
endpoint requires a separate conditioning argument or an extra record.

Because \eqref{eq:step-011-20} gives equal datasets and both calls use the same \(V\),
Lemma~\ref{lem:step-011-shared-randomization} gives equal hypotheses even when
\(A\) is randomized, improper, nonmonotone, or tag-asymmetric. Both sides
then apply the same restriction \(D_J\), so the one-block outputs are
equal. Evaluating the identical functions against the same fixed
\((Q,\tau_t)\) proves \eqref{eq:step-011-21}.

On \(E^c\), Proposition~\ref{prop:step-010-simulator} makes the
actual side return \(g_0\) before reading the input or calling \(A\),
whereas the ideal side still uses \eqref{eq:step-011-9}--\eqref{eq:step-011-11}. The proof makes no comparison
between those datasets or outputs there. \end{proof}

\begin{proposition}[Bounded selected-risk transfer]\label{prop:step-011-risk-transfer}
Under Lemma~\ref{lem:step-011-ideal-marginal},
Proposition~\ref{prop:step-011-no-overflow-identity}, and Proposition~\ref{prop:step-009-overflow}, for every fixed hidden pair
\((t,Q)\) and arbitrary randomized \(A\), define

\[
  L_{\mathrm{id}}:=R_Q(G^{\mathrm{id}},\tau_t),
  \qquad
  L_{\mathrm{act}}:=R_Q(G^{\mathrm{act}},\tau_t).
\tag{S11.23}\label{eq:step-011-23}
\]

Then, on the full coupling space,

\[
  L_{\mathrm{id}}
  \ge L_{\mathrm{act}}-\mathbf1\{U>M\}.
\tag{S11.24}\label{eq:step-011-24}
\]

Taking the exact marginals from
Lemma~\ref{lem:step-011-coupling-kernel} proves \eqref{eq:fixed-hidden-transfer}. Averaging the hidden
pair under \(\mu_{N,M}\) and applying
Lemma~\ref{lem:step-011-ideal-marginal} proves the exact ideal
interface \eqref{eq:prior-transfer}. Proposition~\ref{prop:step-009-overflow}
supplies the strict residual in \eqref{eq:overflow-certificate}.
If \(k=2\) or \(k=3\), then \(U\le M\) pathwise and the same result has
zero residual.
\end{proposition}

\begin{proof}
Population 0-1 risk is a probability, so for every output function,

\[
  0\le L_{\mathrm{id}}\le1,
  \qquad
  0\le L_{\mathrm{act}}\le1.
\tag{S11.25}\label{eq:step-011-25}
\]

On \(E\), Proposition~\ref{prop:step-011-no-overflow-identity} gives
\(L_{\mathrm{id}}=L_{\mathrm{act}}\), so \eqref{eq:step-011-24} holds with a zero
indicator. On \(E^c\), \eqref{eq:step-011-25} gives

\[
  L_{\mathrm{act}}-1\le0\le L_{\mathrm{id}},
\]

which is exactly \eqref{eq:step-011-24} with indicator one. This overflow argument uses only
boundedness. In particular, the actual early-abort output is \(g_0\), but
its risk may be any value in \([0,1]\): it is zero for the all-zero target,
one for the all-one target under any \(Q\), and can take an intermediate
value for other thresholds. No equality or favorable abort-risk property
is assumed.

Taking expectation of \eqref{eq:step-011-24} over the joint variables in
Lemma~\ref{lem:step-011-coupling-kernel} gives

\[
  \mathbb E_{\mathrm{id},\theta}L_{\mathrm{id}}
  \ge
  \mathbb E_{\mathrm{act},\theta}L_{\mathrm{act}}
  -\mathbb E\mathbf1\{U>M\}.
\tag{S11.26}\label{eq:step-011-26}
\]

The actual marginal identity in that lemma identifies the middle term with
the simulator expectation displayed in \eqref{eq:fixed-hidden-transfer}, including all of its
preprocessing and learner randomness. The ideal marginal identifies the
first term with the ideal selected-risk expectation. Finally,

\[
  \mathbb E\mathbf1\{U>M\}=\Pr\{U>M\}<2^{-9}
\tag{S11.27}\label{eq:step-011-27}
\]

by Proposition~\ref{prop:step-009-overflow}. There is no conditioning-to-expectation
conversion and no interchange of a limit, supremum, or integral: the
ideal dataset uses only finitely many coordinates of \eqref{eq:step-011-7}, and all other
spaces are finite.

Because \eqref{eq:step-011-26} holds for every fixed \((t,Q)\), average it over
\((t,Q)\sim\mu_{N,M}\). The overflow term is unchanged because its law
depends only on \(J,I_{1:n}\). Lemma~\ref{lem:step-011-ideal-marginal},
specifically \eqref{eq:step-011-15a}, identifies the averaged ideal term
with the selected-risk quantity in
Proposition~\ref{prop:step-008-selected-risk-identity}. The averaged
actual term is
the right-hand simulator expectation in \eqref{eq:prior-transfer}. This proves \eqref{eq:prior-transfer} without
rederiving or using the equality between ideal selected risk and
product population risk.

For \(k=2,3\), Proposition~\ref{prop:step-009-overflow} gives
\(M\ge n\ge U\), so \eqref{eq:step-011-21} holds everywhere and the final term in \eqref{eq:step-011-26} is
exactly zero. For \(n<k\) or \(M=8\), the coupling and pointwise inequality
are unchanged, while the floor-eight estimate supplies the stated
strict residual. \end{proof}
The shared realization of $A$ synchronizes outputs whenever the ordered
datasets agree.  Proposition~\ref{prop:step-011-risk-transfer} proves
\eqref{eq:fixed-hidden-transfer}--\eqref{eq:prior-transfer} with one and
only one overflow charge.

\subsection{One-charge privacy of the simulator}\label{app:step-012}
Continue under \eqref{eq:local-negation}.  For every
$z\simeq z'$ in $\mathsf Z_M$ and every $E\subseteq\mathsf G_N$, the
target simulator inequalities are
\begin{align}
B_{\mu_{N,M},A}(z,E)
&\le e^\varepsilon B_{\mu_{N,M},A}(z',E)+\delta,
\tag{S12.4}\label{eq:simulator-dp}\\
B_{\mu_{N,M},A}(z,E)
&\le e^{0.1}B_{\mu_{N,M},A}(z',E)+\Delta_M,
\qquad
\Delta_M=\frac{d_*}{M^2\log M}.
\tag{S12.6}\label{eq:simulator-source-cap}
\end{align}
\begin{lemma}[Deterministic one-row preprocessing map]\label{lem:step-012-one-row-map}
Under Assumption~\ref{assump:candidate-regime}, the contradiction hypothesis~\eqref{eq:local-negation},
and Proposition~\ref{prop:step-010-row-construction} and
Lemma~\ref{lem:step-010-one-use}, fix any realization \(\omega\) of
the input-independent preprocessing seed and any arbitrary
\(z,z'\in\mathsf Z_M\) with \(z\simeq z'\).  Then exactly one of the
following holds.

\begin{enumerate}
\item If \(U(\omega)>M\), the preprocessing control flow evaluates no
   coordinate of either input, so its output branch is identical on
   \(z,z'\).
\item If \(U(\omega)\le M\), the two exact size-\(n\) datasets
   \(\mathcal P_\omega(z)\) and \(\mathcal P_\omega(z')\) are equal or
   replacement-adjacent.  More precisely, if the inputs differ only at
   row \(q\), then they are equal when \(q>U\), while for \(q\le U\)
   they can differ only at the single global row \(r_q\), the \(q\)-th
   hidden-tag occurrence.
\end{enumerate}

This assertion includes \(U=0\), a changed unused record, a changed used
record, simultaneous replacement of the record's feature and label, and
arbitrary possibly nonrealizable label patterns.  All nonhidden global
rows and all hidden rows other than \(r_q\) coincide.  In particular, the
map never produces group-adjacent datasets.
\end{lemma}

\begin{proof}
Proposition~\ref{prop:step-010-row-construction} gives

\[
  U=\sum_{r=1}^n\mathbf1\{I_r=J\},
\tag{S12.10}\label{eq:step-012-10}
\]

as a function only of \(J,I_{1:n}\), all of which belong to \(\omega\).
Hence the same value of \(U\), and therefore the same overflow branch, is
used on \(z\) and \(z'\).  When \(U>M\), Proposition~\ref{prop:step-010-row-construction} stops before either input
is read.  This proves the first branch pathwise, not only in distribution.

Suppose \(U\le M\).  Lemma~\ref{lem:step-010-one-use} gives distinct
hidden positions

\[
  1\le r_1<\cdots<r_U\le n
\tag{S12.11}\label{eq:step-012-11}
\]

depending only on \(\omega\), and its exact row rule gives, for
\(z_\ell=(x_\ell,y_\ell)\),

\[
  \bigl(\mathcal P_\omega(z)\bigr)_{r_\ell}
  =((J,x_\ell),y_\ell),
  \qquad \ell=1,\ldots,U.
\tag{S12.12}\label{eq:step-012-12}
\]

Every other position is a nonhidden row already fixed by \(\omega\).
The same formulas, with primes, hold for \(z'\), with exactly the same
positions \(r_\ell\) and exactly the same nonhidden rows.

If \(z=z'\), the synthesized datasets are equal.  Otherwise adjacency in
\(\mathsf Z_M\) supplies one index \(q\in[M]\) such that
\(z_\ell=z'_\ell\) for every \(\ell\ne q\).  If \(q>U\), Lemma~\ref{lem:step-010-one-use} says row \(q\) is unused, and \eqref{eq:step-012-12}
shows equality of every used hidden row; all nonhidden rows also coincide.
Thus

\[
  \mathcal P_\omega(z)=\mathcal P_\omega(z').
\tag{S12.13}\label{eq:step-012-13}
\]

If \(q\le U\), then for every hidden row \(r_\ell\) with \(\ell\ne q\),
the equality \(z_\ell=z'_\ell\) and \eqref{eq:step-012-12} give equal rows.  At the sole
remaining hidden position,

\[
  \bigl(\mathcal P_\omega(z)\bigr)_{r_q}=((J,x_q),y_q),
  \qquad
  \bigl(\mathcal P_\omega(z')\bigr)_{r_q}=((J,x'_q),y'_q).
\tag{S12.14}\label{eq:step-012-14}
\]

Equation \eqref{eq:step-012-14} is one replacement even if both \(x_q\ne x'_q\) and
\(y_q\ne y'_q\).  It is still a legal replacement when either raw label
is inconsistent with every threshold because both labels remain elements
of \(\{0,1\}\), and Assumption~\ref{assump:central-dp} is stated on all
labeled datasets.  Every nonhidden row is a component of the fixed seed,
so no nonhidden row changes.  Thus the synthesized datasets differ in at
most the one row \(r_q\), proving replacement adjacency and excluding any
group-adjacency charge.

Finally, if \(U=0\), there are no positions in \eqref{eq:step-012-11}, every synthesized
row is nonhidden and seed-fixed, and \eqref{eq:step-012-13} holds for every pair of inputs.
This also proves the requested null case. \end{proof}

\begin{proposition}[Fixed-seed privacy kernel with learner coins retained]\label{prop:step-012-fixed-seed-kernel}
Under Assumptions~\ref{assump:candidate-regime} and
\ref{assump:central-dp}, the contradiction hypothesis~\eqref{eq:local-negation}, Proposition~\ref{prop:step-010-simulator}, and
Lemma~\ref{lem:step-012-one-row-map}, fix a preprocessing seed \(\omega\)
but do not fix the internal coins of \(A\).  Define the all-zero global
hypothesis

\[
  h_0(i,x):=0
  \qquad ((i,x)\in X_{k,N}),
\tag{S12.15}\label{eq:step-012-15}
\]

and the global-output kernel \(K_\omega:\mathsf Z_M\to\mathcal H_{k,N}\)
by

\[
  K_\omega(z,F)=
  \begin{cases}
    \mathbf1\{h_0\in F\}, & U(\omega)>M,\\[1mm]
    \Pr_A\!\left[A(\mathcal P_\omega(z))\in F\right],
      & U(\omega)\le M,
  \end{cases}
\tag{S12.16}\label{eq:step-012-16}
\]

for every \(F\subseteq\mathcal H_{k,N}\).  Then, for all arbitrary
\(z\simeq z'\) in \(\mathsf Z_M\),

\[
  K_\omega(z,F)
  \le e^\varepsilon K_\omega(z',F)+\delta
  \qquad(F\subseteq\mathcal H_{k,N}).
\tag{S12.17}\label{eq:step-012-17}
\]

On overflow the two output laws coincide exactly.  On nonoverflow, \eqref{eq:step-012-17}
uses Assumption~\ref{assump:central-dp} at most once.  The statement holds
for arbitrary randomized \(A\); its internal coins are precisely the
randomness under \(\Pr_A\) in \eqref{eq:step-012-16}, not part of the fixed seed.
\end{proposition}

\begin{proof}
The function \(h_0\) belongs to the full output space
\(\mathcal H_{k,N}\), and \(D_Jh_0=g_0\) for every \(J\).  Thus the first
line of \eqref{eq:step-012-16} is an analytical global-output lift of the simulator's actual
fixed overflow output; it does not call \(A\).

If \(U(\omega)>M\), the seed fixes the same branch for both inputs, and

\[
  K_\omega(z,F)=\mathbf1\{h_0\in F\}=K_\omega(z',F).
\tag{S12.18}\label{eq:step-012-18}
\]

Since \(e^\varepsilon\ge1\) and \(\delta\ge0\), \eqref{eq:step-012-18} implies \eqref{eq:step-012-17}.  This
also shows directly that overflow contributes no privacy residual.

Suppose \(U(\omega)\le M\).  By
Lemma~\ref{lem:step-012-one-row-map},
\(\mathcal P_\omega(z)=\mathcal P_\omega(z')\) or
\(\mathcal P_\omega(z)\simeq\mathcal P_\omega(z')\).  In the equality
case the two laws in \eqref{eq:step-012-16} coincide, including when \(U=0\) or the changed
input row is unused.  In the adjacent case, Assumption~\ref{assump:central-dp}
applied once to the exact size-\(n\) pair and the event \(F\) gives

\[
\begin{aligned}
  K_\omega(z,F)
  &=\Pr_A[A(\mathcal P_\omega(z))\in F]\\
  &\le e^\varepsilon
    \Pr_A[A(\mathcal P_\omega(z'))\in F]+\delta\\
  &=e^\varepsilon K_\omega(z',F)+\delta.
\end{aligned}
\tag{S12.19}\label{eq:step-012-19}
\]

No coin of \(A\) was conditioned on or fixed in \eqref{eq:step-012-19}.  Arbitrary feature
and label replacements are covered because the premise of \eqref{eq:step-012-19} is the
replacement adjacency proved in Lemma~\ref{lem:step-012-one-row-map}, not
realizability.  There is one invocation of
Assumption~\ref{assump:central-dp}, so neither group privacy
nor sequential composition appears. \end{proof}

\begin{lemma}[Common input-independent mixture and restriction preserve one charge]\label{lem:step-012-postprocess-mixture}
Under Assumptions~\ref{assump:candidate-regime} and
\ref{assump:central-dp}, the contradiction hypothesis~\eqref{eq:local-negation}, Lemma~\ref{lem:step-010-public-preprocessing} and
Proposition~\ref{prop:step-010-simulator}, and
Proposition~\ref{prop:step-012-fixed-seed-kernel}, let \(\nu\) be the law
of the input-independent preprocessing seed on its finite seed space
\(\mathsf\Omega_{\mathrm{pre}}\).  For
\(G\subseteq\mathsf\Omega_{\mathrm{pre}}\times\mathcal H_{k,N}\), define
its seed section

\[
  G_\omega:=\{h\in\mathcal H_{k,N}:(\omega,h)\in G\},
\tag{S12.20}\label{eq:step-012-20}
\]

and form the joint seed/global-output kernel

\[
  \widehat K(z,G)
  :=\sum_{\omega\in\mathsf\Omega_{\mathrm{pre}}}
       \nu(\omega)K_\omega(z,G_\omega).
\tag{S12.21}\label{eq:step-012-21}
\]

Then \(\widehat K\) is \((\varepsilon,\delta)\)-DP.  Applying only after
this common mixture the deterministic postprocessing map

\[
  \Phi(\omega,h):=D_{J(\omega)}h
\tag{S12.22}\label{eq:step-012-22}
\]

produces exactly the simulator law.  Consequently \eqref{eq:simulator-dp} holds with
exactly one additive \(\delta\).
\end{lemma}

\begin{proof}
The seed space is finite: Lemma~\ref{lem:step-010-public-preprocessing}
samples finitely many tags,
finitely many elements of the finite-support prior, and features from the
finite domain \([N]\).  For each \(G\), every section \(G_\omega\) in
\eqref{eq:step-012-20} is therefore a valid event in the global output space.  Equation
\eqref{eq:step-012-21} is a total kernel: all summands are nonnegative, its value on
\(\mathsf\Omega_{\mathrm{pre}}\times\mathcal H_{k,N}\) is
\(\sum_\omega\nu(\omega)K_\omega(z,\mathcal H_{k,N})=1\), and finite
additivity follows by taking disjoint seed sections and using additivity of
each \(K_\omega\).  The seedwise inequality \eqref{eq:step-012-17} is uniform in \(\omega\)
and in the event, so for arbitrary \(z\simeq z'\),

\[
\begin{aligned}
  \widehat K(z,G)
  &=\sum_\omega\nu(\omega)K_\omega(z,G_\omega)\\
  &\le\sum_\omega\nu(\omega)
       \bigl(e^\varepsilon K_\omega(z',G_\omega)+\delta\bigr)\\
  &=e^\varepsilon\sum_\omega\nu(\omega)K_\omega(z',G_\omega)
    +\delta\sum_\omega\nu(\omega)\\
  &=e^\varepsilon\widehat K(z',G)+\delta.
\end{aligned}
\tag{S12.23}\label{eq:step-012-23}
\]

Here Lemma~\ref{lem:step-010-public-preprocessing} supplies the
same input-independent weights \(\nu(\omega)\) for \(z\) and \(z'\).  Also,
\[
\sum_\omega\nu(\omega)=1.
\]
Thus \eqref{eq:step-012-23} first proves exact DP for the
common joint mixture.  The overflow subset
\(\{\omega:U(\omega)>M\}\), the \(U=0\) subset, and every nonoverflow
subset are determined solely by this common seed.  No input-dependent
conditioning, normalization, or change of measure occurs.

Now fix \(E\subseteq\mathsf G_N\) and use the one joint-output event

\[
  G_E:=\Phi^{-1}(E)
  =\{(\omega,h):D_{J(\omega)}h\in E\}.
\tag{S12.24}\label{eq:step-012-24}
\]

On an overflow seed, \eqref{eq:step-012-15} gives \(\Phi(\omega,h_0)=g_0\), so that seed's
contribution to \(\widehat K(z,G_E)\) is
\(\nu(\omega)\mathbf1\{g_0\in E\}\), independently of \(z\).  On a
nonoverflow seed, its contribution is

\[
  \nu(\omega)\Pr_A\!\left[
    D_{J(\omega)}A(\mathcal P_\omega(z))\in E
  \right],
\tag{S12.25}\label{eq:step-012-25}
\]

with all internal randomness of \(A\) retained.  Summing these two exact
branches is precisely Proposition~\ref{prop:step-010-simulator},
so

\[
  B_{\mu_{N,M},A}(z,E)=\widehat K(z,G_E).
\tag{S12.26}\label{eq:step-012-26}
\]

Finally apply \eqref{eq:step-012-23} to \(G=G_E\) and use \eqref{eq:step-012-26} for both inputs.  This is the
event-preimage proof of postprocessing after the common mixture and gives
\eqref{eq:simulator-dp}.  The additive term in \eqref{eq:step-012-23} is \(\delta\), not \(\delta\) times the
number of seeds, \(n\), \(k\), \(U\), or an overflow factor. \end{proof}

\begin{proposition}[Exact source-cap membership]\label{prop:step-012-source-membership}
Under Assumptions~\ref{assump:candidate-regime} and
\ref{assump:central-dp}, the contradiction hypothesis~\eqref{eq:local-negation}, Propositions~\ref{prop:step-005-certificate} and
\ref{prop:step-010-simulator}, and
Lemma~\ref{lem:step-012-postprocess-mixture}, define the exact source-cap
learner set at the present \(N,M\) by

\[
\begin{aligned}
  \mathcal K^{\mathrm{src}}_{N,M}:=
  \bigl\{K:\;&\mathsf Z_M\longrightarrow\mathsf G_N
      \text{ a total randomized kernel}:\\[-1mm]
    &K(z,E)\le e^{0.1}K(z',E)+\Delta_M\\[-1mm]
    &\text{for all }z\simeq z'\text{ in }\mathsf Z_M
      \text{ and }E\subseteq\mathsf G_N\bigr\}.
\end{aligned}
\tag{S12.27}\label{eq:step-012-27}
\]

Then

\[
  B_{\mu_{N,M},A}\in\mathcal K^{\mathrm{src}}_{N,M}.
\tag{S12.28}\label{eq:step-012-28}
\]

More fully, the same parameters satisfy
\[
N\ge N_*,\qquad 8\le M<b_*\log_2^*N,\qquad
0<\varepsilon\le0.1,\qquad 0<\delta<\Delta_M,
\]
the simulator has exact input size \(M\), full arbitrary one-block
output space \(\mathsf G_N\), and is total on all labeled inputs, and its
event probabilities satisfy \eqref{eq:simulator-source-cap}.  Thus \eqref{eq:step-012-28} is exact source-interface
membership, not an inference from the phrase
"\((\varepsilon,\delta)\)-DP" alone.
\end{proposition}

\begin{proof}
Proposition~\ref{prop:step-010-simulator} proves that
\(B_{\mu_{N,M},A}\) is a total randomized kernel with domain exactly
\(\mathsf Z_M\) and codomain exactly the full improper space
\(\mathsf G_N\).  Its totality includes overflow, \(U=0\), arbitrary
corrupt input labels, and unused trailing records.

Proposition~\ref{prop:step-005-certificate}, applied to the same
\(N,M,\varepsilon,\delta\), gives
$N\ge N_*$, $8\le M<b_*\log_2^*N$,
$0<\varepsilon\le0.1$, and $0<\delta<\Delta_M$.  In particular,
\(M\ge8\) implies \(\log M>0\), so the value
\(\Delta_M=d_*/(M^2\log M)\) in
\eqref{eq:simulator-source-cap} is a well-defined positive source cap.
Lemma~\ref{lem:step-012-postprocess-mixture} gives, for every
defining pair and event in \eqref{eq:step-012-27},

\[
  B_{\mu_{N,M},A}(z,E)
  \le e^\varepsilon B_{\mu_{N,M},A}(z',E)+\delta.
\tag{S12.29}\label{eq:step-012-29}
\]

Because probabilities are nonnegative, \(\varepsilon\le0.1\), and
\(\delta<\Delta_M\),

\[
\begin{aligned}
  e^\varepsilon B_{\mu_{N,M},A}(z',E)+\delta
  &\le e^{0.1}B_{\mu_{N,M},A}(z',E)+\delta\\
  &< e^{0.1}B_{\mu_{N,M},A}(z',E)+\Delta_M.
\end{aligned}
\tag{S12.30}\label{eq:step-012-30}
\]

The weak inequality required in \eqref{eq:step-012-27} follows.  The scalar
inequalities from Proposition~\ref{prop:step-005-certificate},
\eqref{eq:step-012-27}, and
\eqref{eq:step-012-30}, together with the exact total kernel typing, check every component
of the source interface separately.  No accuracy, realizability, or
properness property has been added to that learner set, and no private
eligibility is inferred without the event-level derivation \eqref{eq:step-012-29}--\eqref{eq:step-012-30}.
This proves \eqref{eq:step-012-28}. \end{proof}
The fixed-seed adjacency lemma, the single invocation of
Assumption~\ref{assump:central-dp}, and the common seed mixture prove
\eqref{eq:simulator-dp}.  Proposition~\ref{prop:step-012-source-membership}
then verifies the exact source-cap typing and parameters in
\eqref{eq:simulator-source-cap}.

\subsection{Prior-averaged product risk}\label{app:step-013}
Set $\eta=2^{-8}$ and $B=B_{\mu_{N,M},A}$.  In the ideal experiment let
\begin{equation}
\boldsymbol\Xi\sim\mu_{N,M}^{\otimes k},
\qquad
S^{\mathrm{id}}\mid\boldsymbol\Xi
\sim(P_{\boldsymbol Q}^{c_{\boldsymbol T}})^n,
\qquad
H^{\mathrm{id}}\sim A(S^{\mathrm{id}}).
\tag{S13.3}\label{eq:product-lower-experiment}
\end{equation}
The conclusion sought is
\begin{equation}
\mathbb E R_{P_{\boldsymbol Q}}(H^{\mathrm{id}},c_{\boldsymbol T})
>2^{-9}.
\tag{S13.4}\label{eq:prior-product-lower}
\end{equation}
\begin{proposition}[Hard-prior loss of the certified simulator]\label{prop:step-013-simulator-hardness}
Under Proposition~\ref{prop:step-004-finite-hard-prior} and
Proposition~\ref{prop:step-012-source-membership}, if the total
simulator \(B=B_{\mu_{N,M},A}\) is instantiated at their identical
admissible \(N,M\) and public prior \(\mu_{N,M}\), then

\[
\mathbb E_{\substack{(T,Q)\sim\mu_{N,M}\\
                      Z\sim(Q^{\tau_T})^M\\ B}}
R_Q(B(Z),\tau_T)>\eta=2^{-8}.
\tag{S13.9}\label{eq:step-013-9}
\]

The expectation includes the prior draw, exactly \(M\) i.i.d. realizable
one-block input rows conditional on that draw, all simulator preprocessing,
and any internal coins of \(A\) used by the simulator.
\end{proposition}

\begin{proof}
Proposition~\ref{prop:step-012-source-membership} gives every
membership clause needed by the universal learner quantifier in Proposition~\ref{prop:step-004-finite-hard-prior}: \(B\) is total on the
entire arbitrary labeled input space of exact size \(M\), its output is an
arbitrary member of \(\{0,1\}^{[N]}\), and it satisfies every event-level
one-row replacement inequality at privacy parameters

\[
\left(0.1,\frac{d_*}{M^2\log M}\right).
\tag{S13.10}\label{eq:step-013-10}
\]

Both propositions use the same finite ordered domain \([N]\),
thresholds \(\tau_t\), population-risk convention, source constants, and
integers \(N,M\). Thus no convention, object, sample-size, or privacy bridge
remains to prove.

The prior in Proposition~\ref{prop:step-004-finite-hard-prior} was selected
before the universal learner
quantifier. Although the code of \(B\) uses the public
\(\mu_{N,M}\), it is therefore one of the prior-aware kernels already
covered by that quantifier. Substituting this exact \(B\) into
Proposition~\ref{prop:step-004-finite-hard-prior} yields
\eqref{eq:step-013-9} directly. The substitution neither conditions on overflow nor changes
the simulator on any input, so it introduces no residual. \end{proof}

\begin{proposition}[Prior-averaged product-risk transfer]\label{prop:step-013-product-lower-bound}
Under Propositions~\ref{prop:step-008-selected-risk-identity},
\ref{prop:step-009-overflow}, \ref{prop:step-011-risk-transfer}, and
\ref{prop:step-013-simulator-hardness}, the ideal experiment
\eqref{eq:product-lower-experiment} obeys the exact
product-risk lower bound \eqref{eq:prior-product-lower}.
\end{proposition}

\begin{proof}
Define only for this proof

\[
\mathcal L_{\mathrm{act}}
:=
\mathbb E_{\substack{(T,Q)\sim\mu_{N,M}\\
                      Z\sim(Q^{\tau_T})^M\\ B}}
R_Q(B(Z),\tau_T)
\tag{S13.11}\label{eq:step-013-11}
\]

and

\[
\mathcal L_{\mathrm{id}}
:=
\mathbb E_{\mathrm{ideal}}
R_{Q_J}(D_JH^{\mathrm{id}},\tau_{T_J}).
\tag{S13.12}\label{eq:step-013-12}
\]

These are proof-local names for the two exact quantities already typed by
the dependencies. Proposition~\ref{prop:step-013-simulator-hardness}
gives

\[
\mathcal L_{\mathrm{act}}>2^{-8}.
\tag{S13.13}\label{eq:step-013-13}
\]

Proposition~\ref{prop:step-011-risk-transfer} and the notation
\(p_{\mathrm{ov}}=\Pr(U>M)\), with its bound supplied by
Proposition~\ref{prop:step-009-overflow}, give

\[
\mathcal L_{\mathrm{id}}
\ge \mathcal L_{\mathrm{act}}-p_{\mathrm{ov}}.
\tag{S13.14}\label{eq:step-013-14}
\]

Combining \eqref{eq:step-013-13} with the weak inequality \eqref{eq:step-013-14} preserves strictness:

\[
\mathcal L_{\mathrm{id}}
>2^{-8}-p_{\mathrm{ov}}.
\tag{S13.15}\label{eq:step-013-15}
\]

Proposition~\ref{prop:step-009-overflow} gives the strict upper
bound \(p_{\mathrm{ov}}<2^{-9}\). Since subtraction reverses the comparison
in the subtracted term,

\[
2^{-8}-p_{\mathrm{ov}}
>2^{-8}-2^{-9}
=2\cdot2^{-9}-2^{-9}
=2^{-9}.
\tag{S13.16}\label{eq:step-013-16}
\]

Equations \eqref{eq:step-013-15}--\eqref{eq:step-013-16} therefore yield

\[
\mathcal L_{\mathrm{id}}>2^{-9}.
\tag{S13.17}\label{eq:step-013-17}
\]

The overflow event has been charged exactly once: it occurs only in \eqref{eq:step-013-14},
and no later equality or bound adds another copy.

Finally, Proposition~\ref{prop:step-008-selected-risk-identity}
identifies \eqref{eq:step-013-12}, with no residual, as

\[
\mathcal L_{\mathrm{id}}
=
\mathbb E_{\substack{\boldsymbol\Xi\sim\mu_{N,M}^{\otimes k}\\
                      S^{\mathrm{id}}\sim
                      (P_{\boldsymbol Q}^{c_{\boldsymbol T}})^n\\
                      H^{\mathrm{id}}\sim A(S^{\mathrm{id}})}}
R_{P_{\boldsymbol Q}}(H^{\mathrm{id}},c_{\boldsymbol T}).
\tag{S13.18}\label{eq:step-013-18}
\]

Substituting \eqref{eq:step-013-18} into \eqref{eq:step-013-17} proves \eqref{eq:prior-product-lower}. The equality in \eqref{eq:step-013-18} is taken under
the unconditional ideal experiment; the conditional i.i.d. sample law is
invoked only after fixing \(\boldsymbol\Xi\), and no conditioning on
\(U\le M\) appears. \end{proof}
Proposition~\ref{prop:step-013-simulator-hardness} applies the public prior
only after source-cap membership is established.  The ledger in
Proposition~\ref{prop:step-013-product-lower-bound} preserves strictness and
proves \eqref{eq:prior-product-lower}.

\subsection{A fixed hard product instance}\label{app:step-014}
For $\boldsymbol z=((t_i,Q_i))_{i=1}^k\in\mathcal I_N^k$, define
\begin{equation}
\Phi_A(\boldsymbol z)
=\mathbb E_{\substack{S\sim(P_{\boldsymbol Q}^{c_{\boldsymbol t}})^n\\
\rho_A}}
R_{P_{\boldsymbol Q}}(A_{\rho_A}(S),c_{\boldsymbol t}).
\tag{S14.3}\label{eq:fixed-instance-functional}
\end{equation}
\begin{proposition}[Fixed support point from a strict prior average]\label{prop:step-014-fixed-extraction}
Under Proposition~\ref{prop:step-013-product-lower-bound}, fix its
same \(k,N,M,n,A\), finite prior \(\mu_{N,M}\), and ideal experiment. Then
there exists a deterministic vector
\[
\boldsymbol z^*\in\operatorname{supp}(\mu_{N,M}^{\otimes k})
\subseteq\mathcal I_N^k
\tag{S14.8}\label{eq:step-014-8}
\]
such that
\[
\mathbb E_{\substack{S\sim(P_{\boldsymbol Q^*}^{c_{\boldsymbol t^*}})^n\\
\rho_A}}
R_{P_{\boldsymbol Q^*}}\bigl(A_{\rho_A}(S),c_{\boldsymbol t^*}\bigr)
>2^{-9}.
\tag{S14.9}\label{eq:step-014-9}
\]
Once \(\boldsymbol z^*\) is selected, the only randomness in \eqref{eq:step-014-9} is the
fixed-instance sample and the learner coins.
\end{proposition}

\begin{proof}
Let
\[
\mathcal S=\operatorname{supp}(\mu_{N,M}^{\otimes k}).
\tag{S14.10}\label{eq:step-014-10}
\]
Proposition~\ref{prop:step-013-product-lower-bound} supplies a finite
probability law, so
\(\mathcal S\) is finite. For each \(\boldsymbol z\in\mathcal S\), put
\[
w_{\boldsymbol z}=\mu_{N,M}^{\otimes k}(\{\boldsymbol z\}).
\tag{S14.11}\label{eq:step-014-11}
\]
Then \(w_{\boldsymbol z}>0\) and
\(\sum_{\boldsymbol z\in\mathcal S}w_{\boldsymbol z}=1\).

Conditional on \(\boldsymbol\Xi=\boldsymbol z\), the sample and output
experiment in Proposition~\ref{prop:step-013-product-lower-bound}
is exactly the fixed-instance experiment defining \(\Phi_A(\boldsymbol z)\).
Finite expansion of the outer prior expectation therefore gives
\[
\begin{aligned}
&\mathbb E_{\substack{
  \boldsymbol\Xi\sim\mu_{N,M}^{\otimes k}\\
  S^{\mathrm{id}}\sim(P_{\boldsymbol Q}^{c_{\boldsymbol T}})^n\\
  H^{\mathrm{id}}\sim A(S^{\mathrm{id}})}}
  R_{P_{\boldsymbol Q}}(H^{\mathrm{id}},c_{\boldsymbol T})\\
&\qquad=\sum_{\boldsymbol z\in\mathcal S}
  w_{\boldsymbol z}\Phi_A(\boldsymbol z).
\end{aligned}
\tag{S14.12}\label{eq:step-014-12}
\]
This is finite conditional expectation. It does not assert that the
unconditional prior mixture is iid from one deterministic distribution.

The left side of \eqref{eq:step-014-12} is strictly greater than \(2^{-9}\) by Proposition~\ref{prop:step-013-product-lower-bound}; hence
\[
\sum_{\boldsymbol z\in\mathcal S}w_{\boldsymbol z}\Phi_A(\boldsymbol z)
>2^{-9}.
\tag{S14.13}\label{eq:step-014-13}
\]
If every \(\boldsymbol z\in\mathcal S\) had
\(\Phi_A(\boldsymbol z)\le2^{-9}\), then
\[
\sum_{\boldsymbol z\in\mathcal S}w_{\boldsymbol z}\Phi_A(\boldsymbol z)
\le\sum_{\boldsymbol z\in\mathcal S}w_{\boldsymbol z}2^{-9}
=2^{-9},
\tag{S14.14}\label{eq:step-014-14}
\]
contradicting \eqref{eq:step-014-13}. Thus some \(\boldsymbol z^*\in\mathcal S\) has
\(\Phi_A(\boldsymbol z^*)>2^{-9}\), which is \eqref{eq:step-014-9}. \end{proof}

\begin{proposition}[Same-instance PAC contradiction]\label{prop:step-014-fixed-contradiction}
Under Proposition~\ref{prop:step-014-fixed-extraction} and Proposition~\ref{prop:step-006-pointwise-pac-expectation}, fix the
deterministic vector \(\boldsymbol z^*\) extracted in \eqref{eq:step-014-8}. Then
\[
2^{-9}<\Phi_A(\boldsymbol z^*)\le2^{-12},
\tag{S14.15}\label{eq:step-014-15}
\]
which is impossible.
\end{proposition}

\begin{proof}
The extraction proposition gives
\[
\Phi_A(\boldsymbol z^*)>2^{-9}.
\tag{S14.16}\label{eq:step-014-16}
\]
This lower bound concerns the fixed \(P_{\boldsymbol Q^*}\), fixed target
\(c_{\boldsymbol t^*}\), exact sample size \(n\), fixed randomized learner
\(A\), its internal coin law, and the setting-defined population 0-1 risk.
No component is changed by selecting \(\boldsymbol z^*\).

Because \(\boldsymbol z^*\in\mathcal I_N^k\), Proposition~\ref{prop:step-006-pointwise-pac-expectation}, which is the
pointwise consequence of Assumption~\ref{assump:distribution-free-realizable-pac},
applies to this same vector:
\[
\begin{aligned}
\Phi_A(\boldsymbol z^*)
&=\mathbb E_{\substack{
  S\sim(P_{\boldsymbol Q^*}^{c_{\boldsymbol t^*}})^n\\ \rho_A}}
  R_{P_{\boldsymbol Q^*}}\bigl(A_{\rho_A}(S),c_{\boldsymbol t^*}\bigr)\\
&\le\alpha_0+\beta_0
=2^{-13}+2^{-13}=2^{-12}.
\end{aligned}
\tag{S14.17}\label{eq:step-014-17}
\]
Equations \eqref{eq:step-014-16} and \eqref{eq:step-014-17} concern literally the same deterministic number;
neither includes a prior average. Finally,
\[
2^{-9}=8\cdot2^{-12}>2^{-12},
\tag{S14.18}\label{eq:step-014-18}
\]
so \eqref{eq:step-014-16}--\eqref{eq:step-014-17} contradict one another. \end{proof}
Finite-support extraction produces one deterministic $\boldsymbol z^*$
with $\Phi_A(\boldsymbol z^*)>2^{-9}$, while the pointwise PAC result gives
$\Phi_A(\boldsymbol z^*)\le2^{-12}$.  Proposition~\ref{prop:step-014-fixed-contradiction}
records this same-instance contradiction.

\subsection{Nonasymptotic closure}\label{app:step-015}
The constants are those in \eqref{eq:calibration}, fixed before choosing
\begin{equation}
k\ge2,
\qquad N\ge N_0,
\qquad n\ge1,
\tag{S15.T2}\label{eq:candidate-order}
\end{equation}
the admissible privacy pair, and the arbitrary randomized map
\begin{equation}
A:(X_{k,N}\times\{0,1\})^n\longrightarrow\mathcal H_{k,N}.
\tag{S15.T4}\label{eq:candidate-map}
\end{equation}
\begin{lemma}[Impossibility of the negated candidate bound]\label{lem:step-015-negated-target}
Under Assumptions~\ref{assump:candidate-regime},
\ref{assump:central-dp}, and
\ref{assump:distribution-free-realizable-pac},
Proposition~\ref{prop:step-005-certificate}, and
Proposition~\ref{prop:step-014-fixed-contradiction}, fix any candidate in
the order \eqref{eq:candidate-order}--\eqref{eq:candidate-map}, with the
privacy pair chosen between the integer tuple and the learner. If

\[
n<ak\log_2^*N,
\tag{S15.1}\label{eq:step-015-1}
\]

then the preceding propositions yield a contradiction. Hence no
such fixed candidate can satisfy \eqref{eq:step-015-1}.
\end{lemma}

\begin{proof}
The constants in \eqref{eq:calibration}, including \(a=b_*/16\), and
\(N_0\) were fixed by
Proposition~\ref{prop:step-005-certificate} before the present
candidate was chosen. Apply that proposition to the fixed candidate and the
strict hypothesis \eqref{eq:step-015-1}. It yields, at the exact integer budget

\[
M=\max\left\{8,\left\lceil\frac{4n}{k}\right\rceil\right\},
\]

the strict hard-regime and source-privacy conclusions

\[
N\ge N_*,\qquad 8\le M<b_*\log_2^*N,
\qquad 0<\varepsilon\le0.1,
\qquad 0<\delta<\frac{d_*}{M^2\log M}.
\tag{S15.2}\label{eq:step-015-2}
\]

The same proposition retains, rather than replaces, the primitive
conjunction

\[
0<\delta\le\frac{1}{n\log(n+1)},
\qquad
0<\delta\le\frac{c_\delta}{M^2\log(M+1)}.
\tag{S15.3}\label{eq:step-015-3}
\]

The fixed learner still satisfies its exact
Assumption~\ref{assump:central-dp} and
Assumption~\ref{assump:distribution-free-realizable-pac} premises; neither
is altered by the scalar specialization \eqref{eq:step-015-2}. Therefore the complete
premise package for
Proposition~\ref{prop:step-014-fixed-contradiction} holds at this identical
candidate. Applying that proposition, without reopening any of its internal
dependencies, gives one fixed product instance
\((P_{\boldsymbol Q^*},c_{\boldsymbol t^*})\) for which

\[
\mathbb E_{\substack{
 S\sim(P_{\boldsymbol Q^*}^{c_{\boldsymbol t^*}})^n\\
 \rho_A}}
R_{P_{\boldsymbol Q^*}}
\bigl(A_{\rho_A}(S),c_{\boldsymbol t^*}\bigr)>2^{-9}
\tag{S15.4}\label{eq:step-015-4}
\]

and, for literally the same fixed distribution, target, size-\(n\) iid
sample law, randomized learner, internal-coin law, and population risk,

\[
\mathbb E_{\substack{
 S\sim(P_{\boldsymbol Q^*}^{c_{\boldsymbol t^*}})^n\\
 \rho_A}}
R_{P_{\boldsymbol Q^*}}
\bigl(A_{\rho_A}(S),c_{\boldsymbol t^*}\bigr)\le2^{-12}.
\tag{S15.5}\label{eq:step-015-5}
\]

But

\[
2^{-9}=8\cdot2^{-12}>2^{-12},
\tag{S15.6}\label{eq:step-015-6}
\]

so no real number can satisfy both \eqref{eq:step-015-4} and \eqref{eq:step-015-5}. This contradicts \eqref{eq:step-015-1}.

There is no exceptional boundary branch. Proposition~\ref{prop:step-005-certificate} derives \eqref{eq:step-015-2}--\eqref{eq:step-015-3} with the same
strict inequalities when \(N=N_0\), \(M=8\), \(n=1\), \(n<k\), or
\(k=2\), and Proposition~\ref{prop:step-014-fixed-contradiction} applies to that same
candidate and exact fixed-instance experiment. No floor, ceiling, endpoint,
privacy parameter, or sample size is changed in this invocation. \end{proof}

\begin{proposition}[Nonasymptotic candidate closure]\label{prop:step-015-exact-closure}
Fix the universal constants in \eqref{eq:calibration} and the fixed
\(N_0\) before
candidate quantification. Under Assumptions~\ref{assump:candidate-regime},
\ref{assump:central-dp}, and
\ref{assump:distribution-free-realizable-pac}, every candidate quantified
in that same candidate order satisfies

\[
n\ge a k\log_2^*N.
\tag{S15.7}\label{eq:step-015-7}
\]

This is a deterministic pointwise implication for the exact candidate. It
makes no assertion outside the approved candidate regime.
\end{proposition}

\begin{proof}
Fix an arbitrary candidate covered by the proposition. If \eqref{eq:step-015-7} failed, then
the total order on the real numbers would give its exact strict negation

\[
n<ak\log_2^*N.
\tag{S15.8}\label{eq:step-015-8}
\]

No integer rounding is inserted into this logical negation. But
Lemma~\ref{lem:step-015-negated-target} proves that \eqref{eq:step-015-8} is impossible for a
candidate satisfying the three primitive assumptions. Therefore \eqref{eq:step-015-7} holds.
Because the constants were fixed before the arbitrary candidate and the
candidate was otherwise arbitrary within that order, the conclusion has the
quantifier order stated in the formalized goal. \end{proof}
Lemma~\ref{lem:step-015-negated-target} excludes the exact strict negation
for every candidate in this order.  Proposition~\ref{prop:step-015-exact-closure}
therefore yields the weak nonasymptotic lower bound with no rounding or
scope change.

\subsection{Tower-diagonal specialization}\label{app:step-016}
For compactness in the following proof, write

\[
T_s=\operatorname{Tow}_2(s)\qquad(s\ge1)
\]

\begin{lemma}[Exact tower inversion and floor arithmetic]\label{lem:step-016-tower-arithmetic}
Under the tower and iterated-logarithm definitions in
Section~\ref{sec:setup}, the map
\(s\mapsto T_s\) is strictly increasing and unbounded, with
\(T_s\ge2^s\) for every integer \(s\ge1\). The function
\(u\mapsto\log_2^*u\) is nondecreasing on \((0,\infty)\). Moreover, for
every integer \(r\ge1\),

\[
\log_2^*T_r=r
\tag{S16.1}\label{eq:step-016-1}
\]

and, for every integer \(r\ge2\),

\[
\left\lfloor\log_2(T_r+1)\right\rfloor=T_{r-1}.
\tag{S16.2}\label{eq:step-016-2}
\]
\end{lemma}

\begin{proof}
Because \(T_1=2\) and \(T_{s+1}=2^{T_s}\), induction gives
\(T_s\ge2^s\). Indeed, the claim is true at \(s=1\), and the elementary
induction \(2^s\ge s+1\) gives

\[
T_{s+1}=2^{T_s}\ge2^{2^s}\ge2^{s+1}.
\]

Also \(2^m>m\) for every integer \(m\ge2\), so
\(T_{s+1}>T_s\). Thus the tower is strictly increasing and, from
\(T_s\ge2^s\), unbounded.

For monotonicity of log-star, first note that for \(u>1\),

\[
\log_2^*u=1+\log_2^*(\log_2u).
\tag{S16.3}\label{eq:step-016-3}
\]

If \(0<u\le v\), induction on \(\log_2^*v\), using the monotonicity of
\(\log_2\) and \eqref{eq:step-016-3}, gives \(\log_2^*u\le\log_2^*v\). The base case
\(\log_2^*v=0\) means \(v\le1\), so also \(u\le1\), and both log-star
values are zero. In the induction step, \(u\le1\) is immediate; if
\(u>1\), apply the induction hypothesis to
\(0<\log_2u\le\log_2v\) in \eqref{eq:step-016-3}.

For \(0\le j\le r-1\), repeated use of the tower recursion yields

\[
\log_2^{(j)}T_r=T_{r-j}.
\tag{S16.4}\label{eq:step-016-4}
\]

In particular, the \((r-1)\)-fold iterate equals \(T_1=2>1\), whereas
the \(r\)-fold iterate equals \(\log_2 2=1\). By the minimum in the
definition of log-star, \eqref{eq:step-016-1} follows exactly.

Finally, put \(q=T_{r-1}\). Then \(T_r=2^q>1\), and hence

\[
2^q<T_r+1<2T_r=2^{q+1}.
\]

Taking base-two logarithms gives

\[
q<\log_2(T_r+1)<q+1.
\]

Since \(q\) is an integer, its floor is exactly \(q=T_{r-1}\), proving
\eqref{eq:step-016-2}. \end{proof}

\begin{proposition}[Exact diagonal structural identities]\label{prop:step-016-structural-diagonal}
Under the primitive diagonal specialization \(r\in\mathbb Z_{\ge2}\),
\(k=r\), \(N=T_r\), Lemma~\ref{lem:step-001-cardinality}, Lemma~\ref{lem:step-001-vc}, Proposition~\ref{prop:step-001-product-ld}, and
Lemma~\ref{lem:step-016-tower-arithmetic},

\[
\operatorname{VC}(C_{r,T_r})=r,
\qquad
\operatorname{LD}(C_{r,T_r})=rT_{r-1},
\qquad
|C_{r,T_r}|=(T_r+1)^r.
\tag{S16.5}\label{eq:step-016-5}
\]

With the setting convention that unadorned \(\log\) is natural,

\[
rT_{r-1}\log2
  <\log|C_{r,T_r}|
  <\frac32rT_{r-1}\log2.
\tag{S16.6}\label{eq:step-016-6}
\]

In particular,
\(\log|C_{r,T_r}|=\Theta(rT_{r-1})\) with universal hidden constants.
\end{proposition}

\begin{proof}
The structural identities give, after the exact substitution
\(k=r,N=T_r\),

\[
\operatorname{VC}(C_{r,T_r})=r,
\]

\[
\operatorname{LD}(C_{r,T_r})
=r\left\lfloor\log_2(T_r+1)\right\rfloor,
\]

and \(|C_{r,T_r}|=(T_r+1)^r\). Equation \eqref{eq:step-016-2} in
Lemma~\ref{lem:step-016-tower-arithmetic} turns the middle display into
\(rT_{r-1}\), proving every exact identity in \eqref{eq:step-016-5}.

For \eqref{eq:step-016-6}, write \(T_r=2^{T_{r-1}}\). The same strict comparison used for
the floor gives

\[
T_{r-1}\log2
  <\log(T_r+1)
  <(T_{r-1}+1)\log2.
\]

Since \(r\ge2\) implies \(T_{r-1}\ge T_1=2\),
\(T_{r-1}+1\le(3/2)T_{r-1}\). Multiplication by \(r\) and the exact
identity \(\log|C_{r,T_r}|=r\log(T_r+1)\) prove \eqref{eq:step-016-6}. The constants
\(\log2\) and \((3/2)\log2\) are absolute and independent of \(r\).
\end{proof}

\begin{lemma}[Iterated logarithm of the diagonal Littlestone dimension]\label{lem:step-016-ld-logstar}
Under Lemma~\ref{lem:step-016-tower-arithmetic} and
Proposition~\ref{prop:step-016-structural-diagonal}, for every integer
\(r\ge2\),

\[
r-1
\le
\log_2^*\operatorname{LD}(C_{r,T_r})
\le r.
\tag{S16.7}\label{eq:step-016-7}
\]

Consequently,

\[
\frac r2
\le
\log_2^*\operatorname{LD}(C_{r,T_r})
\le r,
\tag{S16.8}\label{eq:step-016-8}
\]

so \(\log_2^*\operatorname{LD}(C_{r,T_r})=\Theta(r)\) with constants
\(1/2\) and \(1\), independent of \(r\).
\end{lemma}

\begin{proof}
Let \(q=T_{r-1}\). Proposition~\ref{prop:step-016-structural-diagonal}
gives

\[
\operatorname{LD}(C_{r,T_r})=rq\ge q=T_{r-1}.
\tag{S16.9}\label{eq:step-016-9}
\]

We next prove the upper comparison \(rq\le T_r=2^q\). At \(r=2\),
\(q=T_1=2\) and both sides equal \(4\). For \(r\ge3\), induction gives
\(T_{r-1}\ge r\): the base is \(T_2=4\ge3\), and
\(T_s\ge s+1\) implies \(T_{s+1}=2^{T_s}\ge2^{s+1}\ge s+2\).
Also, \(2^x\ge x^2\) for every integer \(x\ge4\); this follows by
induction from equality at \(x=4\) and
\(2x^2\ge(x+1)^2\) for \(x\ge3\). Therefore, for \(r\ge3\),

\[
T_r=2^q\ge q^2\ge rq.
\tag{S16.10}\label{eq:step-016-10}
\]

Combining \eqref{eq:step-016-9}--\eqref{eq:step-016-10}, monotonicity of log-star from
Lemma~\ref{lem:step-016-tower-arithmetic}, and its exact tower inversion
gives

\[
\log_2^*T_{r-1}=r-1
\le\log_2^*(rq)
\le\log_2^*T_r=r,
\]

which is \eqref{eq:step-016-7}. Finally, \(r-1\ge r/2\) for every \(r\ge2\), yielding
\eqref{eq:step-016-8} and the stated uniform \(\Theta\)-constants. \end{proof}

\begin{proposition}[Tower-diagonal Rate Specialization Bridge]\label{prop:step-016-rate-bridge}
Fix the universal constants
\[
a,c_\delta,\varepsilon_0,\alpha_0,\beta_0,N_0
\]
from Proposition~\ref{prop:step-015-exact-closure}, before any candidate is
chosen. Thus \(a=b_*/16\), \(c_\delta=d_*\), \(\varepsilon_0=0.1\), and
\(\alpha_0=\beta_0=2^{-13}\). Define
the finite deterministic index

\[
r_0=\min\{s\in\mathbb Z_{\ge2}:T_s\ge N_0\}.
\tag{S16.11}\label{eq:step-016-11}
\]

Under Assumptions~\ref{assump:candidate-regime},
\ref{assump:central-dp}, and
\ref{assump:distribution-free-realizable-pac}, every diagonal candidate
with integer \(r\ge r_0\), \(k=r\), and \(N=T_r\) satisfies

\[
n\ge ar^2
\tag{S16.12}\label{eq:step-016-12}
\]

and

\[
n\ge
a\,\operatorname{VC}(C_{r,T_r})
   \log_2^*\operatorname{LD}(C_{r,T_r}).
\tag{S16.13}\label{eq:step-016-13}
\]

Moreover,

\[
\frac12r^2
\le
\operatorname{VC}(C_{r,T_r})
\log_2^*\operatorname{LD}(C_{r,T_r})
\le r^2.
\tag{S16.14}\label{eq:step-016-14}
\]

Thus \eqref{eq:step-016-12}--\eqref{eq:step-016-13} are exactly
\(n=\Omega(r^2)=\Omega(\operatorname{VC}\log_2^*\operatorname{LD})\),
with constants independent of \(r\), and only within the approved
candidate/privacy/PAC regime.
\end{proposition}

\begin{proof}
Lemma~\ref{lem:step-016-tower-arithmetic} proves that \(T_s\) is
increasing and unbounded, so the set in \eqref{eq:step-016-11} is nonempty and \(r_0\) is a
finite integer depending only on the already-fixed \(N_0\). For every
\(r\ge r_0\), one has \(r\ge2\) and \(T_r\ge T_{r_0}\ge N_0\).
Therefore \(k=r,N=T_r\) meet the structural range of Proposition~\ref{prop:step-015-exact-closure}.

No privacy condition is absorbed in this substitution. The candidate has
the exact budget

\[
m_{n,r}=\max\left\{8,\left\lceil\frac{4n}{r}\right\rceil\right\},
\]

and Assumption~\ref{assump:candidate-regime} remains exactly

\[
n\in\mathbb Z_{\ge1},
\quad 0<\varepsilon\le\varepsilon_0,
\quad
0<\delta\le
\min\left\{
\frac1{n\log(n+1)},
\frac{c_\delta}{m_{n,r}^2\log(m_{n,r}+1)}
\right\}.
\tag{S16.15}\label{eq:step-016-15}
\]

The learner remains an arbitrary randomized map

\[
A:(X_{r,T_r}\times\{0,1\})^n
  \longrightarrow\{0,1\}^{X_{r,T_r}}
\]

satisfying the exact central replacement-DP and distribution-free
realizable-PAC assumptions at this same candidate. Proposition~\ref{prop:step-015-exact-closure} therefore applies without a
mode or scope conversion and gives

\[
n\ge ar\log_2^*T_r.
\]

Equation \eqref{eq:step-016-1} in Lemma~\ref{lem:step-016-tower-arithmetic} makes the
right-hand side exactly \(ar^2\), proving \eqref{eq:step-016-12}. No asymptotic term is
dropped in this equality.

By Proposition~\ref{prop:step-016-structural-diagonal} and
Lemma~\ref{lem:step-016-ld-logstar},

\[
\operatorname{VC}(C_{r,T_r})=r,
\qquad
\frac r2
\le\log_2^*\operatorname{LD}(C_{r,T_r})\le r.
\]

Multiplication proves \eqref{eq:step-016-14}. In particular, the upper half of \eqref{eq:step-016-14} and
\eqref{eq:step-016-12} give

\[
n\ge ar^2
\ge
a\,\operatorname{VC}(C_{r,T_r})
   \log_2^*\operatorname{LD}(C_{r,T_r}),
\]

which is \eqref{eq:step-016-13} with the same inherited constant \(a\).

The finite initial indices are not hidden inside an \(r\)-dependent
constant. At \(r=2\), the arithmetic itself is exact:

\[
T_1=2,\qquad T_2=4,\qquad \log_2^*T_2=2,\qquad
\lfloor\log_2(5)\rfloor=2,
\]

and hence \(\operatorname{VC}=2\), \(\operatorname{LD}=4\),
\(|C|=25\), and \(\log_2^*\operatorname{LD}=2\). The same structural
identities hold for every \(2\le r<r_0\). By the minimality of \(r_0\)
and monotonicity of the tower, every such index has \(T_r<N_0\). Therefore
Proposition~\ref{prop:step-015-exact-closure} is inapplicable, so no learning
lower-bound conclusion is asserted for those indices. The asymptotic statements begin
at the single fixed index \(r_0\). \end{proof}

\begin{proposition}[Diagonal separation and unresolved scales]\label{prop:step-016-scale-comparison}
Under Proposition~\ref{prop:step-016-structural-diagonal} and
Lemma~\ref{lem:step-016-ld-logstar}, for every integer \(r\ge2\),

\[
2r-1
\le
\operatorname{VC}(C_{r,T_r})
+\log_2^*\operatorname{LD}(C_{r,T_r})
\le2r,
\tag{S16.16}\label{eq:step-016-16}
\]

and

\[
r2^{r-1}
\le
\operatorname{VC}(C_{r,T_r})
2^{\log_2^*\operatorname{LD}(C_{r,T_r})}
\le r2^r.
\tag{S16.17}\label{eq:step-016-17}
\]

Consequently, the established lower-bound scale \(r^2\) exceeds the
additive expression by an unbounded factor, but

\[
r^2
=o\!\left(
\operatorname{VC}(C_{r,T_r})
2^{\log_2^*\operatorname{LD}(C_{r,T_r})}
\right)
\tag{S16.18}\label{eq:step-016-18}
\]

and

\[
r^2=o(\log|C_{r,T_r}|).
\tag{S16.19}\label{eq:step-016-19}
\]

These comparisons concern the size of the proved lower-bound expression;
they are not upper bounds on the true sample complexity.
\end{proposition}

\begin{proof}
Substitute \(\operatorname{VC}=r\) and the two bounds in \eqref{eq:step-016-7}. Addition
gives \eqref{eq:step-016-16}, and exponentiation followed by multiplication by \(r\) gives
\eqref{eq:step-016-17}. In particular,

\[
\frac{r^2}{
 \operatorname{VC}(C_{r,T_r})
 +\log_2^*\operatorname{LD}(C_{r,T_r})}
\ge\frac r2\longrightarrow\infty,
\]

which proves the unbounded-factor comparison with the additive scale.

For the first remaining scale, the lower half of \eqref{eq:step-016-17} gives

\[
0\le
\frac{r^2}{
 \operatorname{VC}(C_{r,T_r})
 2^{\log_2^*\operatorname{LD}(C_{r,T_r})}}
\le\frac r{2^{r-1}}.
\tag{S16.20}\label{eq:step-016-20}
\]

The right-hand side tends to zero: its ratio at consecutive indices is
\((r+1)/(2r)\le3/4\) for \(r\ge2\). This proves \eqref{eq:step-016-18}.

For cardinality, \eqref{eq:step-016-6} and \(T_{r-1}\ge2^{r-1}\) from
Lemma~\ref{lem:step-016-tower-arithmetic} yield

\[
0\le\frac{r^2}{\log|C_{r,T_r}|}
<\frac{r}{T_{r-1}\log2}
\le\frac{r}{2^{r-1}\log2}\longrightarrow0,
\]

which proves \eqref{eq:step-016-19}. \end{proof}

\begin{proof}[Assembly of the tower-diagonal specialization]
Lemma~\ref{lem:step-016-tower-arithmetic} uses the exact convention
\[
T_1=2,\qquad T_{s+1}=2^{T_s},
\]
and the minimum-based definition of log-star to
prove
\[
\log_2^*T_r=r,
\qquad
\left\lfloor\log_2(T_r+1)\right\rfloor=T_{r-1}.
\]
Combining the floor identity with Lemmas~\ref{lem:step-001-cardinality} and
\ref{lem:step-001-vc} and
Proposition~\ref{prop:step-001-product-ld},
Proposition~\ref{prop:step-016-structural-diagonal} gives the exact formulas
\[
\operatorname{VC}(C_{r,T_r})=r,
\qquad
\operatorname{LD}(C_{r,T_r})=rT_{r-1},
\qquad
|C_{r,T_r}|=(T_r+1)^r.
\]
Lemma~\ref{lem:step-016-ld-logstar} sharpens the asymptotic notation to the
two-sided finite-index inequality
\[
r-1\le
\log_2^*\operatorname{LD}(C_{r,T_r})
\le r,
\]
and hence to
\[
\frac12r^2\le
\operatorname{VC}(C_{r,T_r})
\log_2^*\operatorname{LD}(C_{r,T_r})
\le r^2.
\]

Let $r_0$ be the fixed index in
Proposition~\ref{prop:step-016-rate-bridge}.  For every $r\ge r_0$, that
proposition verifies the unchanged candidate, privacy, and PAC conditions
before applying Proposition~\ref{prop:step-015-exact-closure}; it therefore
gives, with the same universal $a$,
\[
n\ge ar^2
\ge a\,\operatorname{VC}(C_{r,T_r})
       \log_2^*\operatorname{LD}(C_{r,T_r}).
\]
All structural identities remain valid for the finite indices
$2\le r<r_0$, while no learning conclusion is asserted there because
$T_r<N_0$.

Finally, Proposition~\ref{prop:step-016-scale-comparison} proves
\[
2r-1\le
\operatorname{VC}(C_{r,T_r})+
\log_2^*\operatorname{LD}(C_{r,T_r})
\le2r,
\]
so the quadratic lower bound exceeds the additive expression by an
unbounded factor.  The same proposition proves
\[
r^2=o\!\left(
\operatorname{VC}(C_{r,T_r})
2^{\log_2^*\operatorname{LD}(C_{r,T_r})}
\right),
\qquad
r^2=o(\log|C_{r,T_r}|).
\]
Thus the result remains a partial resolution: it proves neither the
stronger exponential milestone nor the $\Omega(\log|C|)$ lower bound,
and it supplies neither a universal upper bound nor a matching
combinatorial characterization for all finite-Littlestone classes.
\end{proof}

\subsection{Proof of the Main Theorem}\label{app:main-proof}

\begin{proof}[Proof of Theorem~\ref{thm:main}]
Choose $b_*,d_*,N_*$ from
Proposition~\ref{prop:step-002-wrapper}, and then choose
$a,c_\delta,\varepsilon_0,\alpha_0,\beta_0,N_0$ as in
Lemma~\ref{lem:step-005-calibration} and \eqref{eq:calibration}.  These
constants are universal and precede every candidate quantifier.

Fix an arbitrary candidate satisfying
Assumptions~\ref{assump:candidate-regime},
\ref{assump:central-dp}, and
\ref{assump:distribution-free-realizable-pac}.  Proposition~\ref{prop:step-015-exact-closure}
gives
\[
n\ge ak\log_2^*N,
\]
which proves \eqref{eq:main-nonasymptotic} with the exact fixed-size,
replacement-private, population-risk, and unrestricted-improper scope.

For $T_s=\operatorname{Tow}_2(s)$, Lemma~\ref{lem:step-016-tower-arithmetic}
proves $\log_2^*T_r=r$ and the exact floor identity.
Proposition~\ref{prop:step-016-structural-diagonal} and
Lemma~\ref{lem:step-016-ld-logstar} then give
\eqref{eq:main-structure}--\eqref{eq:main-product-scale}.  The tower is
increasing and unbounded, so the displayed $r_0$ is finite.  For every
$r\ge r_0$, Proposition~\ref{prop:step-016-rate-bridge} applies the
nonasymptotic bound to the same admissible diagonal candidate and proves
\eqref{eq:main-diagonal}.  Finally,
Proposition~\ref{prop:step-016-scale-comparison} gives
\eqref{eq:main-additive}--\eqref{eq:main-remaining-scales}.  This proves all
parts of the theorem.
\end{proof}
